\section{Modularity arguments for the window bound}\label{sec:pf:modularity}

\providecommand{\modularitycyc}{\chi_{\mathrm{cyc}}}
\providecommand{\modularityG}{\mathcal G}
\providecommand{\modularityQZ}{\Q_\ell/\Z_\ell}
\providecommand{\modularitySha}{\mathrm{sha}}
\providecommand{\modularitycot}{\lambda^{\mathrm{cot}}}

This section proves three of the nine inputs of the proof of \cref{thm:window:HE} given in
\cref{sec:pf:vwin-he:proof}, namely the steps \ref{pf:vwin-he:b}, \ref{pf:vwin-he:e} and \ref{pf:vwin-he:g} there:
\begin{itemize}
\item $\sfname{HypothesesH'}$ implies the automorphy statements $(\mathrm{U}_{\mathrm{ord}})$ and $(\mathrm{U}_i)$
(\cref{pf:modularity:ends0});
\item the component \ref{pf:vwin-he:H1} of $\sfname{HypothesesH'}$ and $(\mathrm{U}_{\mathrm{ord}})$ imply
$(\mathrm{U}_{\mathrm{ur}})$ (\cref{pf:modularity:unr-main});
\item $\sfname{HypothesesH'}$, Lemma K and $(\mathrm{Sp})$ imply $\sfname{CountII}$ and $\sfname{CountIII}$
(\cref{pf:modularity:count23}).
\end{itemize}
The four automorphy statements $(\mathrm{Sp})$, $(\mathrm{U}_{\mathrm{ord}})$, $(\mathrm{U}_i)$, $(\mathrm{U}_{\mathrm{ur}})$
were only described in \cref{pf:vwin-he:autstmts}; here they are defined precisely
(\cref{sec:pf:modularity:stmts}). The section also contains the results on standard $L$-functions of cuspidal
representations of $\GL_n$ (the Godement--Jacquet functional equation and the unramified computation) which the
formal development proves in the same group of modules and which are used in the Eisenstein congruence arguments
(\cref{sec:pf:modularity:GJ}).

\emph{The idea.} All three results belong to the quantitative form of Ribet's method \cite{Ribet1976} described in
\cref{sec:window:eis}. Let us recall it for cases II and III, which are treated in full in
\cref{sec:pf:modularity:count23}. For a character $\psi$ of a host and a twist $e$, one chooses a \emph{raw datum}
$D$: a four-dimensional reducible ``Eisenstein'' representation
$r_0=\theta\oplus\Psi\oplus\alpha\oplus1$ of $G_\Q$ with values in $\GSp_4$, built from $\Psi=\modularitycyc^e\psi$
and auxiliary characters, together with local conditions. The Selmer-type group whose length is to be bounded is
compared (Lemma M$'$) with the cotangent space in degree $-1$ of a presentation of the deformation ring of $r_0$,
graded by the action of a torus; Lemma M bounds this cotangent length by twice the sum of the valuations of
$\tr r(\sigma)-\tr r_0(\sigma)$ over representatives $r$ of the orbits of the characteristic-zero points. The
hypotheses of Lemma M are provided by the automorphy statement $(\mathrm{Sp})$ (finiteness, absence of endoscopic
points) and by a rigidity theorem (the points have no non-trivial infinitesimal deformations), itself a consequence of an adjoint
vanishing theorem. Finally the points $r$ are automorphic of bounded level, so that the product of the differences
$\tr r(\sigma)-\tr r_0(\sigma)$ divides an explicit non-zero integer, an Eisenstein congruence number, whose size is
polynomially bounded. The statements $(\mathrm{U}_{\mathrm{ord}})$, $(\mathrm{U}_{\mathrm{ur}})$ and $(\mathrm{U}_i)$
play the same role for the three-dimensional unitary deformation problems used for cases I and V, which are treated
elsewhere (\lref{UpperTheory.V.Z.J.count_i_of_printed_bound_k_kato_automorphy},
\lref{UpperTheory.V.Z.J.count_v_of_printed_bound_k_reduced_unitary_lift}).

\emph{Organisation.} \Cref{sec:pf:modularity:conv} fixes conventions and the notions of automorphy.
\Cref{sec:pf:modularity:stmts} defines the four automorphy statements. \Cref{sec:pf:modularity:unr} deduces
$(\mathrm{U}_{\mathrm{ur}})$ from $(\mathrm{U}_{\mathrm{ord}})$; the argument is a specialisation argument from the
universal deformation ring to ordinary points with integral coefficients. \Cref{sec:pf:modularity:count23} proves
$\sfname{CountII}$ and $\sfname{CountIII}$: the main argument (\cref{pf:modularity:mainbig}), the key inequality
(\cref{pf:modularity:ineq}), the Eisenstein congruence number (\cref{pf:modularity:frob},
\cref{pf:modularity:numberSp4}) and rigidity (\cref{pf:modularity:rig}). \Cref{sec:pf:modularity:ends} proves
$(\mathrm{U}_{\mathrm{ord}})$ and $(\mathrm{U}_i)$ from the linked-chain statements, and the consequence $(\mathrm{Unl})$ of
these statements used for cases I and V (\cref{pf:modularity:unlend}). \Cref{sec:pf:modularity:compl} contains auxiliary
results on algebraic characters, window conditions and newforms, and \cref{sec:pf:modularity:modU} proves the modularity
statements for universal ordinary determinant rings and the other inputs of the linked-chain statements.
\Cref{sec:pf:modularity:numU3} proves the
Eisenstein congruence number for $\Uni(3)$ used for cases I and V, and \cref{sec:pf:modularity:GJ}
contains the results on standard $L$-functions. Several long auxiliary results of the formal development are cited
with a precise statement; their proofs are not reproduced here.

\subsection{Conventions}\label{sec:pf:modularity:conv}

We use the conventions of \cref{sec:pf:vwin-he:counts}: $\Omega$ is an algebraic closure of $\Q$, $G_F=\Gal(\Omega/F)$
for a number field $F\subset\Omega$ (more generally for a number field $F$ with a given embedding into $\Omega$), and
for a valuation ring $A$ of $\Omega$ we write $D_A\supset I_A$ for its decomposition and inertia groups in $G_\Q$ (or
in $G_F$, when $F$ is specified). We say that $A$ \emph{lies over} the prime $p$ if $p$ is in the maximal ideal of $A$.
An element $\sigma\in D_A$ is a \emph{Frobenius element} at $A$ if it acts on the residue field of $A$ as $x\mapsto x^p$.
A \emph{complex conjugation} is an element $c\in G_\Q$ such that $j\circ c=\overline{j}$ for some embedding
$j:\Omega\to\C$.

Fix a prime $\ell$, an algebraic closure $\Qlbar$ of $\Ql$ with its absolute value $\|\cdot\|$ normalised by
$\|\ell\|=\ell^{-1}$, and put $v_\ell(x)=-\log\|x\|/\log\ell$ for $x\ne0$. Coefficient rings $\OO\subset\Qlbar$ and the
normalised lengths $\|n\|_\OO=n/e_\OO$ are as in \cref{pf:vwin-he:coef}; $\mathfrak m_\OO$ is the maximal ideal and
$k_\OO$ the residue field of $\OO$. We write $\modularitycyc:G_\Q\to\Zl^\times$ for the cyclotomic character, also
viewed with values in $\OO^\times$ or $\Qlbar^\times$, and similarly for its restrictions to the groups $G_F$. For a
Dirichlet character $\psi$ modulo $N$ with values in $\OO$, $\psi_G=\psi\circ\kappa_N:G_\Q\to\OO^\times$ is the
associated Galois character (\cref{pf:vwin-he:V}), and for $e\ge0$ we put
\begin{equation}\label{pf:modularity:Psi}
\Psi_{e,\psi}=\modularitycyc^{e}\,\psi_G:G_\Q\to\OO^\times .
\end{equation}
For a continuous character $\psi:G_\Q\to\OO^\times$, $e\ge0$ and $n\ge1$, let
$\modularitySha_n(\psi,e)=\len_\OO\Sha_n(\Q,\OO\otimes\modularityQZ(\psi^{-1}\modularitycyc^{-e}))\in\N\cup\{\infty\}$,
with $\Sha_n$ as in \cref{pf:vwin-he:V}. Thus $V_\ell(\theta,j)=\modularitySha_{\ell N}(\theta_G,j)$ for a Dirichlet
character $\theta$ modulo $N$. Two representations $r,r':\Gamma\to\GL_n(k)$ are \emph{isomorphic} if
$r'=uru^{-1}$ for some $u\in\GL_n(k)$; for a character $\chi$, $r\otimes\chi$ is the twist $g\mapsto\chi(g)r(g)$.
A \emph{complete local $\OO$-algebra} is a Noetherian local $\OO$-algebra $R$, complete for its maximal ideal, such
that $\OO\to R$ is local and $R=\OO+\mathfrak m_R$. Hosts $h=(\chi,m,q_1,d)$, the sets $\mathcal X_h(\OO,j)$ and the
cases are those of \cref{pf:vwin-he:host,pf:vwin-he:cases}.

\begin{definition}[Automorphy]\label{pf:modularity:aut}
Let $F$ be a CM field, $n\ge1$, $r:G_F\to\GL_n(\Qlbar)$ a homomorphism and $\iota:\Qlbar\to\C$ an isomorphism.
\begin{enumerate}[label=(\alph*)]
\item $r$ is \emph{attached} via $\iota$ to a cuspidal automorphic representation $\pi$ of $\GL_n(\A_F)$ if $r$ is
semisimple and for all but finitely many finite places $w$ of $F$ there is a Satake parameter $B$ of $\pi$ at $w$ such
that $r$ is unramified at $w$ and $\iota$ carries the characteristic polynomial of $r(\Frob_w^{-1})$ to
$\prod_{b}(X-q_w^{(n-1)/2}b)$, the product over $b\in B$.
\item For a place $v\mid\ell$ of $F$, $r$ is \emph{unlinked} at $v$ if for every valuation ring $A$ of $\Omega$ over $v$
the restriction $r|_{D_A}$ stabilises a full flag whose graded characters $\kappa_1,\dots,\kappa_n:D_A\to\Qlbar^\times$
satisfy $\kappa_i\ne\kappa_j\modularitycyc$ for all $i,j$; it is \emph{potentially unlinked} at $v$ if for every such
$A$ there is such a flag with $\kappa_i^m\ne(\kappa_j\modularitycyc)^m$ for all $i,j$ and $m\ge1$. A cuspidal
automorphic representation $\pi$ of $\GL_n(\A_F)$ satisfies the \emph{window condition with centre $t\in\R$} at a finite
place $v$ if the local component of $\pi$ at $v$ has cuspidal support consisting of characters of $F_v^\times$ (it is a
constituent of a principal series), and the absolute values of these characters at a uniformiser lie strictly between
$q_v^{t-1/2}$ and $q_v^{t+1/2}$. The representation $r$ satisfies the \emph{window condition} $\sfname{WitA}(\iota)$ if
(i) $r$ is potentially unlinked at every place $v\mid\ell$ at which it is unlinked; and (ii) for every regular algebraic
cuspidal $\pi$ attached to $r$ there is $t\in\R$ such that $\pi$ satisfies the window condition with centre $t$ at every
finite place $v$ which does not divide $\ell$ or at which $r$ is potentially unlinked.
\item $r$ is \emph{conjugate self-dual automorphic} via $\iota$ if it is attached via $\iota$ to a regular algebraic
cuspidal $\pi$ with $\pi^\vee\cong\pi^c$, and, if $r$ is continuous, $r$ satisfies $\sfname{WitA}(\iota)$.
\item For $F=\Q$, $r$ is \emph{automorphic over imaginary quadratic fields} via $\iota$ if, when $r$ is continuous,
for every imaginary quadratic field $K$ such that $r|_{G_K}$ is irreducible, $r|_{G_K}$ is attached via $\iota$ to a
regular algebraic cuspidal representation of $\GL_n(\A_K)$ and satisfies $\sfname{WitA}(\iota)$.
\item For $F=\Q$ and $n=2$, $r$ satisfies $\sfname{WitAQ}(\iota)$ if $r|_{G_K}$ satisfies $\sfname{WitA}(\iota)$ for
every imaginary quadratic field $K$ such that $r|_{G_K}$ is irreducible.
\end{enumerate}
\formal{UpperTheory.V.Af.AttachedSys}
\formal{UpperTheory.V.Z.Win.Unl1}
\formal{UpperTheory.V.Z.Win.PotUnl}
\formal{UpperTheory.V.Z.ImaginaryQuadraticEtaCharacter.WinAt}
\formal{UpperTheory.V.Z.Win.WitA}
\formal{UpperTheory.V.Z.Win.IsAutCsdA}
\formal{UpperTheory.V.Z.Win.IsAutQA}
\formallabel{UpperTheory.V.Z.Win.WitAQ}
\end{definition}

Condition (e) is the window condition $\sfname{WitAQ}$ of \cref{def:window:kato}. The notions of
\cref{pf:modularity:aut} are invariant under isomorphism of $r$; we record the case used below.

\begin{lemma}\label{pf:modularity:iso}
Let $F$ be a CM field, $\iota:\Qlbar\cong\C$, and $r\cong r':G_F\to\GL_n(\Qlbar)$. If $r$ is conjugate self-dual
automorphic via $\iota$, so is $r'$.
\formallabel{UpperTheory.GalN.V.autCsd_iso_alt}
\end{lemma}

\begin{proof}
Write $r'=uru^{-1}$. Being attached to $\pi$ only involves semisimplicity and characteristic polynomials, which are
invariant under conjugation (\lref{UpperTheory.GalN.V.attC_iso} states this for attachment). If $r'$ is
continuous, so is $r=u^{-1}r'u$, so $r$ satisfies $\sfname{WitA}(\iota)$. A flag of $r|_{D_v}$ is carried by $u$ to a
flag of $r'|_{D_v}$ with the same graded characters, and conversely, so condition (i) transfers; condition (ii) only
involves the representations $\pi$ attached to $r$, which are those attached to $r'$. Hence $r'$ satisfies
$\sfname{WitA}(\iota)$.
\end{proof}

\begin{lemma}\label{pf:modularity:isirr}
Let $\Gamma,\Gamma'$ be groups, $k$ a field, and $r:\Gamma\to\GL_n(k)$, $r':\Gamma'\to\GL_n(k)$ homomorphisms such that
every $r'(g')$, $g'\in\Gamma'$, is of the form $r(g)$ with $g\in\Gamma$. If $r'$ is irreducible, so is $r$. (Here a
representation is irreducible if $k^n\ne0$ and $0$ and $k^n$ are the only subspaces stable under its image.)
\formallabel{GlobalTheory.Ord2.IrrTransfer.isIrr_of_range}
\end{lemma}

\begin{proof}
$k^n\ne0$ since $r'$ is irreducible. A subspace stable under all $r(g)$ is stable under all $r'(g')$, as these are among
the $r(g)$; hence it is $0$ or $k^n$.
\end{proof}

\begin{lemma}\label{pf:modularity:mkcsd}
Let $F$ be a CM field, $\iota:\Qlbar\cong\C$ and $r:G_F\to\GL_n(\Qlbar)$. If $r$ is attached via $\iota$ to a regular
algebraic cuspidal $\pi$ with $\pi^\vee\cong\pi^c$, and $r$ satisfies $\sfname{WitA}(\iota)$ provided it is continuous,
then $r$ is conjugate self-dual automorphic via $\iota$.
\formallabel{UpperTheory.V.Z.Win.Birth.mkCsdW'}
\end{lemma}

\begin{proof}
This is \cref{pf:modularity:aut}(c).
\end{proof}

\subsection{The automorphy statements}\label{sec:pf:modularity:stmts}

\subsubsection{Symplectic data}

\begin{definition}[Raw data and their points]\label{pf:modularity:raw}
Let $\ell$ be a prime and $\OO$ a coefficient ring for $\ell$. A \emph{raw datum} $D$ over $\OO$ consists of a finite
set $S$ of natural numbers, an integer $e\ge0$, homomorphisms $\psi,\xi,\omega:G_\Q\to\OO^\times$, elements
$c,\varphi,g_0\in G_\Q$, a valuation ring $B_p$ of $\Omega$ and an element $\sigma_p\in G_\Q$ for every $p\in\N$, and an
element $s_0\in D_{B_\ell}$. Put
\[
 \Psi=\modularitycyc^e\psi,\quad \alpha=\modularitycyc\,\omega^{-1}\xi,\quad \beta=\Psi\alpha^{-1},\quad
 \theta=\alpha\Psi,
\]
\[
 r_0=\operatorname{diag}(\theta,\Psi,\alpha,1):G_\Q\to\GL_4(\OO),
\]
the \emph{Eisenstein representation}, with Hodge--Tate weights $(e+1,e,1,0)$; let $\bar r_0$ be its reduction modulo
$\mathfrak m_\OO$. Index the basis by $0,1,2,3$ and give the index $i$ the weight $w_i=(\lambda_i,\mu_i)\in\Z^2$, with
$(w_0,w_1,w_2,w_3)=((1,1),(0,1),(1,0),(0,0))$; the entry $(i,j)$ has weight $w_i-w_j$. Let
$S_b=\{p\in S:p\ne\ell,\ \ell\mid\#\psi(I_{B_p})\}$ and $T=S_b\cup\{\ell\}$. A matrix $X\in M_4$ is in
$\mathfrak{gsp}$-form with multiplier $\nu$ if $X^\ast X=\nu$, where $X^\ast_{ij}=\epsilon_i\epsilon_jX_{3-j,3-i}$ for
the signs $\epsilon=(1,1,-1,-1)$; $\mathfrak{sp}_4$ is the space of $X$ with $X^\ast=-X$.

Let $R$ be a complete local $\OO$-algebra. A \emph{$D$-point} over $R$ is a continuous lift $r:G_\Q\to\GL_4(R)$ of $\bar
r_0$ such that: $r(g)$ is in $\mathfrak{gsp}$-form with multiplier $\theta(g)$ for all $g$; $r(g)=1$ for every $g$ in
an inertia group $I_A$ with $A$ not over any prime of $T$ and $\bar r_0(g)=1$; $r(\varphi)$ is block diagonal for the
partition $\{0,2\}\cup\{1,3\}$ and $h_{02}=h_{20}=0$ for $h=r(g_0)+\theta(g_0)r(g_0^{-1})$; the restriction of $r$ to
$D_{B_\ell}$ lies, as a local lift, on a component through $r_0$ of the ordinary condition (all $r(g)$ in
$\mathfrak{gsp}$-form, the characteristic polynomial of $r(s_0)$ equal to that of $r_0(s_0)$, the flag defined by the
eigenspaces of $r(s_0)$ stable, and $I_{B_\ell}$ acting on the graded pieces through the diagonal of $r_0$); and for
$p\in S_b$ the restriction to $D_{B_p}$ lies on a component through $r_0$ of the condition ``$\mathfrak{gsp}$-form and
trivial on $I_{B_p}\cap\ker\psi$''. A \emph{characteristic-zero point} of $D$ is a homomorphism
$r:G_\Q\to M_4(\Qlbar)$ of the form $y\circ r_R$ for a $D$-point $r_R$ over some $R$ and an $\OO$-algebra map
$y:R\to\Qlbar$; it is \emph{irreducible} if the $\Qlbar$-span of $r(G_\Q)$ is $M_4(\Qlbar)$. Two characteristic-zero
points are \emph{conjugate} if they are isomorphic. For $g\in G_\Q$ put
\[
 a_g(r)=\tr r(g)-\tr r_0(g).
\]
\formal{UpperTheory.V.GradeZeroInterpolatingElement.Raw}
\formal{UpperTheory.V.GradeZeroInterpolatingElement.Raw.IsPt}
\formal{UpperTheory.V.GradeZeroInterpolatingElement.Raw.EPt}
\formallabel{UpperTheory.V.GradeZeroInterpolatingElement.Raw.aVal}
\end{definition}

\begin{definition}[Good data, and the conditions on points]\label{pf:modularity:good}
A raw datum $D$ is \emph{good} if: $\ell$ is odd; $S$ consists of primes and contains $2$ and $\ell$; $e\ge2$; $B_p$
lies over $p$ for every prime $p$; $c$ is a complex conjugation; $\psi$ and $\xi$ have open kernels and are unramified
outside $S$; $\Psi(c)=1$; $\omega^{\ell-1}=1$ and $\omega\equiv\modularitycyc\pmod{\mathfrak m_\OO}$; $\xi$ has order
prime to $\ell$, is trivial on $I_{B_v}$ for $v\in T$, and $\xi(c)=-1$; $\sigma_v$ is a Frobenius element at $B_v$ for
$v\in T$ and $\modularitycyc(\sigma_\ell)=1$; $(\xi\Psi^i)(\sigma_v)\not\equiv1\pmod{\mathfrak m_\OO}$ for $v\in T$ and
$i\in\Z$; for each $\gamma\in\{\alpha,\beta,\theta\}$ and $\gamma^{-1}$, every locally constant $1$-cocycle of $G_\Q$
with values in $k_\OO(\bar\gamma)$ which vanishes on the inertia elements away from $T$ acting trivially through
$\bar r_0$ and is a coboundary on $D_{B_v}$ for $v\in T$ is a coboundary; $\varphi=\lim_n\sigma_\ell^{\ell^{n!}}$;
$\bar\alpha(g_0)\ne1$ and $\bar\Psi(g_0)\ne1$; and $s_0\in I_{B_\ell}\cap\ker\psi\cap\ker\omega$ with
$\modularitycyc(s_0)$ of infinite order. For a good datum we use the conditions:
\begin{itemize}
\item $\sfname{Fin}_0(D)$: for every $D$-point $r$ over a complete local $\OO$-algebra $R$, the $\OO$-subalgebra of $R$
generated by the products of matrix entries $r(g_1)_{i_1j_1}\cdots r(g_k)_{i_kj_k}$ of total weight
$\sum(w_{i_\nu}-w_{j_\nu})=0$ is a finite $\OO$-module;
\item $\sfname{NoEnd}(D)$: for every characteristic-zero point $r$ all of whose values are block diagonal for
$\{0,2\}\cup\{1,3\}$, the $\Qlbar$-span of the blocks $(r(g)_{ij})_{i,j\in\{0,2\}}$ is not all of $M_2(\Qlbar)$;
\item $\sfname{Rigid}(D,r)$, for a characteristic-zero point $r$: for every point $r'$ of $D$ with values in the dual
numbers $\Qlbar[\epsilon]$ reducing to $r$, there is $X\in\mathfrak{sp}_4(\Qlbar)$ with
$r'=(1+\epsilon X)r(1-\epsilon X)$.
\end{itemize}
Finally let $\modularitySha(D)=\modularitySha_n(\psi,e)$ with $n=\prod_{p\in S}p$. For $g\in G_\Q$, $\sfname{Ineq}(D,g)$
is the statement: \emph{$\modularitySha(D)<\infty$, and there is a finite set $X$ of pairwise non-conjugate irreducible
characteristic-zero points of $D$ such that}
\begin{equation}\label{pf:modularity:ineqdef}
\|\modularitySha(D)\|_\OO\log\ell\ \le\ 2\sum_{r\in X}-\log\|a_g(r)\|\ =\ 2\log\ell\sum_{r\in X}v_\ell(a_g(r)).
\end{equation}
\formal{UpperTheory.V.GradeZeroInterpolatingElement.Raw.Good}
\formal{UpperTheory.V.GradeZeroInterpolatingElement.Raw.Fin0}
\formal{UpperTheory.V.GradeZeroInterpolatingElement.Raw.NoEnd}
\formal{UpperTheory.V.GradeZeroInterpolatingElement.Raw.RigidAt}
\formal{UpperTheory.V.GradeZeroInterpolatingElement.Raw.Ineq}
\formallabel{UpperTheory.V.GradeZeroInterpolatingElement.Raw.sha}
\end{definition}

The raw data relevant to a character of a host are those ``of the case'':

\begin{definition}[Data of the case]\label{pf:modularity:ofcase}
Let $h$ be a host with character $\chi$, $\OO$ a coefficient ring, $e\ge1$ and $\psi$ a Dirichlet character modulo
$N_h$ with values in $\OO$. For $n',p_1,m_0\in\N$, a raw datum $D$ is \emph{of character} $(e,\psi,n',p_1,m_0)$ if its
twist is $e$, its character $\psi$ is $\psi_G$, $S$ is the set of prime divisors of $2\ell N_hp_1$, and
$\xi=\xi_{0,G}\,\xi'_G$ with $\xi'$ a primitive odd Dirichlet character modulo $2^{n'}$ and $\xi_0$ a Dirichlet
character modulo $p_1$ of order $m_0$. It is \emph{of the case} of $(h,e,\psi)$ if either $(\ell,e,\psi)$ is in case
II and $D$ is of character $(e,\psi,n_{\mathrm{II}},1,1)$, or $(\ell,e,\psi)$ is in case III and $D$ is of character
$(e,\psi,n_{\mathrm{III}},p_1(\chi),\ord\psi)$. Here
\[
 n_{\mathrm{II}}=\bigl\lceil\log_2\max\bigl(\ell N_h4^{n_0(\ell,\psi)},\,2^{c_2(\psi)+6}o_\chi\ell^3q^2\bigr)\bigr\rceil,\qquad
 n_{\mathrm{III}}=\bigl\lceil\log_2\max\bigl(\ell q\,p_1(\chi)4^{n_0(\ell,\psi)},\,8\ell^2\bigr)\bigr\rceil,
\]
$q=N_\chi$ is the conductor of $\chi$, $n_0$ and $c_2$ are as in \cref{pf:vwin-he:cases}, and $p_1(\chi)$ is the least
prime $p\equiv1$ modulo the odd part of $o_\chi$ with $p\nmid2N_\chi p_K(\chi)$.
\formal{UpperTheory.V.Aut.OfChar}
\formal{UpperTheory.V.Aut.OfCase}
\formal{UpperTheory.V.HV.Host.nII}
\formal{UpperTheory.V.HV.Host.nIII}
\end{definition}

\begin{definition}[Symplectic automorphy $(\mathrm{Sp})$]\label{pf:modularity:Sp}
$(\mathrm{Sp})$ is the statement: for every $\Omega$, every host $h$ big for $t$, every odd prime $\ell$, every
coefficient ring $\OO$, every $0<e<t$, every $\psi\in\mathcal X_h(\OO,e)$ and every good raw datum $D$ over $\OO$ of the
case of $(h,e,\psi)$, the conditions $\sfname{Fin}_0(D)$ and $\sfname{NoEnd}(D)$ hold and there is an isomorphism
$\iota:\Qlbar\cong\C$ such that every irreducible characteristic-zero point of $D$ is automorphic over imaginary
quadratic fields via $\iota$ (\cref{pf:modularity:aut}(d)).
\formallabel{UpperTheory.V.Aut.SymplecticFourImageDense}
\formal{UpperTheory.V.Aut.IsAut4Alt}
\end{definition}

\subsubsection{Unitary data}

The three-dimensional deformation problems take values in the group
$\modularityG(R)=\GL_3(R)\rtimes\{\pm1\}$, where $-1$ acts on $\GL_3(R)$ by the involution
$\Theta(g)=J\,{}^tg^{-1}J$, $J$ being the antidiagonal matrix with entries $1$.
For $x\in\modularityG(R)$ we write $x=(x^\flat,\epsilon(x))$ (\emph{matrix part} and \emph{sign}). The index $i\in\{0,1,2\}$
has weight $1-i$, the entry $(a,b)$ has weight $(1-a)-(1-b)=b-a$, and for a unit $t$ we put
$\mathrm{tor}(t)=\operatorname{diag}(t,1,t^{-1})$. For a character $\eta:G_\Q\to\{\pm1\}$ and
$\Psi:G_\Q\to\OO^\times$ let $r_0^{\eta,\Psi}(g)=(\mathrm{tor}(\Psi(g)),\eta(g))$. A homomorphism
$r:G_\Q\to\modularityG(R)$ has \emph{multiplier} $\eta$ if $\epsilon\circ r=\eta$. It is of \emph{type I} if its matrix
part is irreducible on $\ker\eta$, and of \emph{type II} if on $\ker\eta$ the entries of odd weight vanish (so that
$r^\flat$ is block diagonal for $\{0,2\}\cup\{1\}$) and the $2\times2$ \emph{corner} $(r^\flat_{ab})_{a,b\in\{0,2\}}$ is
irreducible on $\ker\eta$.

\begin{definition}[Unitary data and ordinary points]\label{pf:modularity:dat}
Let $K\subset\Omega$ be an imaginary quadratic field (more precisely a CM field with an embedding $K\to\Omega$), $\ell$
a prime, $\OO$ a coefficient ring, $N\ge1$, $e\ge0$ and $\psi$ a Dirichlet character modulo $N$ with values in $\OO$.
A \emph{unitary datum} $D$ for $(K,\ell,\OO,e,\psi)$ consists of: the good-data conditions
$\sfname{GD}(K,\ell,e,\psi)$ of \cref{sec:window:eis} ($\ell$ odd, $e\ge1$, $N$ even, $\ell\nmid\ord\psi$, the
$\ell$-component of $\psi$ has order at most $2$, $\psi(-1)=(-1)^e$, $K$ is imaginary quadratic with
$\ell\nmid d_Kh_K$ and $(d_K,N)=1$, the intersection of the kernels of $\modularitycyc\bmod\ell$ and of
$\psi^\flat_G$ is not contained in $G_K$, and $(\ell-1)\nmid2e$); the quadratic character
$\eta=\eta_D$ with $\ker\eta=G_K$; a complex conjugation $c$; an embedding $j:\Omega\to\Qlbar$ and a continuous embedding
$\iota_\ell:\Gal(\Qlbar/\Ql)\to G_\Q$ compatible with $j$; a finite extension $K_v/\Ql$ with
$\eta(\iota_\ell(g))=1\iff g\in G_{K_v}$; and the conditions that $\eta$ is unramified at $\ell$, that
$\Psi=\Psi_{e,\psi}$ is continuous, and that on the inertia group at $\ell$, $\Psi(\iota_\ell(g))$ is
$\modularitycyc(g)^e$ times the $\ell$-component of $\psi_G$. Let $B$ be a valuation ring over $\ell$ and $g_1\in G_\Q$.
\begin{enumerate}[label=(\alph*)]
\item $g_1$ is \emph{marked} if $c^2=1$, $\eta(c)=-1$, $\Psi(c)=1$, $\eta(g_1)=1$ and $\Psi(g_1)^2-1\in\OO^\times$.
\item Let $A$ be a local $\OO$-algebra. An \emph{ordinary point} over $A$ (for $(\eta,\Psi,c,g_1,B)$) is a homomorphism
$r:G_\Q\to\modularityG(A)$ with multiplier $\eta$, congruent to $r_0^{\eta,\Psi}$ modulo $\mathfrak m_A$, continuous for the
$\mathfrak m_A$-adic topology, \emph{normalised} ($r(c)^\flat=1$ and the entries $(0,1)$ and $(1,0)$ of
$r(g_1)^\flat-r(g_1^\ast)^\flat$ vanish, where $g_1^\ast$ is a fixed element built from $c$ and $g_1$), trivial on the
elements of $\ker\eta\cap\ker\Psi$ in the inertia groups
of valuation rings not over $\ell$, and \emph{ordinary at $B$}: on $D_B$ it is conjugate, by a matrix congruent to $1$,
to an upper triangular representation on which $I_B\cap\ker\eta$ acts on the diagonal through
$(\Psi,1,\Psi^{-1})$.
\end{enumerate}
\formal{UpperTheory.V.U3.Dat}
\formal{UpperTheory.V.HV.GD}
\formal{UpperTheory.V.U3.Marks}
\formal{UpperTheory.V.U3.Normalised}
\formal{UpperTheory.V.Pr.OrdPt}
\end{definition}

\begin{definition}[Unitary automorphy at the ordinary points $(\mathrm{U}_{\mathrm{ord}})$]\label{pf:modularity:Uord}
Let $D$ be a unitary datum, $\iota:\Qlbar\cong\C$, $B$ a valuation ring over $\ell$ and $g_1\in G_\Q$. Consider:
\begin{itemize}
\item $\sfname{FinEnd}(g_1,B)$: there is a subgroup $U\subset G_\Q$ of finite index such that for every complete
discrete valuation ring $A$ over $\OO$ with finite residue field in which $\mathfrak m_\OO$ maps to $0$, every ordinary
point $r$ over $A$ and every $g\in U\cap\ker\eta$, the characteristic polynomial of $r(g)^\flat$ equals that of
$r_0^{\eta,\Psi}(g)^\flat$;
\item $\sfname{AutEnd}_3(g_1,B,\iota)$: for every discrete valuation ring $A\subset\Qlbar$ finite over $\OO$ and every
ordinary point $r$ over $A$ whose matrix part is irreducible on every open subgroup of $\ker\eta$, the representation
$\rho=r^\flat|_{G_K}:G_K\to\GL_3(\Qlbar)$ is such that $\rho\otimes\modularitycyc^{-1}$ is conjugate self-dual
automorphic via $\iota$;
\item $\sfname{AutEnd}_2(g_1,B,\iota)$: for every such $A$ and every ordinary point $r$ over $A$ of type II with
$r(g)^\flat_{11}=1$ for $g\in\ker\eta$, and every $W:G_\Q\to\GL_2(\Qlbar)$ which agrees with the corner of $r$ on
$\ker\eta$ and satisfies $\tr W(c)=0$: $W\otimes\modularitycyc^{-e}$ is attached via $\iota$ to a newform of weight
$2e+1$, and $W$ satisfies $\sfname{WitAQ}(\iota)$.
\end{itemize}
$(\mathrm{U}_{\mathrm{ord}})$ is the statement: for every $\Omega$, every CM field $K$ with $K\to\Omega$, every prime
$\ell$, every $\OO,N,e,\psi$, every unitary datum $D$, every $\iota$ and every valuation ring $B$ over $\ell$, there is
a marked $g_1$ satisfying $\sfname{FinEnd}(g_1,B)$, $\sfname{AutEnd}_3(g_1,B,\iota)$ and $\sfname{AutEnd}_2(g_1,B,\iota)$.
\formallabel{UpperTheory.V.Pr.GdEndAlt}
\formallabel{UpperTheory.V.Pr.AutEnd3Alt}
\formallabel{UpperTheory.V.Z.Win.Dn.AutEnd2Ao}
\formal{UpperTheory.V.Pr.FinEnd}
\end{definition}

For $(\mathrm{U}_{\mathrm{ur}})$ we need the universal objects. For a unitary datum $D$, a universal lifting ring $U$
of $r_0^{\eta,\Psi}|_{D_\ell}$ (without local condition) and $g_1\in G_\Q$, the formal development defines a
\emph{universal normalised deformation} $S=(R_S,r_S)$ of $r_0^{\eta,\Psi}$ relative to $U$ and $g_1$: $R_S$ is a complete
local $\OO$-algebra and $r_S:G_\Q\to\modularityG(R_S)$ is universal among normalised lifts which are trivial on the
inertia elements away from $\ell$ acting trivially residually and whose restriction to the decomposition group at $\ell$
satisfies the box condition defined by $U$ \cite{LeanSource2026}. The
\emph{weight-zero subalgebra}
of $S$ is the $\OO$-subalgebra of $R_S$ generated by the products of entries of $r_S^\flat$ of total weight $0$; the
\emph{entry algebra} is the $\OO$-subalgebra generated by all entries. A \emph{point} of $S$ is a pair $(\xi,r)$ of an
$\OO$-algebra map $\xi$ from the entry algebra to $\Qlbar$ and $r=\xi_\ast r_S:G_\Q\to\modularityG(\Qlbar)$.

\begin{definition}[Unitary automorphy with unramified coefficients $(\mathrm{U}_{\mathrm{ur}})$]\label{pf:modularity:Uur}
For a unitary datum $D$, $\iota$ and a marked $g_1$, consider:
\begin{itemize}
\item $\sfname{FinA}(g_1)$: for every $U$ and every universal normalised deformation $S$ relative to $U$ and $g_1$, the
weight-zero subalgebra of $S$ is a finite $\OO$-module;
\item $\sfname{AutB}(g_1,\iota)$: for every $U$, $S$ and every point $(\xi,r)$ of $S$ of type I, the representation
$r^\flat|_{G_K}\otimes\modularitycyc^{-1}$ is conjugate self-dual automorphic via $\iota$;
\item $\sfname{LemW}(g_1,\iota)$: for every $U$, $S$, every point $(\xi,r)$ of $S$ of type II with $r(g)^\flat_{11}=1$
for $g\in\ker\eta$, and every $W:G_\Q\to\GL_2(\Qlbar)$ agreeing with the corner of $r$ on $\ker\eta$ with
$\tr W(c)=0$: $W\otimes\modularitycyc^{-e}$ is attached via $\iota$ to a newform of weight $2e+1$ and $W$ satisfies
$\sfname{WitAQ}(\iota)$ (we then say that $W$ is \emph{modular}).
\end{itemize}
$(\mathrm{U}_{\mathrm{ur}})$ is the statement: for every $\Omega$, CM field $K$, prime $\ell$, $\OO$, $N$, $e$, $\psi$
and unitary datum $D$ with $e_\OO=1$, and every $\iota$, there is a marked $g_1$ satisfying $\sfname{FinA}(g_1)$,
$\sfname{AutB}(g_1,\iota)$ and $\sfname{LemW}(g_1,\iota)$.
\formallabel{UpperTheory.V.Aut.AutomorphyUnramifiedCoefficientsAlt}
\formallabel{UpperTheory.V.Aut.AutBAlt}
\formallabel{UpperTheory.V.Aut.LemWAlt}
\formallabel{UpperTheory.V.Z.Win.Dn.IsModularA}
\formal{UpperTheory.V.Aut.FinA}
\end{definition}

\begin{definition}[Unitary automorphy over $\Q(i)$, $(\mathrm{U}_i)$]\label{pf:modularity:Ui}
$(\mathrm{U}_i)$ is the statement: for every imaginary quadratic field $K$ containing a square root of $-1$ (that is,
$K=\Q(i)$), every $\Omega$, prime $\ell$, $\iota:\Qlbar\cong\C$, coefficient ring $\OO$, $N$, $e$, $\psi$ and $o_1$
satisfying the host-data conditions ($\ell$ odd, $e\ge1$, $N$ even, $\psi(-1)=(-1)^e$, $o_1\ge1$, and conditions
relating $o_1$ to the $2$-adic conductor exponent $c_2(\psi)\ge4$), every finite set $S$ of primes with chosen
valuation rings over them, and every host frame $F$ (a framed deformation problem for three-dimensional
representations of $G_{\Q,S}$ with values in $\modularityG$, with multiplier the quadratic character of $\Q(i)$ and
residual representation $r_0^{\eta,\Psi_{e,\psi}}$, $S$ being the set of primes dividing $\ell N$), there is a marked
element $g_1\in G_{\Q,S}$ such that for every
complete local $\OO$-algebra $R$ and every universal framed deformation $u:G_{\Q,S}\to\modularityG(R)$ of $F$ at $g_1$:
the part of weight $0$ of $R$ (for the torus action) is a finite $\OO$-module, and for every $\OO$-algebra map $\xi$
from the subalgebra of $R$ generated by the homogeneous parts to $\Qlbar$ such that $r=\xi_\ast u$ is of type I, the
representation $r^\flat|_{G_K}\otimes\modularitycyc^{-1}$ is conjugate self-dual automorphic via $\iota$.
\formallabel{UpperTheory.V.Aut.ThmS6Core_s}
\formallabel{UpperTheory.V.Af.AutOnAlt}
\formal{UpperTheory.V.UnitaryThreeInertia.HostData}
\formal{UpperTheory.V.UnitaryThreeInertia.IsHostFrame}
\end{definition}

\subsection{From ordinary points to unramified coefficients}\label{sec:pf:modularity:unr}

In this subsection we prove step \ref{pf:vwin-he:e} of \cref{sec:pf:vwin-he:proof}. The only arithmetic input besides
$(\mathrm{U}_{\mathrm{ord}})$ is the package of \emph{local theorems}: local finiteness, local Tate duality and the
local Euler characteristic formula for finite extensions of $\Ql$, and the local theory of Bloch and Kato
\cite{BlochKato1990}. These are components of \ref{pf:vwin-he:H1}(c):

\begin{lemma}\label{pf:modularity:loc}
\ref{pf:vwin-he:H1} implies: (a) the local theorems; (b) Poitou--Tate duality; (c) the package $\sfname{CohP}$ (local
finiteness, local Tate duality and local Euler characteristics for $\Ql$, global finiteness, the global Euler
characteristic formula, the Greenberg--Wiles formula, normalised Poitou--Tate duality, and tame local duality); (d) the
global compatibility package (the local statements of (c), global finiteness, global Euler characteristics, normalised
Poitou--Tate duality, unramified orthogonality and the Greenberg--Wiles formula); (e) local class field theory at all
places.
\formallabel{UpperTheory.V.HV.U.locPrints_of_litN}
\formallabel{UpperTheory.V.W.locPrints_of_litN}
\formallabel{UpperTheory.V.HV.U.pt_s}
\formallabel{UpperTheory.V.W.cohP_alt}
\formallabel{UpperTheory.V.W.global_compat_printed_of_literal_automorphy_n}
\formallabel{UpperTheory.V.W.local_class_field_theory_all_of_literal_automorphy}
\formal{UpperTheory.V.Loc.LocalPublishedTheorems}
\end{lemma}

\begin{proof}
The statements in (a), (c) and (d) are among the Galois-cohomological components of
$\sfname{LiteratureRankNHypothesis}$ listed in \ref{pf:vwin-he:H1}(c), except tame local duality, which is deduced from
local Tate duality (\lref{GlobalTheory.Published.TameLocal_of_duality}). Poitou--Tate duality in the form used
is deduced from its normalised form (\lref{GlobalTheory.PoitouTateConj.poitouTateDuality_of_normalised}). Local
class field theory at all places is deduced from global Artin reciprocity, a component of \ref{pf:vwin-he:H1}(b), and
(c) (\lref{UpperTheory.V.W.locCFT}).
\end{proof}

Fix a unitary datum $D$ for $(K,\ell,\OO,e,\psi)$ and let $B$ be the valuation ring
$\{x\in\Omega:\|j(x)\|\le1\}$ defined by the embedding $j$ of $D$. By
\lref{UpperTheory.V.U3.dat_loc}, $B$ lies over $\ell$, the decomposition group $D_B$ is the image of
$\iota_\ell$, and $I_B$ is the image of the inertia group of $\Gal(\Qlbar/\Ql)$.

\begin{lemma}[Finiteness from finite endomorphisms]\label{pf:modularity:finA}
Assume the local theorems. Let $g_1$ be marked and assume $\sfname{FinEnd}(g_1,B)$. Then $\sfname{FinA}(g_1)$ holds.
\formallabel{UpperTheory.V.Pr.finA_of_end}
\end{lemma}

\begin{proof}
Let $U$ be a universal local lifting ring and $S=(R_S,r_S)$ a universal normalised deformation relative to $U$ and
$g_1$. We use three auxiliary results.
\begin{enumerate}[label=(\roman*)]
\item (\lref{gf_pres}) $R_S$ has a $\Z$-weighted presentation $W$ (a presentation as a quotient of a power series
ring over $\OO$ with homogeneous variables and relations, the grading coming from the action of the torus
$t\mapsto\mathrm{tor}(t)$ by conjugation on normalised lifts) whose generators $z_i$ lie in $\mathfrak m_{R_S}$ and are of
the form $r_S(g)^\flat_{ab}-r_0^{\eta,\Psi}(g)^\flat_{ab}$, of weight $b-a$; moreover every entry $r_S(g)^\flat_{ab}$ is
homogeneous of weight $b-a$. The proof uses that conjugation by $\mathrm{tor}(t)$ preserves the normalisation and the
local conditions (the local theorems are used to see that the box condition at $\ell$ is stable under conjugation,
\lref{UpperTheory.V.Loc.locGood}).
\item Since the entries are homogeneous, every product of entries of total weight $0$ lies in the part of degree $0$
of $W$, so the weight-zero subalgebra is an $\OO$-submodule of it; as $\OO$ is Noetherian, it suffices to show that
the part of degree $0$ is a finite $\OO$-module.
\item (Finiteness criterion, \lref{UpperTheory.V.Pr.finCrit}) For a $\Z$-weighted presentation $W$ of a complete
local $\OO$-algebra with generators in the maximal ideal, the part of degree $0$ is a finite $\OO$-module provided:
for every complete discrete valuation ring $A$ over $\OO$ with finite residue field in which $\mathfrak m_\OO$ maps to
$0$, and every local $\OO$-algebra map $f:R_S\to A$, $f$ kills the generators of weight $0$ and either all generators of
negative weight or all generators of positive weight.
\end{enumerate}
Let $A$ and $f$ be as in (iii). By \lref{UpperTheory.V.Pr.ordPt_of_gd} (which uses the local theorems),
$r=f_\ast r_S$ is an ordinary point over $A$. By $\sfname{FinEnd}(g_1,B)$ there is a subgroup $U'$ of finite index such
that the characteristic polynomials of $r(g)^\flat$ and $r_0^{\eta,\Psi}(g)^\flat$ agree for $g\in U'\cap\ker\eta$.
Since $2$ is a unit in $\OO$ and $g_1$ is marked, the triangularity criterion \lref{UpperTheory.V.Pr.constOfChar}
(for a normalised lift over a local domain of characteristic $\ell$ of the residual representation, with multiplier
$\eta$ and Eisenstein characteristic polynomials on $U'\cap\ker\eta$: the lift is upper or lower triangular on
$\ker\eta$ and its diagonal entries are those of the residual representation) applies to $r$, viewed over the residue
field of $\OO$. For $g\notin\ker\eta$, $gc\in\ker\eta$ and $r(g)^\flat=r(gc)^\flat$ because $r(c)^\flat=1$. Hence the
diagonal entries of $r(g)^\flat$ are those of $r_0^{\eta,\Psi}(g)^\flat$ for all $g$, and either all entries of
negative weight of $r^\flat$ vanish, or all entries of positive weight vanish. The
generators of weight $0$ are differences of diagonal entries, since $b-a=0$ only for $a=b$, so $f$ kills them; and
since $r_0^{\eta,\Psi}(g)^\flat_{ab}=0$ for $a\ne b$, $f$ kills all generators of one sign. Hence (iii) applies and
$S$ satisfies the finiteness required in $\sfname{FinA}(g_1)$.
\end{proof}

\begin{lemma}[Points come from integral ordinary points]\label{pf:modularity:ordpt}
Assume the local theorems. Let $g_1$ be marked and assume $\sfname{FinA}(g_1)$. Let $U$, $S$ be as above and
$(\xi,r)$ a point of $S$. Then there are a discrete valuation ring $A\subset\Qlbar$ finite over $\OO$, a unit
$t\in\Qlbar^\times$ and an ordinary point $r':G_\Q\to\modularityG(A)$ (for $(\eta,\Psi,c,g_1,B)$) such that, in
$\modularityG(\Qlbar)$,
\begin{equation}\label{pf:modularity:conjtor}
r'(g)=\mathrm{tor}(t)\,r(g)\,\mathrm{tor}(t)^{-1}\qquad(g\in G_\Q).
\end{equation}
\formallabel{UpperTheory.V.Pr.exists_ordPt_of_point}
\end{lemma}

\begin{proof}
By \lref{exists_pres} $R_S$ has a $\Z$-weighted presentation $W$ in which all entries $r_S(g)^\flat_{ab}$ are
homogeneous of weight $b-a$. By \lref{part_zero_le} and $\sfname{FinA}(g_1)$ the part of degree $0$ of $W$ lies in
the weight-zero subalgebra (here one uses that $\ell\in\mathfrak m_{R_S}$), so it is a finite $\OO$-module and is
contained in the entry algebra (\lref{A_le}). Hence the image under $\xi$ of the part of degree $0$ is a finitely
generated $\OO$-module. A point of a $\Z$-weighted algebra whose restriction to degree $0$ has finitely generated
image can be moved by the torus into a finite extension: by \lref{UpperTheory.V.Alg.points_finite} there are a
character $t:\Z\to\Qlbar^\times$ and a point $\Xi$ with $\Xi(x)=t(n)\xi(x)$ for $x$ of degree $n$, whose image is a
finitely generated $\OO$-module. Then (\lref{exists_range}) the image of $\Xi$ is contained in a discrete
valuation ring $A\subset\Qlbar$ finite over $\OO$, and $\Xi$ factors through an $\OO$-algebra map $f:R_S\to A$. The
point $r'=f_\ast r_S$ is ordinary by \lref{UpperTheory.V.Pr.ordPt_of_gd} (and \lref{loc_of_int}). Finally,
an entry of weight $b-a$ is multiplied by $t^{b-a}$; with $t=t(1)$ this is \eqref{pf:modularity:conjtor} for the matrix
parts, and both sides have sign $\eta(g)$.
\end{proof}

Note the direction of \eqref{pf:modularity:conjtor}: conjugation by $\mathrm{tor}(t)$ multiplies the entry $(a,b)$ by
$t^{(1-a)-(1-b)}$. For $g\in\ker\eta$ the sign is trivial, so \eqref{pf:modularity:conjtor} is an identity of matrices
$r'(g)^\flat=\mathrm{tor}(t)r(g)^\flat\mathrm{tor}(t)^{-1}$; for type II points, the corners satisfy
$\mathrm{corner}(r'(g)^\flat)=d(t)\,\mathrm{corner}(r(g)^\flat)\,d(t)^{-1}$ with $d(t)=\operatorname{diag}(t,t^{-1})$, and
the middle entries agree.

\begin{lemma}\label{pf:modularity:modend}
Let $g_1$ be marked, $A\subset\Qlbar$ a discrete valuation ring finite over $\OO$ and $r'$ an ordinary point over $A$
of type II with $r'(g)^\flat_{11}=1$ for $g\in\ker\eta$. Assume $\sfname{AutEnd}_2(g_1,B,\iota)$. If
$W':G_\Q\to\GL_2(\Qlbar)$ agrees with the corner of $r'$ on $\ker\eta$, $\tr W'(c)=0$, and $W\cong W'$, then $W$ is
modular (\cref{pf:modularity:Uur}).
\formallabel{UpperTheory.GalN.V.modOfEnd'_alt}
\end{lemma}

\begin{proof}
By $\sfname{AutEnd}_2$, $W'\otimes\modularitycyc^{-e}$ is attached to a newform $f$ of weight $2e+1$ and $W'$
satisfies $\sfname{WitAQ}(\iota)$. Attachment to $f$ and $\sfname{WitAQ}$ are invariant under isomorphism (as in
\cref{pf:modularity:iso}), and $W\otimes\modularitycyc^{-e}\cong W'\otimes\modularitycyc^{-e}$.
\end{proof}

\begin{theorem}\label{pf:modularity:unr-kisin}
The local theorems and $(\mathrm{U}_{\mathrm{ord}})$ imply $(\mathrm{U}_{\mathrm{ur}})$.
\formallabel{UpperTheory.V.Pr.automorphyUnramified_of_locPrints_kisin_gdEnd}
\end{theorem}

\begin{proof}
Let $D$ be a unitary datum for $(K,\ell,\OO,e,\psi)$, $\iota:\Qlbar\cong\C$, and $B$ as above. Apply
$(\mathrm{U}_{\mathrm{ord}})$ to $(D,\iota,B)$: there is a marked $g_1$ satisfying $\sfname{FinEnd}(g_1,B)$,
$\sfname{AutEnd}_3(g_1,B,\iota)$ and $\sfname{AutEnd}_2(g_1,B,\iota)$. By \cref{pf:modularity:finA},
$\sfname{FinA}(g_1)$ holds.

$\sfname{AutB}(g_1,\iota)$. Let $(\xi,r)$ be a point of $S$ of type I (it has multiplier $\eta$, since $r_S$ has). Let
$A,t,r'$ be as in \cref{pf:modularity:ordpt}. By \eqref{pf:modularity:conjtor}, $r'^\flat$ and $r^\flat$ are conjugate
by $\mathrm{tor}(t)$ on $\ker\eta$, so $r'^\flat$ is irreducible on $\ker\eta$, and therefore, $r'$ being an ordinary
point, on every open subgroup of $\ker\eta$ (\lref{UpperTheory.V.Pr.matIrr_open_gd}). By
$\sfname{AutEnd}_3(g_1,B,\iota)$, $r'^\flat|_{G_K}\otimes\modularitycyc^{-1}$ is conjugate self-dual automorphic via
$\iota$. Since $G_K=\ker\eta$, $r^\flat|_{G_K}\otimes\modularitycyc^{-1}$ is isomorphic to it, hence conjugate self-dual
automorphic by \cref{pf:modularity:iso}.

$\sfname{LemW}(g_1,\iota)$. Let $(\xi,r)$ be a point of $S$ of type II with $r(g)^\flat_{11}=1$ on $\ker\eta$, and $W$
agreeing with the corner of $r$ on $\ker\eta$ with $\tr W(c)=0$. Let $A,t,r'$ be as in \cref{pf:modularity:ordpt} and
put $W'=d(t)Wd(t)^{-1}$. By the remark after \cref{pf:modularity:ordpt}, $r'$ is of type II (its entries of odd weight
on $\ker\eta$ are multiples of those of $r$, and its corner is conjugate to that of $r$), $r'(g)^\flat_{11}=1$ on
$\ker\eta$, $W'$ agrees with the corner of $r'$ on $\ker\eta$, and $\tr W'(c)=\tr W(c)=0$. Since $W\cong W'$,
\cref{pf:modularity:modend} shows that $W$ is modular.
\end{proof}

\begin{remark}\label{pf:modularity:rem-unr}
The hypothesis $e_\OO=1$ in $(\mathrm{U}_{\mathrm{ur}})$ is not used in the proof: the argument gives the conclusion of
$(\mathrm{U}_{\mathrm{ur}})$ for every coefficient ring. Also, the name of the formal theorem mentions a Kisin-type
input, but its only hypotheses are the local theorems and $(\mathrm{U}_{\mathrm{ord}})$.
\end{remark}

\begin{theorem}[Step \ref{pf:vwin-he:e} of \cref{sec:pf:vwin-he:proof}]\label{pf:modularity:unr-main}
\ref{pf:vwin-he:H1} and $(\mathrm{U}_{\mathrm{ord}})$ imply $(\mathrm{U}_{\mathrm{ur}})$.
\formallabel{UpperTheory.V.HV.Core0.automorphyUnramified_of_litN_of_gdEnd}
\end{theorem}

\begin{proof}
Combine \cref{pf:modularity:loc} and \cref{pf:modularity:unr-kisin}.
\end{proof}

\subsection{Cases II and III}\label{sec:pf:modularity:count23}

In this subsection we prove step \ref{pf:vwin-he:g} of \cref{sec:pf:vwin-he:proof}. Besides $(\mathrm{Sp})$, the proof
uses the following inputs; the first four are statements of the formal development proved elsewhere.
\begin{itemize}
\item \emph{The cohomological hypotheses}: the local theorems of \cref{sec:pf:modularity:unr}, Poitou--Tate duality (in
normalised form), global finiteness and the global Euler characteristic formula, orthogonality of unramified classes
and the Greenberg--Wiles formula \cite{Wiles1995}; all are components of \ref{pf:vwin-he:H1}(c).
\item Lemma M (\cref{pf:vwin-he:lemmaM}); Lemma M$'$ (\lref{UpperTheory.V.Alg.M.mprime_holds}), which bounds the
length of a module $M$ over a discrete valuation ring with finite residue field by the cotangent length
$\modularitycot_j(x;F)$ whenever $M$ embeds into the homogeneous point derivations of degree $j$ at $x$ modulo some
$c\ne0$ which vanish on the $F_i$; and the presentation Lemma P (\lref{UpperTheory.V.Alg.lemmaP}), which
presents a universal framed deformation ring of a residual representation as a power series ring over the local
lifting rings modulo relations, with the numbers of variables and relations of each weight given by the dimensions of
the Galois cohomology groups of the graded pieces of the adjoint representation.
\item $\sfname{CorX}$ (\lref{UpperTheory.V.X.corX}): for the residual characters arising from a Dirichlet
character of the type considered and an odd character modulo a power of $2$, every residual $1$-cocycle satisfying
the relevant unramifiedness conditions is a coboundary.
\item For every $\Omega$ and every prime $\ell$, the statements $\sfname{SoV}(\Omega,\ell)$ and
$\sfname{SoFin}(\Omega,\ell)$ of \cref{pf:modularity:so} below, which follow from Lemma K
(\lref{UpperTheory.V.X.soBar} and \lref{UpperTheory.V.X.soFin}).
\item The Eisenstein congruence number $\sfname{NumberSp4}$ and the rigidity statement $\sfname{RigSp4}$
(\cref{pf:modularity:number,pf:modularity:rigdef}); they are proved from $\sfname{HypothesesH'}$ in
\cref{pf:modularity:numberSp4,pf:modularity:rig}.
\end{itemize}

\begin{definition}\label{pf:modularity:so}
Let $\Omega$ be fixed and $\ell$ a prime.
\begin{enumerate}[label=(\alph*)]
\item $\sfname{SoV}(\Omega,\ell)$: for every continuous $\psi:G_\Q\to\Qlbar^\times$ with open kernel, every $e\ge1$, and
$\theta=\psi^{-1}\modularitycyc^{-e}$, every continuous $1$-cocycle $z:G_\Q\to\Qlbar(\theta)$ which is a coboundary on
$I_A$ for every valuation ring $A$ of $\Omega$, and on $D_A$ for every $A$ over $\ell$, is a coboundary.
\item $\sfname{SoFin}(\Omega,\ell)$: for every coefficient ring $\OO$, every continuous $\psi:G_\Q\to\OO^\times$ with
open kernel and every $e\ge1$, $\modularitySha_\ell(\psi,e)<\infty$.
\end{enumerate}
\formallabel{UpperTheory.V.GradeZeroInterpolatingElement.SoV}
\formallabel{UpperTheory.V.GradeZeroInterpolatingElement.SoFin}
\end{definition}

\begin{definition}[The Eisenstein congruence number for $\GSp_4$]\label{pf:modularity:number}
For a prime $p$, $e\in\N$ and $z_1,z_2,z_3\in\C$ put $\mathrm{eis}_4(p,e;z)=p^{e+1}z_1+p^ez_2+pz_3+1$. Let $\iota$,
$M\ge1$ and $m=(e+1,e,1,0)$. A continuous $r:G_\Q\to\GL_4(\Qlbar)$ is \emph{admissible} of level $M$ (with
$\iota$) if it is irreducible, automorphic over imaginary quadratic fields via $\iota$, of level $M$ away from $\ell$
and of type $(M,m)$ at $\ell$, preserves a non-degenerate bilinear form up to a
character $\theta=\modularitycyc^{e+1}\theta'$ with $\theta'$ trivial on $G_{\Q(\mu_{M'})}$ for some $M'$, and
satisfies the side condition $\sfname{SafeRA}$: either $\ell\mid M$ and the condition $\sfname{UnlAll}$ of
\cite{LeanSource2026} holds, or no two of the weights $m_i$ differ by $1$ (here the second alternative is excluded).
$\sfname{NumberSp4}$ is the statement: \emph{there are $c,A\in\N$ such that for all $e\ge2$, $M\ge1$, primes $p\ge5$
with $p\nmid M$, $k\ge1$ and $k$-th roots of unity $z_1,z_2,z_3\in\C$, there is an integer $\mathcal N\ne0$ with
\[
 \log|\mathcal N|\le c\,\varphi(k)\,(M(e+1))^A\log p
\]
such that for every prime $\ell\ne2,p$, every $\iota$, every Frobenius element $\sigma$ at a valuation ring over $p$,
and every finite set $X$ of pairwise non-isomorphic admissible representations of level $M$, one has
$\tr r(\sigma)\ne\iota^{-1}\mathrm{eis}_4(p,e;z)$ for $r\in X$ and
\[
 \|\mathcal N\|\le\prod_{r\in X}\bigl\|\tr r(\sigma)-\iota^{-1}\mathrm{eis}_4(p,e;z)\bigr\|.
\]}
\formallabel{UpperTheory.V.Af.NumberSp4Alt}
\formallabel{UpperTheory.V.Af.eis4}
\formallabel{UpperTheory.V.Af.AdmRAlt}
\formallabel{UpperTheory.V.Z.Win.A4.SafeRA}
\formal{UpperTheory.V.Af.Level}
\formal{UpperTheory.V.Af.DualTw}
\end{definition}

\begin{definition}[Rigidity]\label{pf:modularity:rigdef}
$\sfname{RigSp4}$ is the statement: for every $\Omega$, $\ell$, $\OO$ and every good raw datum $D$ of the case of some
$(h,e,\psi)$ in case II or III with $h$ big for some $t$ and $0<e<t$, every $N\ge1$ such that $\psi$, $\xi$ and $\omega$
are trivial on $G_{\Q(\mu_N)}$, every $\iota$, and every irreducible characteristic-zero point $r$ of $D$ which is
automorphic over imaginary quadratic fields via $\iota$, the condition $\sfname{Rigid}(D,r)$ holds.
\formallabel{UpperTheory.V.Th.RigiditySymplecticFour}
\end{definition}

\subsubsection{The statement and the reduction to large coefficients}

\begin{theorem}[Step \ref{pf:vwin-he:g} of \cref{sec:pf:vwin-he:proof}]\label{pf:modularity:count23}
$\sfname{HypothesesH'}$, the existence of $C_K$ with $\sfname{LemmaK}(C_K)$, and $(\mathrm{Sp})$ imply $\sfname{CountII}$
and $\sfname{CountIII}$.
\formallabel{UpperTheory.V.HV.DescentVanishingReductions.countII_and_countIII_of_printsH_of_sp4In}
\end{theorem}

\begin{proof}
We apply \cref{pf:modularity:sp4main} below; its hypotheses hold as follows. The cohomological hypotheses are components
of \ref{pf:vwin-he:H1} (\cref{pf:modularity:loc} and \cref{pf:modularity:adjin}). Lemma M, Lemma M$'$ and Lemma P are
theorems (\lref{UpperTheory.V.Alg.M.main_holds}, \lref{UpperTheory.V.Alg.M.mprime_holds}, and \lref{UpperTheory.V.Alg.lemmaP}).
$\sfname{NumberSp4}$ follows from $\sfname{HypothesesH'}$ by \cref{pf:modularity:numberSp4}, and $\sfname{RigSp4}$ by
\cref{pf:modularity:rig}, applied to the inputs of \cref{pf:modularity:adjin}. $\sfname{CorX}$ follows from the
residual vanishing lemma and the non-vanishing of residual Bernoulli numbers
(\lref{UpperTheory.V.X.corX}, \lref{UpperTheory.V.X.lemmaX}, and \lref{UpperTheory.V.X.residual_bernoulli_nonvanishing}).
Finally, given $C_K$ with $\sfname{LemmaK}(C_K)$, for every $\Omega$ and $\ell$ the statement $\sfname{SoV}(\Omega,\ell)$
follows from $\sfname{LemmaK}(C_K)$, Poitou--Tate duality, local finiteness and local Tate duality
(\lref{UpperTheory.V.X.soBar}), and $\sfname{SoFin}(\Omega,\ell)$ from the same inputs
(\lref{UpperTheory.V.X.soFin}); Poitou--Tate duality and the local statements are consequences of
\ref{pf:vwin-he:H1} by \cref{pf:modularity:loc}.
\end{proof}

\begin{theorem}\label{pf:modularity:sp4main}
Assume $(\mathrm{Sp})$, the cohomological hypotheses, Lemma M, Lemma M$'$, Lemma P, $\sfname{NumberSp4}$,
$\sfname{RigSp4}$, $\sfname{CorX}$, and $\sfname{SoV}(\Omega,\ell)$ and $\sfname{SoFin}(\Omega,\ell)$ for all $\Omega$ and
$\ell$. Then $\sfname{CountII}$ and $\sfname{CountIII}$ hold.
\formallabel{UpperTheory.V.Th.sp4Main_alt}
\end{theorem}

The heart of the proof is the following proposition, proved in \cref{sec:pf:modularity:mainarg}.

\begin{proposition}\label{pf:modularity:mainbig}
Under the hypotheses of \cref{pf:modularity:sp4main}, for every primitive Dirichlet character $\chi$ there is $C\in\N$
with the following property. Let $h$ be a host with character $\chi$, $t\in\N$ such that $h$ is big for $t$, $\ell$ an
odd prime, $\OO$ a coefficient ring, $0<e<t$ and $\psi\in\mathcal X_h(\OO,e)$ such that $(\ell,e,\psi)$ is in case II
or III. If $\OO$ contains the $2^{n_{\mathrm{II}}}$-th and the $2^{n_{\mathrm{III}}}$-th roots of unity
(\cref{pf:modularity:ofcase}), then $(\psi,e)$ is bounded by $(\ell N_he)^C$.
\end{proposition}

\begin{proof}[Proof of \cref{pf:modularity:sp4main}]
It suffices to prove $\mathrm{Count}(P)$ for the property $P$: ``$(\ell,e,\psi)$ is in case II or in case III''; the
two statements follow by restricting $P$. Let $\chi$ be given and $C$ as in \cref{pf:modularity:mainbig}. Let $h$, $t$,
$\Omega$, $\ell$ odd, $\OO$, $0<e<t$ and $\psi\in\mathcal X_h(\OO,e)$ in case II or III be given, and put
$n=n_{\mathrm{II}}+n_{\mathrm{III}}$ (computed for $\psi$). There is a coefficient ring $\OO_1$ containing the $2^n$-th roots of unity
(\lref{UpperTheory.V.HV.exists_coef}), and a coefficient ring $\OO'$ containing $\OO$ and $\OO_1$
(\lref{UpperTheory.V.HV.coefSup}); $\OO'$ contains the $2^{n_{\mathrm{II}}}$-th and $2^{n_{\mathrm{III}}}$-th roots
of unity. Let $f:\OO\to\OO'$ be the inclusion and $\psi'=f\circ\psi$. The invariants entering the definitions are
unchanged when $\psi$ is replaced by $\psi'$: $f$ is injective, so $\ker\psi'=\ker\psi$, $\psi'(-1)=f(\psi(-1))$, the
orders of $\psi$ and of its components (in particular of $\psi^\flat$) are unchanged, and so are the conductors of the
components; hence $n_0(\ell,\psi')=n_0(\ell,\psi)$, $c_2(\psi')=c_2(\psi)$, $n_{\mathrm{II}}$ and $n_{\mathrm{III}}$
are unchanged, $\psi'\in\mathcal X_h(\OO',e)$, and $(\ell,e,\psi')$ is in the same case as $(\ell,e,\psi)$ (also
$\psi'\ne1$ if $\psi\ne1$). By \cref{pf:modularity:mainbig}, $(\psi',e)$ is bounded by $(\ell N_he)^C$. Since $\psi$
and $\psi'$ have the same values in $\Qlbar$, \lref{UpperTheory.V.HV.coefIndep} shows that $V_\ell(\psi,e)$ is
finite if and only if $V_\ell(\psi',e)$ is, and that $\widetilde V_\ell(\psi,e)=\widetilde V_\ell(\psi',e)$ (the normalised
length does not depend on the coefficient ring). Hence $(\psi,e)$ is bounded by $(\ell N_he)^C$.
\end{proof}

\subsubsection{Models and the key inequality}\label{sec:pf:modularity:ineq}

The deformation rings of a good raw datum are studied through a \emph{model}: a presentation of the universal
$D$-point compatible with the action of the diagonal torus. Recall the weights $w_i\in\Z^2$ of
\cref{pf:modularity:raw}; put $w_\alpha=(1,0)$, $w_\beta=(-1,1)$, $w_\theta=(1,1)$ and $w_\Psi=(0,1)$, so that the entry
$(i,j)$ of weight $w_\alpha$, $w_\beta$, $w_\theta$, $w_\Psi$ is acted on by the torus as the corresponding character
of $r_0$ (for instance the entries $(0,1)$ and $(2,3)$ have weight $w_\alpha$, the entries $(0,2)$ and $(1,3)$ weight
$w_\Psi$).

\begin{definition}[Models]\label{pf:modularity:model}
Let $D$ be a good raw datum over $\OO$. A \emph{model} of $D$ consists of:
\begin{enumerate}[label=(\alph*)]
\item a complete local $\OO$-algebra $R$ which is a domain, with a $\Z^2$-weighted presentation (a presentation as a
quotient of a power series ring over $\OO$ in homogeneous variables whose weights are weights of entries, modulo
homogeneous relations), and an \emph{origin} $x:R\to\OO$ killing the variables; we write $R_w$ for the homogeneous
parts and $A=\bigoplus_{n\in\Z}R_{(0,n)}$, graded by $n$, with $x_0=x|_A$;
\item integers $r\ge1$, $r',m\ge0$ such that the origin is a smooth point of the $\Z^2$-graded algebra with tangent
dimension $1$ in each of the degrees $w_\alpha,w_\beta,w_\theta$, $r'$ in degree $w_\Psi$, $r$ in degree $-w_\Psi$, $m$
in degree $0$, and $0$ otherwise;
\item elements $F_0,\dots,F_r\in R_{-w_\Psi}$, $E_1,\dots,E_{r'}\in R_{w_\Psi}$ and $Z_1,\dots,Z_m\in R_0$ with
$x(Z_i)=0$, and a surjection $\pi:R\to S$ onto a complete local $\OO$-algebra with kernel generated by the $F_i$, $E_i$,
$Z_i$;
\item a $D$-point $r_u$ over $S$ whose entries generate $S$, such that every entry $r_u(g)_{ij}$ is the image of an
element of $R_{w_i-w_j}$, and $x(f)=r_0(g)_{ij}$ whenever $\pi(f)=r_u(g)_{ij}$;
\item for every $0\ne a\in\OO$ and every locally constant $1$-cocycle $z:G_\Q\to(\OO/a)(\Psi^{-1})$ which vanishes on
the inertia elements away from $T$ acting trivially through $\bar r_0$, vanishes at $g_0$, and is a coboundary on
$D_{B_v}$ for $v\in T$ (a \emph{$\Sha$-cocycle modulo $a$}), a point derivation $\delta:R\to\OO$ at the origin,
homogeneous of degree $-w_\Psi$, mapping $\ker\pi$ into $a\OO$, such that $\delta(f)\equiv\theta(g)z(g)$ whenever
$\pi(f)=r_u(g)_{20}$ and $\delta(f)\equiv\Psi(g)z(g)$ whenever $\pi(f)=r_u(g)_{31}$, modulo $a$.
\end{enumerate}
Then $(A,x_0)$ is a weighted ring in the sense of \cref{pf:vwin-he:wring} and $(F,E,Z)$ is a system of equations of type
$(r,0,r',0,m)$ in the sense of \cref{pf:vwin-he:eqns}, generating an ideal $I\subset A$. Let $S^\circ\subset S$ be the
subalgebra generated by the images of the homogeneous parts; it contains all entries of $r_u$.
\formallabel{UpperTheory.V.GradeZeroInterpolatingElement.Model}
\end{definition}

The torus $(\Qlbar^\times)^2$ acts on $\Qlbar$-points of $A$ and $S^\circ$; for $t=(t_1,t_2)$ we write $t\cdot r$ for
the corresponding conjugate of a representation by a diagonal matrix. The degree $n$ of $A$ is the second coordinate
of the weight; thus the $\Qlbar^\times$-orbits of \cref{pf:vwin-he:wring} are the orbits of the second factor.

\begin{lemma}\label{pf:modularity:aelt}
Let $D$ be good with a model and $g\in G_\Q$. There is $a\in A_0$ with $x_0(a)=0$ and $\xi(a)=a_g(\xi\circ r_u)$ for every
$\OO$-algebra map $\xi:S^\circ\to\Qlbar$ (where $\xi(a)$ means $\xi(\pi(a))$).
\formallabel{UpperTheory.V.GradeZeroInterpolatingElement.aElt}
\end{lemma}

\begin{proof}
For $i=0,\dots,3$ choose $f_i\in R_{w_i-w_i}=R_0$ with $\pi(f_i)=r_u(g)_{ii}$, and put
$a=\sum_if_i-\tr r_0(g)\in A_0$. Then $x_0(a)=\sum_ir_0(g)_{ii}-\tr r_0(g)=0$ by (d), and
$\xi(a)=\sum_i\xi(r_u(g)_{ii})-\tr r_0(g)=\tr(\xi\circ r_u)(g)-\tr r_0(g)$.
\end{proof}

For a homomorphism $r$ with values in $M_4$ write $i\to_rj$ if $r(g)_{ij}\ne0$ for some $g$, and let $\lambda_i$ be the
first coordinate of $w_i$, so that $(\lambda_0,\dots,\lambda_3)=(1,0,1,0)$. For $t=(t_1,t_2)\in(\Qlbar^\times)^2$ and
$X\in M_4(\Qlbar)$ let $t\cdot X$ be the matrix whose entry $(i,j)$ is $X_{ij}$ multiplied by $t_1^{a}t_2^{b}$, where
$(a,b)=w_i-w_j$; this is the action of the torus used in Step 1 of \cref{pf:modularity:ineq} below.

\begin{lemma}\label{pf:modularity:twoirr}
Let $D$ be a good raw datum with a model, and assume the local theorems, $\sfname{SoV}(\Omega,\ell)$, $\sfname{Fin}_0(D)$
and $\sfname{NoEnd}(D)$. Let $\xi:S^\circ\to\Qlbar$ be an $\OO$-algebra map whose restriction to $A$ is two-signed. Then
the $\Qlbar$-span of the image of $\xi\circ r_u$ is $M_4(\Qlbar)$.
\formallabel{UpperTheory.V.GradeZeroInterpolatingElement.twoIrr}
\end{lemma}

\begin{proof}
By \lref{UpperTheory.V.GradeZeroInterpolatingElement.ptScale} (which uses $\sfname{Fin}_0(D)$: there are
$t\in(\Qlbar^\times)^2$ and a characteristic-zero point $r$ of $D$ with $r(g)=t\cdot(\xi\circ r_u)(g)$ for all $g$). As
$t\cdot$ multiplies each entry by a non-zero scalar, $i\to_rj$ if and only if $i\to_{\xi\circ r_u}j$; and $X\mapsto t\cdot X$
is a linear automorphism of $M_4(\Qlbar)$ carrying the span of the image of $\xi\circ r_u$ onto that of $r$, so if $r$ is
irreducible, so is $\xi\circ r_u$. By \lref{UpperTheory.V.GradeZeroInterpolatingElement.signs} (which uses
$\sfname{Fin}_0(D)$), if $i\to_{\xi\circ r_u}j$ fails for all $(i,j)$ such that the second coordinate of $w_i-w_j$ is
negative (resp.\ positive), then $\xi|_A$ is negatively (resp.\ positively) trivial; as $\xi|_A$ is two-signed, neither
happens. By \lref{UpperTheory.V.GradeZeroInterpolatingElement.tri} (which uses the local theorems), one of the
following holds.
\begin{enumerate}[label=(\alph*)]
\item $r$ is irreducible. Then so is $\xi\circ r_u$.
\item $r$ is of \emph{diagonal type}: for $i\ne j$, $i\to_rj$ excludes $j\to_ri$. By
\lref{UpperTheory.V.GradeZeroInterpolatingElement.diagCases} (which uses $\sfname{SoV}$), $i\to_rj$ then fails either
for all $(i,j)$ with second coordinate of $w_i-w_j$ negative, or for all with it positive. Transferring to $\xi\circ r_u$,
this contradicts the two-signedness of $\xi|_A$.
\item $0\to_r2$, $2\to_r0$, and $i\to_rj$ fails either for all $(i,j)$ with $\lambda_i<\lambda_j$ or for all with
$\lambda_i>\lambda_j$. The same holds for $\xi\circ r_u$, so by
\lref{UpperTheory.V.GradeZeroInterpolatingElement.ptLimit} (which uses $\sfname{Fin}_0(D)$) there is
$\xi_0:S^\circ\to\Qlbar$ with $(\xi_0\circ r_u)(g)_{ij}=(\xi\circ r_u)(g)_{ij}$ if $\lambda_i=\lambda_j$ and $=0$ otherwise.
By \lref{UpperTheory.V.GradeZeroInterpolatingElement.ptScale} again, $r_0'=t_0\cdot(\xi_0\circ r_u)$ is a
characteristic-zero point of $D$ for some $t_0$. It is block diagonal for $\{0,2\}\cup\{1,3\}$, and $i\to_{r_0'}j$ whenever
$\lambda_i=\lambda_j$ and $i\to_rj$; in particular $0\to_{r_0'}2$ and $2\to_{r_0'}0$. The values of a characteristic-zero
point are in $\mathfrak{gsp}$-form and the span of its image contains the diagonal matrix units
(\lref{UpperTheory.V.GradeZeroInterpolatingElement.ptIn}); for such a representation, $0\to2$, $2\to0$ and block
diagonality imply that the blocks $(r_0'(g)_{ij})_{i,j\in\{0,2\}}$ span $M_2(\Qlbar)$
(\lref{UpperTheory.V.GradeZeroInterpolatingElement.geAlg}). This contradicts $\sfname{NoEnd}(D)$.\qedhere
\end{enumerate}
\end{proof}

\begin{lemma}\label{pf:modularity:m1b}
Let $D$ be a good raw datum with a model, and assume $\sfname{SoV}(\Omega,\ell)$. Then the $F_i$, the empty family, and
the $Z_i$ span the cotangent space of $(A,x_0)$ in degrees $-1$, $-2$ and $0$ respectively.
\formallabel{UpperTheory.V.GradeZeroInterpolatingElement.m1b}
\end{lemma}

\begin{proof}
By \lref{UpperTheory.V.GradeZeroInterpolatingElement.aSmooth}, $x_0$ is a smooth point of $A$ with tangent
dimensions $n_{r,0,r',0,m}$. For a chart as in \cref{pf:vwin-he:chart}, the rank of the cotangent lattice in degree $j$ is
at most the number of variables of weight $j$ (\lref{cotangent_rank_le_weight_count}); so it is at most $r$, $0$ and
$m$ in degrees $-1$, $-2$ and $0$. In degree $-2$ the cotangent lattice has rank $0$, and every family spans it
(\lref{spans_cotangent_of_rank_zero}). In degrees $j\in\{-1,0\}$ we compare $A$ with the $\Z^2$-graded ring $R$ of the
model: $A_j=R_{(0,j)}$, and the inclusion $A\to R$ is surjective onto $R_{(0,j)}$ and compatible with the origins. By
condition (b) of \cref{pf:modularity:model}, the cotangent lattice of $R$ at the origin has rank exactly $r$ in weight
$(0,-1)=-w_\Psi$ and $m$ in weight $0$ (\lref{cotangent_rank_at_weight_eq_dims}). As the ranks agree, it suffices that
the images of the $F_i$, resp.\ the $Z_i$, span the cotangent space of $R$ in weight $-w_\Psi$, resp.\ $0$
(\lref{spans_cotangent_transfer_equal_rank}). By the derivation test \lref{spans_cotangent_of_derivation_test}
(which uses that $D$ is good and $\sfname{SoV}$), a family of elements of weight $w\in\{-w_\Psi,0\}$ vanishing at the origin
spans the cotangent space of $R$ in weight $w$ as soon as every $\OO$-linear point derivation of $R$ at the origin which is
homogeneous of weight $w$ and vanishes on the family also vanishes on all the generators $F_i,E_i,Z_i$ of $\ker\pi$. A
homogeneous derivation of weight $-w_\Psi$ vanishes on the $E_i$ and the $Z_i$, whose weights $w_\Psi$ and $0$ differ from
$-w_\Psi$; a homogeneous derivation of weight $0$ vanishes on the $F_i$ and the $E_i$. Hence both families pass the test.
\end{proof}

\begin{proposition}[Key inequality]\label{pf:modularity:ineq}
Let $D$ be a good raw datum with a model. Assume Lemma M, Lemma M$'$, the local theorems, $\sfname{SoV}(\Omega,\ell)$,
$\sfname{SoFin}(\Omega,\ell)$, $\sfname{Fin}_0(D)$, $\sfname{NoEnd}(D)$, and $\sfname{Rigid}(D,r)$ for every irreducible
characteristic-zero point $r$. Let $g\in G_\Q$ be such that $a_g(r)\ne0$ for every irreducible characteristic-zero point
$r$. Then $\sfname{Ineq}(D,g)$ holds.
\formallabel{UpperTheory.V.GradeZeroInterpolatingElement.galoisRep_ineq_of_allIrreducible_nonvanishing}
\end{proposition}

\begin{proof}
The grading of $A$ is a direct sum decomposition (\lref{UpperTheory.V.GradeZeroInterpolatingElement.aGraded}).
Let $a\in A_0$ be as in \cref{pf:modularity:aelt}. In this proof a \emph{point} is an $\OO$-algebra map $A\to\Qlbar$ which
vanishes on $I$, and two-signed points and orbits are as in \cref{pf:vwin-he:wring}.

\emph{Step 1: points give Galois representations.} For every two-signed point $x$ there is an irreducible
characteristic-zero point $\rho_x$ of $D$ with $a_g(\rho_x)=x(a)$; and if $\rho_x$ and $\rho_{x'}$ are conjugate, then $x$
and $x'$ lie in the same orbit. Indeed: $x$ extends to a map
$\xi:S^\circ\to\Qlbar$ (\lref{UpperTheory.V.GradeZeroInterpolatingElement.ptLift}); by $\sfname{Fin}_0(D)$ there is
$t\in(\Qlbar^\times)^2$ such that $\rho_x=t\cdot(\xi\circ r_u)$ is a characteristic-zero point of $D$
(\lref{UpperTheory.V.GradeZeroInterpolatingElement.ptScale}); since $x$ is two-signed, $\xi\circ r_u$, hence
$\rho_x$, is irreducible, by the local theorems, $\sfname{SoV}$, $\sfname{Fin}_0$ and $\sfname{NoEnd}$
(\cref{pf:modularity:twoirr}); traces are invariant under conjugation, so
$a_g(\rho_x)=a_g(\xi\circ r_u)=x(a)$; and conjugate representations $\rho_x,\rho_{x'}$ come from points in the same
orbit (\lref{same_orbit_of_conjugate_reps}).

\emph{Step 2: Lemma M.} The hypotheses (M1)--(M5) of Lemma M hold for $(A,x_0)$, the system $(F,E,Z)$ and $a$:
(M1) by \lref{UpperTheory.V.GradeZeroInterpolatingElement.aSmooth} and, for the spanning conditions, by
\cref{pf:modularity:m1b} (which uses $\sfname{SoV}$); (M2) by
\lref{UpperTheory.V.GradeZeroInterpolatingElement.m2} (which uses $\sfname{Fin}_0$); (M3) by
\lref{UpperTheory.V.GradeZeroInterpolatingElement.m3} (local theorems, $\sfname{SoV}$, $\sfname{Fin}_0$); (M4) by
\lref{UpperTheory.V.GradeZeroInterpolatingElement.m4} (local theorems, $\sfname{SoV}$, $\sfname{Fin}_0$,
$\sfname{NoEnd}$ and the rigidity of all irreducible points); and (M5) holds because $a\in A_0$, $x_0(a)=0$, and for a
two-signed point $x$, $x(a)=a_g(\rho_x)\ne0$ by Step 1 and the hypothesis on $g$. We apply Lemma M with $K=\Qlbar$ and
the normalised valuation $v=e_\OO v_\ell$ (which restricts to $\ord_\OO$ on $\OO$). It gives: a finite set covering $I$
exists, $\modularitycot_{-1}(x_0;F)<\infty$, and for every finite set $T$ covering $I$,
\begin{equation}\label{pf:modularity:LM}
\modularitycot_{-1}(x_0;F)\le2e_\OO\sum_{x\in T}v_\ell(x(a)).
\end{equation}

\emph{Step 3: a minimal cover.} Choose a covering set $T$ of minimal cardinality. Its elements lie in pairwise distinct
orbits: if $x\ne x'$ in $T$ were in the same orbit, $T\smallsetminus\{x'\}$ would still cover $I$, since orbits are
equivalence classes. Let $X=\{\rho_x:x\in T\}$. By Step 1 the elements of $X$ are irreducible characteristic-zero
points; they are pairwise non-conjugate, since conjugate $\rho_x,\rho_{x'}$ come from $x,x'$ in the same orbit, hence
$x=x'$; and for the same reason $x\mapsto\rho_x$ is injective on $T$, so that
$\sum_{r\in X}v_\ell(a_g(r))=\sum_{x\in T}v_\ell(x(a))$.

\emph{Step 4: Lemma M$'$.} By $\sfname{SoFin}(\Omega,\ell)$ and the monotonicity of $\Sha_n$ in $n$,
$\modularitySha(D)<\infty$; moreover there are $0\ne a'\in\OO$ and an injective homomorphism from the $\Sha$-group
defining $\modularitySha(D)$ to the group of $\Sha$-cocycles modulo $a'$
(\lref{UpperTheory.V.GradeZeroInterpolatingElement.shaCoc}). By \cref{pf:modularity:model}(e) these cocycles give
homogeneous point derivations of degree $-1$ at $x_0$ modulo $a'$ which vanish on $\ker\pi$, in particular on the
$F_i$, and Lemma M$'$ gives $\modularitySha(D)\le\modularitycot_{-1}(x_0;F)$
(\lref{UpperTheory.V.GradeZeroInterpolatingElement.mPrime}).

\emph{Conclusion.} By Steps 2--4,
$\|\modularitySha(D)\|_\OO=\modularitySha(D)/e_\OO\le\modularitycot_{-1}(x_0;F)/e_\OO\le2\sum_{r\in X}v_\ell(a_g(r))$,
which is \eqref{pf:modularity:ineqdef}.
\end{proof}

The existence of models is the deformation-theoretic part of the argument. It rests on a computation of residual Galois
cohomology, which we formulate first. Let $D$ be a good raw datum and $k=k_\OO$. The \emph{frame} of $D$ consists of the
places of $T$, with distinguished place $\ell$ and decomposition groups $D_{B_v}$, and the places $q\in S\smallsetminus T$
with the subgroups $I_{B_q}\cap\ker\bar r_0$ of their inertia groups. For a finite $k[G_{\Q,S}]$-module $V$ let
$H^1_c(V)$ be the space of pairs $(z,(y_v)_{v\in T})$ of a locally constant $1$-cocycle $z$ of $G_{\Q,S}$ vanishing on the
subgroups $I_{B_q}\cap\ker\bar r_0$ and elements $y_v\in V$ with $z(g)=gy_v-y_v$ on $D_{B_v}$ and $y_\ell=0$, and let
$\mathrm{Ob}(V)$ be the corresponding space of framed $2$-cocycles modulo framed coboundaries.

\begin{definition}\label{pf:modularity:numsdef}
For $j\in\Z^2$ let $\mathrm{ad}_j$ be the subspace of weight $j$ of $\mathfrak{gsp}_4(k)$ (for the torus of
\cref{pf:modularity:raw}), with the adjoint action of $G_{\Q,S}$ through $\bar r_0$. $\sfname{Nums}(D)$ is the statement: for
every $j$, $\mathrm{Ob}(\mathrm{ad}_j)$ is finite-dimensional with
$\dim\mathrm{Ob}(\mathrm{ad}_j)=\dim H^1_c(\mathrm{ad}_j)+\varepsilon_j$, where $\varepsilon_0=2$,
$\varepsilon_{w_\Psi}=\varepsilon_{-w_\Psi}=1$ and $\varepsilon_j=0$ otherwise; and $H^1_c(\mathrm{ad}_j)=0$ for
$j\notin\{0,\pm w_\Psi\}$.
\formal{UpperTheory.V.GradeZeroInterpolatingElement.Raw.Nums}
\formal{UpperTheory.V.GradeZeroInterpolatingElement.Raw.frame}
\end{definition}

\begin{proposition}\label{pf:modularity:nums}
The cohomological hypotheses imply $\sfname{Nums}(D)$ for every good raw datum $D$.
\formallabel{UpperTheory.V.GradeZeroInterpolatingElement.nums}
\end{proposition}

\begin{proof}
Let $j=(j_1,j_2)$ and $\gamma=\bar\alpha^{j_1}\bar\Psi^{j_2}:G_\Q\to k^\times$. We collect the following facts.
\begin{enumerate}[label=(\alph*)]
\item The kernel of $\bar r_0$ is an open normal subgroup of index prime to $\ell$, on which $\gamma$ is trivial
(\lref{UpperTheory.V.X.spKer}); in particular $\ell$ does not divide the index of $I_{B_q}\cap\ker\bar r_0$ in
$I_{B_q}$ for $q\in S\smallsetminus T$.
\item $G_\Q$ acts on $\mathrm{ad}_j$ through the scalar character $\gamma$, since $\bar r_0$ is diagonal
(\lref{UpperTheory.V.X.spAct}), and $\dim\mathrm{ad}_j$ is the multiplicity of $j$ as a weight of $\mathfrak{gsp}_4$
(\lref{UpperTheory.V.X.spDim}).
\item $\gamma(c)=(-1)^{j_1}$, because $\bar\Psi(c)=1$ and $\bar\alpha(c)=-1$ for the complex conjugation $c$ of $D$; the
reduction of $\modularitycyc$ is $-1$ at $c$ and $1$ at $\sigma_\ell$ (as $\modularitycyc(\sigma_\ell)=1$); and $\sigma_v$
lies in $D_{B_v}$ for $v\in T$.
\item If $j$ is one of the six roots $\pm w_\alpha,\pm w_\beta,\pm w_\theta$ (the weights of the entries $(i,j)$ with
$\lambda_i\ne\lambda_j$), then every locally constant $1$-cocycle with values in $k(\gamma)$ which vanishes on the inertia
elements away from $T$ acting trivially through $\bar r_0$ and is a coboundary on $D_{B_v}$ for $v\in T$ is a coboundary,
and $\gamma(\sigma_v)\ne1$ for $v\in T$ (\lref{nums_odd}: for these $j$, $\gamma$ is $\bar\alpha^{\pm1}$,
$\bar\beta^{\pm1}$ or $\bar\theta^{\pm1}$, so both statements are among the conditions of goodness).
\end{enumerate}
Given these facts, \lref{nums_all} gives the assertion for $j$: under the cohomological hypotheses, for every finite
$k[G_{\Q,S}]$-module $M$ on which $G_\Q$ acts through a character $\gamma$ trivial on an open normal subgroup $K$ as in (a),
with $\dim M$ the multiplicity of $j$ in $\mathfrak{gsp}_4$, $\gamma(c)=(-1)^{j_1}$, and the conditions (c) and (d),
$\mathrm{Ob}(M)$ is finite-dimensional of dimension $\dim H^1_c(M)+\varepsilon_j$, and $H^1_c(M)=0$ if
$j\notin\{0,\pm w_\Psi\}$.
\end{proof}

\begin{proposition}[Existence of models]\label{pf:modularity:exmodel}
Let $D$ be a good raw datum. Assume Lemma P, that every open subgroup of $G_{\Q,S}$ has only finitely many open
homomorphisms to $\Z/\ell$, $\sfname{Nums}(D)$, and that a local base for $D$ exists: a complete local $\OO$-algebra
$\mathcal B$ with a $\Z^2$-weighted presentation (with weights among the weights of entries) and an origin, together with
lifts $\rho_v$ ($v\in T$) of the local residual representations $\bar r_0|_{D_{B_v}}$ satisfying the local conditions of
\cref{pf:modularity:raw}, universal for such tuples, whose entries are homogeneous of the weights of the corresponding
entries and specialise to those of $r_0$ at the origin, and such that the origin is smooth with the expected tangent
dimensions. Then $D$ has a model.
\formallabel{UpperTheory.V.GradeZeroInterpolatingElement.exModel}
\formal{UpperTheory.V.Alg.PhiCond}
\formal{UpperTheory.V.GradeZeroInterpolatingElement.LocBase}
\end{proposition}

\begin{proof}
Let $\mathcal B$ be a local base. The model is built in three steps.
\begin{enumerate}[label=(\arabic*)]
\item (\lref{UpperTheory.nonempty_local_frames_of_good_data}, which uses goodness, Lemma P, the finiteness of
characters and $\sfname{Nums}(D)$.) There is a \emph{relative presentation} of the universal $D$-point over $\mathcal B$: a
complete local $\OO$-algebra $U$ with a $\Z^2$-weighted presentation extending that of $\mathcal B$ by finitely many
homogeneous variables, a finite family of homogeneous relations, and a $D$-point over $U$ which is universal for
$D$-points whose local restrictions come from $\mathcal B$; the numbers of relative variables and relations of each weight are
controlled by the dimensions in $\sfname{Nums}(D)$.
\item (\lref{UpperTheory.smoothPt_part_of_smoothPt_associatedGraded} and \lref{UpperTheory.nonempty_cover_of_relative_presentation_smoothPt}.)
As the origin of $\mathcal B$ is smooth, the relative presentation has a \emph{cover}: a complete local $\OO$-algebra $P$
with a weighted presentation, smooth at its origin with tangent dimensions those of $\mathcal B$ increased by the numbers
of relative variables of each weight, and a surjection $p:P\to U$ compatible with $\mathcal B$ whose kernel is generated by
homogeneous elements of the weights of the relations, such that a point derivation modulo $a$ of $U$ which is induced on
$\mathcal B$ by a homogeneous point derivation of some weight $j$, and vanishes on the relative variables of weight other
than $j$, lifts to a homogeneous point derivation of $P$ of weight $j$.
\item One divides $P$ and $U$ by the images of the six local coordinates of odd weight $\pm w_\alpha,\pm w_\beta,\pm
w_\theta$ at each place of $T$ (which are independent in the cotangent space,
\lref{UpperTheory.V.GradeZeroInterpolatingElement.ExM.tOdd_ind}) and by the local equations, obtaining $\pi:R\to S$.
The generators of $\ker\pi$ are then homogeneous of weights $-w_\Psi$, $w_\Psi$ and $0$ only, and counting them with the
tangent dimensions (\lref{UpperTheory.V.GradeZeroInterpolatingElement.ExM.dims}) gives $r+1$ generators $F_i$ of
weight $-w_\Psi$, $r'$ generators $E_i$ of weight $w_\Psi$ and $m$ generators $Z_i$ of weight $0$, where
$r=1+|S_b|+a_-$, $r'=2+|S_b|+a_+$ and $m=2+2|S_b|+a_0$, the $a_\bullet$ being the numbers of relative variables of
weights $-w_\Psi$, $w_\Psi$, $0$; and $R$ is smooth at its origin with the tangent dimensions required in (b) of
\cref{pf:modularity:model}. The $D$-point $r_u$ is the image of the universal $D$-point of $U$: its entries generate $S$
and are images of homogeneous elements of $R$ of the right weights, with values at the origin those of $r_0$ (the local
data become trivial in $S$, \lref{UpperTheory.V.GradeZeroInterpolatingElement.ExM.al_one}).
\end{enumerate}
It remains to check (e) of \cref{pf:modularity:model}. Let $z$ be a $\Sha$-cocycle modulo $a$. By
\lref{UpperTheory.exists_quotientDerivation_of_shaCocycle} there is a point derivation $\delta_U$ modulo $a$ of $U$ with
$\delta_U(r(g)_{ij})=(Xr_0(g))_{ij}\,z(g)$ for the universal $D$-point $r$ and the matrix $X$ whose entries $(2,0)$ and
$(3,1)$ are $1$ and the others $0$, induced on $\mathcal B$ by a homogeneous point derivation of weight $-w_\Psi$ and
vanishing on the local data. The derivation $\delta_U$ vanishes on the relative variables of weight other than $-w_\Psi$
(the entries of $Xr_0(g)$ vanish outside the weight $-w_\Psi$), on the odd local coordinates and on the local equations
(which involve $g_0$, and $z(g_0)=z(g_0^{-1})=0$), so (2) and the construction in (3) give a point derivation $\delta$ of
$R$ at the origin, homogeneous of weight $-w_\Psi$, with $\delta(f)\equiv\delta_U(\pi(f))$ modulo $a$; in particular
$\delta$ maps $\ker\pi$ into $a\OO$. Since $(Xr_0(g))_{20}=\theta(g)$ and $(Xr_0(g))_{31}=\Psi(g)$, $\delta$ has the
properties required in (e).
\end{proof}

\begin{corollary}\label{pf:modularity:exmodel2}
Assume the cohomological hypotheses and Lemma P. Every good raw datum has a model.
\end{corollary}

\begin{proof}
The finiteness of characters follows from the cohomological hypotheses
(\lref{UpperTheory.V.GradeZeroInterpolatingElement.phiGam}), $\sfname{Nums}(D)$ holds by \cref{pf:modularity:nums},
and a local base exists by the local theorems, which are among the cohomological hypotheses
(\lref{UpperTheory.V.GradeZeroInterpolatingElement.locBase}). Apply \cref{pf:modularity:exmodel}.
\end{proof}

\subsubsection{The Frobenius element and the Eisenstein congruence number}

\begin{proposition}\label{pf:modularity:frob}
Assume $\sfname{NumberSp4}$. There are $C,A\in\N$ with the following property. Let $D$ be a good raw datum over $\OO$,
$N\ge1$ such that $\psi$, $\xi$ and $\omega$ are trivial on $G_{\Q(\mu_N)}$, $k\ge1$ with $\psi^k=\xi^k=\omega^k=1$,
$p\ge5$ a prime with $p\nmid N$ and $p\notin S$, and $\iota$ such that every irreducible characteristic-zero point of $D$
is automorphic over imaginary quadratic fields via $\iota$. Then there is $\sigma\in G_\Q$ such that $a_\sigma(r)\ne0$ for
every irreducible characteristic-zero point $r$, and such that $\sfname{Ineq}(D,\sigma)$ implies
\begin{equation}\label{pf:modularity:frobbound}
\|\modularitySha(D)\|_\OO\log\ell\ \le\ 2Ck\,(N(e+1))^A\log p.
\end{equation}
\formallabel{UpperTheory.V.Z.Win.A4.exists_frob_nonvanishing_with_sha_bound}
\end{proposition}

\begin{proof}
Let $(C,A)$ be the constants $(c,A)$ of $\sfname{NumberSp4}$. Since $\ell\in S$ and $p\notin S$, $p\ne\ell$. Let
$\sigma$ be a Frobenius element at $B_p$ (which lies over $p$ since $D$ is good). Then $\modularitycyc(\sigma)=p$, as
$\sigma$ acts on $\ell$-power roots of unity by $\zeta\mapsto\zeta^p$, and with the $k$-th roots of unity
$z_1=\iota((\omega^{-1}\xi\psi)(\sigma))$, $z_2=\iota(\psi(\sigma))$, $z_3=\iota((\omega^{-1}\xi)(\sigma))$ we have
\[
 \tr r_0(\sigma)=\theta(\sigma)+\Psi(\sigma)+\alpha(\sigma)+1=\iota^{-1}\mathrm{eis}_4(p,e;z_1,z_2,z_3),
\]
since $\theta=\modularitycyc^{e+1}\omega^{-1}\xi\psi$, $\Psi=\modularitycyc^e\psi$ and $\alpha=\modularitycyc\,\omega^{-1}\xi$.
Let $\mathcal N$ be the integer given by $\sfname{NumberSp4}$ for $(e,N,p,k,z_1,z_2,z_3)$; note $e\ge2$ since $D$ is
good.

Every irreducible characteristic-zero point $r$ of $D$ is admissible of level $N$: it is irreducible, automorphic over
imaginary quadratic fields by hypothesis, of level $N$ and type $(N,(e+1,e,1,0))$, it preserves a bilinear form up to a
character of cyclotomic weight $e+1$ (\lref{UpperTheory.V.GradeZeroInterpolatingElement.level}, which uses that
$\psi,\xi,\omega$ are trivial on $G_{\Q(\mu_N)}$), and it satisfies the side condition
(\lref{UpperTheory.V.Z.Win.A4.safe4'}). Applying $\sfname{NumberSp4}$ to $X=\{r\}$ gives
$\tr r(\sigma)\ne\tr r_0(\sigma)$, i.e.\ $a_\sigma(r)\ne0$.

Assume $\sfname{Ineq}(D,\sigma)$, with the set $X$. Its elements are admissible and pairwise non-isomorphic, so
$\|\mathcal N\|\le\prod_{r\in X}\|a_\sigma(r)\|$, and all factors are non-zero. Taking $-\log$ and using
$-\log\|\mathcal N\|=v_\ell(\mathcal N)\log\ell\le\log|\mathcal N|$ for a non-zero integer $\mathcal N$,
\[
 \sum_{r\in X}-\log\|a_\sigma(r)\|\le\log|\mathcal N|\le C\varphi(k)(N(e+1))^A\log p\le Ck(N(e+1))^A\log p .
\]
Together with \eqref{pf:modularity:ineqdef} this gives \eqref{pf:modularity:frobbound}.
\end{proof}

\begin{proposition}\label{pf:modularity:numberSp4}
$\sfname{Hypotheses'}$ (the part of $\sfname{HypothesesH'}$ formed by \ref{pf:vwin-he:H1} and \ref{pf:vwin-he:H2})
implies $\sfname{NumberSp4}$.
\formallabel{UpperTheory.GalN.V.numberSp4_n_alt}
\formallabel{UpperTheory.GalN.V.Win.N.number_sp4_of_admissible_to_admissible_n}
\end{proposition}

\begin{proof}
We use two consequences of $\sfname{Hypotheses'}$. First, an admissible representation in the sense of
\cref{pf:modularity:number} is also admissible in the stronger sense $\sfname{AdmRN}$ of the formal development, in which
automorphy over imaginary quadratic fields is replaced by the automorphy statement over $\Q$ deduced from it and from the
polarisation (\lref{UpperTheory.GalN.V.autQN'_alt}). Second, the counting statement
\lref{UpperTheory.GalN.V.Win.N.countQ_alt}, proved from $\sfname{Hypotheses'}$ and two level statements
(\lref{UpperTheory.V.Keep.levelPA_alt} and \lref{UpperTheory.V.Z.Win.A4.admissible_restriction_level_n_of_printed_results}),
asserts: \emph{for $n\ge2$ there are $c,A$ such that for every strictly decreasing $m\in\Z^n$ with
$m_i+m_{n+1-i}=\mathrm{wt}$ and $|m_i|\le H$, every $M\ge1$ and every prime $p\nmid M$, there are a finite set $\mathcal S$
with $\#\mathcal S\le c(M(H+1))^A$ and multisets $\beta_s$ ($s\in\mathcal S$) of $n$ complex numbers such that: the family
$(\beta_s)$ is stable under $\Aut(\C)$; for $a\ge0$ with $-a\le m_i$ for all $i$, $p^ab$ is an algebraic integer for
$b\in\beta_s$; the $\beta_s$ satisfy a window condition at $p$; and for every prime $\ell\ne2,p$ and every $\iota$ there
is a map $\mathrm{cl}$ from the set of $\sfname{AdmRN}$-representations of level $M$ and weights $m$ to $\mathcal S$,
with $\mathrm{cl}(r)=\mathrm{cl}(r')$ only if $r\cong r'$, such that the characteristic polynomial of $r(\sigma^{-1})$ for
a Frobenius element $\sigma$ above $p$ has the roots $\iota^{-1}(b^{-1})$, $b\in\beta_{\mathrm{cl}(r)}$.}

Let $c,A$ be the constants for $n=4$, and put $c'=2^{A+1}c$, $A'=A+1$. Let $e\ge2$, $M$, $p\ge5$, $k$ and $z_1,z_2,z_3$
be as in $\sfname{NumberSp4}$; apply the counting statement with $m=(e+1,e,1,0)$, $\mathrm{wt}=e+1$, $H=e+1$. For
$(s,u)\in\mathcal S\times(\Z/k)^\times$ put
\[
 f(s,u)=\sum_{b\in\beta_s}b-\mathrm{eis}_4\bigl(p,e;z_1^u,z_2^u,z_3^u\bigr).
\]
Each $f(s,u)$ is an algebraic integer (the $b$ are, by the integrality condition with $a=0$). The family is stable
under $\Aut(\C)$: an automorphism $\tau$ permutes the $\beta_s$ and acts on $k$-th roots of unity by $z\mapsto z^v$ for some
$v\in(\Z/k)^\times$, so $\tau f(s,u)=f(s',vu)$. Hence $\mathcal N=\prod_{s,u}f(s,u)$ is fixed by $\Aut(\C)$, so it is a rational number which is an algebraic integer,
i.e.\ a rational integer. The window condition gives
$|\sum_{b\in\beta_s}b|\le3\sqrt p\,p^{(e+1)/2}+p^{(e+1)/2}$ (\lref{UpperTheory.GalN.V.WinC.sum_le}), and an
elementary estimate for $p\ge5$, $e\ge2$ and
$|z_i|=1$ shows that then $f(s,u)\ne0$ and $|f(s,u)|\le p^{2(e+1)}$: indeed $p^{e+1}$ exceeds the sum of the absolute
values of the other terms, which is at most $\frac72p^e+p^e+p+1$. Hence $\mathcal N\ne0$ and
\[
 \log|\mathcal N|\le\#\mathcal S\cdot\varphi(k)\cdot2(e+1)\log p\le c'\varphi(k)(M(e+1))^{A'}\log p,
\]
using $\#\mathcal S\le c(M(e+2))^A\le c2^A(M(e+1))^A$ and $2(e+1)\le2M(e+1)$.

Let $\ell\ne2,p$, $\iota$, $\sigma$ a Frobenius element above $p$ and $X$ a finite set of pairwise non-isomorphic
admissible representations. For $r\in X$, the eigenvalues of $r(\sigma)$ are the $\iota^{-1}(b)$, $b\in\beta_{\mathrm{cl}(r)}$,
so $\tr r(\sigma)-\iota^{-1}\mathrm{eis}_4(p,e;z)=\iota^{-1}f(\mathrm{cl}(r),1)\ne0$. The map $r\mapsto(\mathrm{cl}(r),1)$
is injective on $X$. Since $\iota^{-1}\mathcal N=\prod_{s,u}\iota^{-1}f(s,u)$ and every factor is an algebraic integer, hence
of absolute value $\le1$ in $\Qlbar$,
$\|\mathcal N\|\le\prod_{r\in X}\|\iota^{-1}f(\mathrm{cl}(r),1)\|$. This
proves $\sfname{NumberSp4}$ with the constants $(c',A')$.
\end{proof}

\subsubsection{Rigidity}

\begin{lemma}\label{pf:modularity:adjin}
$\sfname{HypothesesH'}$ implies the cohomological hypotheses and the inputs of the unramified adjoint lifting
argument (\lref{UpperTheory.V.Ws.UnramifiedInputs}: the cohomological hypotheses, $\sfname{HypothesesH'}$ itself,
the modularity statement $\sfname{ModU}$, and the existence statements $\sfname{LevEx}$, $\sfname{LamEx}$ and
$\sfname{ProModularRingExists}$); and it implies the adjoint lifting theorem.
\formallabel{UpperTheory.V.HV.DescentVanishingReductions.adjoint_lift_of_prints}
\formallabel{UpperTheory.V.HV.DescentVanishingReductions.adjoint_lifting_of_cohomology}
\formallabel{UpperTheory.V.Ws.UnramifiedInputs}
\formallabel{UpperTheory.V.Th.AdjointLiftingTheorem}
\end{lemma}

\begin{proof}
The cohomological hypotheses are components of \ref{pf:vwin-he:H1} (\cref{pf:modularity:loc}). $\sfname{ModU}$ holds by
\cref{pf:modularity:modU}; $\sfname{LevEx}$ is a theorem
(\lref{UpperTheory.V.W.levEx}); $\sfname{LamEx}$ follows from local class field theory, a component of \ref{pf:vwin-he:H1}
(\lref{UpperTheory.V.W.lamEx} and \cref{pf:modularity:loc}(e)); and
$\sfname{ProModularRingExists}$ follows from the Galois-cohomological part of \ref{pf:vwin-he:H1}
(\lref{UpperTheory.V.W.b0Ex} and \lref{UpperTheory.V.W.phiAll} and \cref{pf:modularity:loc}(c)). The adjoint lifting
theorem is \lref{UpperTheory.V.Th.adjoint_lifting_of_cohomology}, applied to $\sfname{HypothesesH'}$ and the
same cohomological consequence of \ref{pf:vwin-he:H1}; it asserts that for a CM field $L$ not containing the $\ell$-th
roots of unity, a large coefficient ring, $2\le n\le4$ and a universal ordinary (Hida) deformation ring of an
$n$-dimensional residual representation, the adjoint Bloch--Kato Selmer group vanishes at every point which is
automorphic and locally generic.
\end{proof}

\begin{theorem}[Rigidity]\label{pf:modularity:rig}
The inputs of \cref{pf:modularity:adjin} and the adjoint lifting theorem imply $\sfname{RigSp4}$.
\formallabel{UpperTheory.V.Th.rigSp4_alt}
\end{theorem}

\begin{proof}
Let $D$, $N$, $\iota$ and $r$ be as in \cref{pf:modularity:rigdef}, and let $r'$ be a point of $D$ with values in
$\Qlbar[\epsilon]$ reducing to $r$.
\begin{enumerate}[label=(\roman*)]
\item (\lref{UpperTheory.V.Z.Win.Dn.exists_cm_field_twist_h1}, which uses the inputs of \cref{pf:modularity:adjin}
and the adjoint lifting theorem) There are an imaginary quadratic field $K\subset\Omega$ and a character
$\theta:G_\Q\to\Qlbar^\times$ such that $r(g)$ is in $\mathfrak{gsp}$-form with multiplier $\theta(g)$ for all $g$, and such
that every representation of $G_\Q$ on $M_4(\Qlbar)$ which is an adjoint representation of
$(r|_{G_K},\theta|_{G_K})$ has vanishing Bloch--Kato Selmer group.
\item (\lref{UpperTheory.V.GradeZeroInterpolatingElement.tanBK}, using the local theorems) The infinitesimal
deformation $r'$ of $r$ defines a continuous $1$-cocycle $z$ of $G_\Q$ with values in $\mathfrak{sp}_4(\Qlbar)$, for the
adjoint action of $r$, which satisfies the Bloch--Kato local conditions.
\item Transported to the adjoint representation of (i), $z$ becomes a Bloch--Kato cocycle, hence a coboundary.
Translating back and projecting to $\mathfrak{sp}_4$ (\lref{rg_bk} and \lref{rgP_mem}) gives $X\in\mathfrak{sp}_4(\Qlbar)$
with $r'(g)=r(g)+\epsilon\,(Xr(g)-r(g)X)$ for all $g$, that is, $r'=(1+\epsilon X)r(1-\epsilon X)$; this is
$\sfname{Rigid}(D,r)$.
\end{enumerate}
\end{proof}

\subsubsection{The main argument}\label{sec:pf:modularity:mainarg}

We first record three elementary facts.

\begin{lemma}\label{pf:modularity:pzero}
For $m\ge1$ let $p_0(m)$ be the least prime $p\ge5$ with $p\nmid m$. Then $p_0(m)\le2m+10$.
\formallabel{UpperTheory.V.HV.pZeroLe}
\formallabel{UpperTheory.V.HV.pZero}
\end{lemma}

\begin{proof}
By Bertrand's postulate there is a prime $p$ with $m+5<p\le2(m+5)$; then $p\ge5$ and $p\nmid m$ since $p>m$.
\end{proof}

\begin{lemma}\label{pf:modularity:shaeq}
Let $h$ be a host big for $t$, $\ell$ odd, $0<e<t$, $\psi\in\mathcal X_h(\OO,e)$ and $D$ a good raw datum of the case of
$(h,e,\psi)$. Then $V_\ell(\psi,e)\le\modularitySha(D)$.
\formallabel{UpperTheory.V.GradeZeroInterpolatingElement.shaEq}
\end{lemma}

\begin{proof}
$D$ has twist $e$ and character $\psi_G$, and $S$ consists of primes. By \lref{vlen_le},
$V_\ell(\psi,e)\le\modularitySha_n(\psi_G,e)$ for every product $n$ of primes and $e\ne0$; take $n=\prod_{p\in S}p$.
\end{proof}

\begin{lemma}\label{pf:modularity:levof}
In the situation of \cref{pf:modularity:shaeq}, let $D$ be of character $(e,\psi,n',p_1,m_0)$ with $c_2(\psi)\le n'$ and
$m_0\mid o_\chi$, and put $N'=\ell2^{n'}N_hp_1$ and $k=(\ell-1)2^{n'}\cdot2o_\chi$. Then $\psi$, $\xi$ and $\omega$ are
trivial on $G_{\Q(\mu_{N'})}$, and $\psi^k=\xi^k=\omega^k=1$.
\formallabel{UpperTheory.V.GradeZeroInterpolatingElement.levOf}
\end{lemma}

\begin{proof}
$\psi_G$ factors through $(\Z/N_h)^\times$, $\xi=\xi_{0,G}\xi'_G$ through $(\Z/p_12^{n'})^\times$, and $\omega$, which
satisfies $\omega^{\ell-1}=1$ and $\omega\equiv\modularitycyc\pmod{\mathfrak m_\OO}$, is the Teichm\"uller lift of
$\modularitycyc\bmod\ell$ and factors through $(\Z/\ell)^\times$; all these moduli divide $N'$. Next
$\omega^{\ell-1}=1$; $\xi_0^{m_0}=1$ with $m_0\mid o_\chi\mid k$, and $\xi'^{\varphi(2^{n'})}=1$ with
$\varphi(2^{n'})\mid2^{n'}\mid k$. Finally $\psi^k=1$: by the structure of the characters $\psi\in\mathcal X_h(\OO,e)$
(\lref{UpperTheory.V.HV.hostComp}), the $2$-component of $\psi$ has order dividing $\varphi(2^{c_2(\psi)})$, which
divides $2^{n'}$, and the $r$-components for odd $r$ have order dividing $2o_\chi$.
\end{proof}

\begin{proof}[Proof of \cref{pf:modularity:mainbig}]
Let $(C_1,A_1)$ be the constants of \cref{pf:modularity:frob}. By an elementary estimate
(\lref{UpperTheory.V.GradeZeroInterpolatingElement.arith}), there is $C\in\N$, depending only on $\chi$, $C_1$ and
$A_1$, such that for every host $h$ with character $\chi$, every odd $\ell$, $\OO$, $e\ge2$, $\psi\in\mathcal X_h(\OO,e)$ and
$(n',p_1)\in\{(n_{\mathrm{II}},1),(n_{\mathrm{III}},p_1(\chi))\}$,
\begin{equation}\label{pf:modularity:arith}
2C_1\,(\ell-1)2^{n'}2o_\chi\,\bigl(\ell2^{n'}N_hp_1(e+1)\bigr)^{A_1}\log p_0(\ell N_hp_1)\le(\ell N_he)^C.
\end{equation}
(The proof bounds $2^{n'}$ by a power of $\ell N_he$, using $o_\chi\le N_\chi\le N_h$, $2^{c_2(\psi)}\le N_h$,
$16\mid N_h$ and \cref{pf:modularity:pzero}.)

Let $h,t,\Omega,\ell,\OO,e,\psi$ be as in \cref{pf:modularity:mainbig}.
\begin{enumerate}[label=(\arabic*)]
\item By \lref{UpperTheory.V.GradeZeroInterpolatingElement.dataEx} (which uses $\sfname{CorX}$ and the roots of
unity in $\OO$) there is a good raw datum $D$ over $\OO$ of the case of $(h,e,\psi)$. (In case II, $\xi$ is a primitive
odd character modulo $2^{n_{\mathrm{II}}}$; in case III, $\xi$ is the product of a primitive odd character modulo
$2^{n_{\mathrm{III}}}$ and a character modulo $p_1(\chi)$ of order $\ord\psi$.)
\item By $(\mathrm{Sp})$, $\sfname{Fin}_0(D)$ and $\sfname{NoEnd}(D)$ hold, and there is $\iota$ such that all irreducible
characteristic-zero points of $D$ are automorphic over imaginary quadratic fields via $\iota$.
\item By \cref{pf:modularity:shaeq}, $V_\ell(\psi,e)\le\modularitySha(D)$.
\item By \cref{pf:modularity:exmodel2}, $D$ has a model.
\item Let $(n',p_1,m_0)$ be $(n_{\mathrm{II}},1,1)$ in case II and $(n_{\mathrm{III}},p_1(\chi),\ord\psi)$ in case III,
so that $D$ is of character $(e,\psi,n',p_1,m_0)$. Then $c_2(\psi)\le n'$ and $m_0\mid o_\chi$: in case II,
$2^{c_2(\psi)}\le2^{c_2(\psi)+6}o_\chi\ell^3N_\chi^2\le2^{n_{\mathrm{II}}}$ and $m_0=1$; in case III, $\ord\psi$ is odd,
so that $c_2(\psi)=0$ and $\ord\psi$ divides the odd part of $o_\chi$ (\lref{UpperTheory.V.HV.hostOdd}).
\item Put $N'=\ell2^{n'}N_hp_1$ and $k=(\ell-1)2^{n'}2o_\chi$; \cref{pf:modularity:levof} applies. Let
$p_0=p_0(\ell N_hp_1)$. By \cref{pf:modularity:pzero}, $p_0\ge5$ is a prime not dividing $\ell N_hp_1$; hence
$p_0\nmid N'$ (as $p_0$ is odd) and $p_0\notin S$ (the primes of $S$ divide $2\ell N_hp_1$).
\item By \cref{pf:modularity:frob} (with $N'$, $k$, $p_0$, $\iota$) there is $\sigma\in G_\Q$ with $a_\sigma(r)\ne0$
for all irreducible characteristic-zero points $r$, such that $\sfname{Ineq}(D,\sigma)$ implies
\eqref{pf:modularity:frobbound}.
\item $D$ lies in context II or III (it is good and of the case of $(h,e,\psi)$ with $h$ big for $t$ and $0<e<t$), so by
$\sfname{RigSp4}$ (with $N'$ and $\iota$) every irreducible characteristic-zero point $r$ satisfies $\sfname{Rigid}(D,r)$.
By \cref{pf:modularity:ineq} (with the model of (4), $\sfname{SoV}(\Omega,\ell)$, $\sfname{SoFin}(\Omega,\ell)$, and
(2), (7)), $\sfname{Ineq}(D,\sigma)$ holds.
\item Hence $\modularitySha(D)<\infty$, so $V_\ell(\psi,e)<\infty$ by (3), and, since normalised lengths are monotone,
using \eqref{pf:modularity:frobbound}, $e=e_D\ge2$ and \eqref{pf:modularity:arith},
\[
 \widetilde V_\ell(\psi,e)\log\ell\le\|\modularitySha(D)\|_\OO\log\ell\le2C_1k(N'(e+1))^{A_1}\log p_0\le(\ell N_he)^C.
\]
\end{enumerate}
This is the assertion of \cref{pf:modularity:mainbig}.
\end{proof}

\begin{remark}\label{pf:modularity:rem-count23}
(1) Lemma K enters the proof of \cref{pf:modularity:count23} only through $\sfname{SoV}$ and $\sfname{SoFin}$.
(2) In \cref{pf:modularity:mainbig} the auxiliary prime at which the Frobenius element is taken is the single prime
$p_0(\ell N_hp_1)$, which depends on $\ell$; no second auxiliary prime is needed because the bound is not summed over
$\ell$.
(3) The rigidity hypothesis of \cref{pf:modularity:ineq} is used only for condition (M4) of Lemma M.
\end{remark}

\subsection{The endpoint statements \texorpdfstring{$(\mathrm{U}_{\mathrm{ord}})$}{(U-ord)} and \texorpdfstring{$(\mathrm{U}_i)$}{(U-i)}}\label{sec:pf:modularity:ends}

In this subsection we prove step \ref{pf:vwin-he:b} of \cref{sec:pf:vwin-he:proof}. The proof is an assembly of a
number of long results of the formal development, which we state; the automorphy of the three-dimensional points is
obtained over an auxiliary abelian CM field $L$, through a ``Galois link'' between the deformation problem over $\Q$
and an abstract residual frame over $L$, and then descended to the imaginary quadratic field $K$.

\begin{definition}[Frames and Galois links]\label{pf:modularity:link}
A \emph{frame} $\Phi$ (for $\ell$, with coefficients in a finite field $k$) is a compact totally disconnected group
$\Gamma_\Phi$ with a sign character $\delta:\Gamma_\Phi\to\{\pm1\}$, an element $c_\Phi$, a finite set $V_\Phi$ of
``places above $\ell$'' with local groups $G_v$, homomorphisms $G_v\to\ker\delta$ and inertia subgroups $I_v\subset G_v$,
a cyclotomic character $\epsilon:\Gamma_\Phi\to\Zl^\times$, a character $\bar\Psi:\Gamma_\Phi\to k^\times$ and an element
$g_{0,\Phi}\in\ker\delta$, together with some numerical data; $\Phi$ may satisfy the conditions $\sfname{Ok}(C_0)$ and
$\sfname{Big}(C_0)$ for an integer $C_0$. Let $\eta:G_\Q\to\{\pm1\}$, $\Psi:G_\Q\to\OO^\times$, $c,g_1\in G_\Q$, $B$ a
valuation ring over $\ell$, $M\ge1$, and $L$ a CM field with maximal totally real subfield $L^+$. A \emph{Galois link}
from $(\eta,\Psi,c,g_1,B,M)$ to $\Phi$ over $L$ consists of: the facts that $L/\Q$ is abelian, $G_L=G_{L^+}\cap\ker\eta$,
$\eta$ is trivial on $I_{B'}\cap G_{L^+}$ for all $B'$ and on $D_{B'}\cap G_{L^+}$ for $B'$ over primes dividing $M$,
and that $\Psi$ is trivial on the inertia elements of $G_{L^+}$ away from $\ell$ on which $\bar\Psi$ is trivial; a
continuous surjection $p:G_{L^+}\to\Gamma_\Phi$ whose kernel is the closed normal subgroup generated by the inertia
elements away from $\ell$ on which $\eta$ and $\bar\Psi$ are trivial, compatible with $(\delta,\eta)$,
$(\epsilon,\modularitycyc)$, $(\bar\Psi,\Psi\bmod\mathfrak m_\OO)$, and sending $c$ to $c_\Phi$ and $g_1$ to $g_{0,\Phi}$;
and an identification of the places $v\in V_\Phi$ with the $G_{L^+}$-orbits of valuation rings over $\ell$ (conjugate
to $B$ by elements of $\ker\eta$), with embeddings $G_v\to G_{L^+}$ onto the decomposition groups, matching inertia
groups and local data. The \emph{link hypotheses} for $(N,e,\eta,c,B,M)$ are: $e\ge1$, $\ell N\mid M\ne0$, $\eta$ is
unramified at the primes not dividing $M$, $B$ lies over $\ell$ and $c$ is a complex conjugation.
\formal{UpperTheory.V.Fr.Frame}
\formal{UpperTheory.V.Pr.GalLink}
\formal{UpperTheory.V.Pr.LinkHyp}
\formal{UpperTheory.V.Pr.SeedClause}
\end{definition}

A \emph{point of rank $n$} of a link $k$ (with coefficients in a discrete valuation ring $A\subset\Qlbar$ finite over
$\OO$) is a lift to $A$ of the residual representation of rank $n$ of the frame whose characters on the inertia groups
$I_v$ are the prescribed powers of $\Psi$; composing with $p$ and restricting to $G_L$ gives a representation
$r:G_L\to\GL_n(\Qlbar)$.

\begin{definition}[Linked-chain statements]\label{pf:modularity:linkedstmts}
\begin{enumerate}[label=(\alph*)]
\item $\sfname{LinkedEnd}$: for every $\Omega$, $\ell$, $\iota$, $\OO$ (with residue characteristic $\ell$), $N$, $e\ge1$
with $N$ even, $\psi$, and all $(\eta,c,g_1,B,M)$ with $\ell N\mid M$, $\eta$ unramified outside $M$, every CM field $L$,
frame $\Phi$, integer $C_0$ and Galois link $k$ from $(\eta,\Psi_{e,\psi},c,g_1,B,M)$ to $\Phi$ over $L$ with $B$ over
$\ell$, $\Phi$ satisfying $\sfname{Ok}(C_0)$ and $\sfname{Big}(C_0)$, and the seed clause for $L$: (i) the characteristic
polynomials of all lifts of rank $3$ of the residual representation of $\Phi$ to complete discrete valuation rings of
characteristic $\ell$, whose local characters are trivial on the inertia elements on which $\bar\Psi_\Phi$ is trivial,
are those of the residual representation (finiteness); (ii) for every irreducible point $r$ of
rank $3$ of $k$, $r\otimes\modularitycyc^{-1}$ is conjugate self-dual automorphic via $\iota$; and (iii) there is an
algebraic character $\kappa$ of $G_L$ such that $r\otimes\kappa$ is conjugate self-dual automorphic via $\iota$ for
every irreducible point $r$ of rank $2$ of $k$.
\item $\sfname{LinkedUnl}$: under the same hypotheses, for every point $r$ of rank $3$ of $k$ which is irreducible on
every open subgroup, and every place $v\in V_\Phi$ at which $r$ is potentially crystalline, the local characters
$\psi_{v,i}$ of the point (characters of $G_v$) satisfy $\psi_{v,i}\ne\psi_{v,j}\modularitycyc$ for all $i,j$.
\end{enumerate}
\formallabel{UpperTheory.V.Z.Sym.Pl.LinkedEndEAlt}
\formallabel{UpperTheory.V.Z.J.LinkedUnlE}
\formallabel{UpperTheory.V.Pr.FinPt}
\formallabel{UpperTheory.V.Pr.AutPtAlt}
\formallabel{UpperTheory.V.Z.J.UnlPt}
\end{definition}

\begin{definition}[Descent and host endpoint statements]\label{pf:modularity:downdef}
\begin{enumerate}[label=(\alph*)]
\item $\sfname{DownCsd}$: for CM fields $K\subset L$ with $L/K$ Galois with solvable group, $n\ge1$, $\iota$, a complex
conjugation $c$, and a continuous $r:G_K\to\GL_n(\Qlbar)$ which is polarised with respect to $c$ and
$\modularitycyc^{1-n}$, irreducible on $G_L$, and satisfies the local conditions $\sfname{IneAb}$, $\sfname{Flagged}$
and $\sfname{Keeps}(\cdot,L)$ of \cref{pf:modularity:loccond}: if $r|_{G_L}$ is conjugate self-dual automorphic via $\iota$, so is
$r$.
\item $\sfname{HostEnd}$: for $K=\Q(i)$ with quadratic character $\eta_4$, every $\Omega$, $\ell$, $\iota$, $\OO$, $N$,
$e$, $\psi$, $o_1$ satisfying the host-data conditions, every complex conjugation $c$ and every valuation ring $B$ over
$\ell$, there is $g_1$, marked for $(\eta_4,\Psi_{e,\psi},c)$, such that $\sfname{FinEnd}(g_1,B)$ and
$\sfname{AutEnd}_3(g_1,B,\iota)$ hold for $(\eta_4,\Psi_{e,\psi},c)$ (\cref{pf:modularity:Uord}).
\end{enumerate}
\formallabel{UpperTheory.V.Z.Win.Dn.DownCsdKAlt}
\formallabel{UpperTheory.V.Pr.HostEndAlt}
\end{definition}

The main ingredients are the following results, proved in the formal development from $\sfname{HypothesesH'}$; we
state them in the form used.
\begin{enumerate}[label=(E\arabic*)]
\item\label{pf:modularity:E1} \emph{Seeds.} $\sfname{HypothesesH'}$ implies $\sfname{Seed2}$
(\lref{UpperTheory.V.Z.Win.Birth.seed2_alt}): for a pair of fields $(K,L)$ and characters $(\chi',\eta,\theta)$
satisfying the conditions of \cite{LeanSource2026}, and every bound $B$, there are $m\ge B$ with $m\equiv j\pmod{\ell-1}$
and a continuous $\sigma:G_\Q\to\GL_2(\Qlbar)$ with $\det\sigma=\eta\theta$, residually the sum of
$\modularitycyc^j\chi'\eta$ and $(\modularitycyc^j\chi')^{-1}$, whose restriction to $G_L$ is a \emph{seed} of weights
$(m,-m)$ (an automorphic, ordinary representation with prescribed residual and inertial characters). The statement
$\sfname{U3Dense}$, proved from $\sfname{HypothesesH'}$ in
\lref{UpperTheory.V.Z.Sym.W.unitary_three_automorphy_dense_of_prints}, gives three-dimensional seeds
$r:G_L\to\GL_3(\Qlbar)$, irreducible on all open subgroups, of weights $(m,0,-m)$.
\item\label{pf:modularity:E2} \emph{Linked chains.} $\sfname{HypothesesH'}$, $\sfname{Seed2}$ and $\sfname{U3Dense}$
imply $\sfname{LinkedEnd}$ and $\sfname{LinkedUnl}$ (\cref{pf:modularity:linked} below).
\item\label{pf:modularity:E3} \emph{Frames.} For every unitary datum $D$ and every $B$ over $\ell$ there are an abelian
CM field $L$, $g_1$, a frame $\Phi$ with $\sfname{Ok}(C_0)$, $\sfname{Big}(C_0)$ and the seed clause, and a Galois link
from $(\eta_D,\Psi,c_D,g_1,B,\ell N|d_K|)$ to $\Phi$ over $L$ (\lref{UpperTheory.V.Fr.gd_ex}); similarly for host
data over $\Q(i)$ with $M=\ell N$ (\lref{UpperTheory.V.Fr.host_ex}). Both use the global compatibility statement
deduced from \ref{pf:vwin-he:H1} (\cref{pf:modularity:loc}(d)).
\item\label{pf:modularity:E4} \emph{Local conditions} (\cref{pf:modularity:loccond}). For ordinary points $r$ over finite discrete valuation rings in
the presence of a Galois link: $r^\flat|_{G_K}$ satisfies $\sfname{IneAb}$ and $\sfname{Flagged}$
(\lref{UpperTheory.V.Z.Win.Dn.localConditionRankThreeLink_of_inertia_ordinary}); if moreover $N$ is even and
$r^\flat$ is irreducible on open subgroups of $\ker\eta$, it satisfies $\sfname{Keeps}(L)$
(\lref{UpperTheory.V.Z.Win.Dn.ends_rank_three_of_local'}, which uses $\sfname{LinkedUnl}$ and the local theorems);
the two-dimensional analogues are \lref{UpperTheory.V.Z.Win.Dn.potTwo_of} and
\lref{UpperTheory.V.Z.Win.Dn.inertiaConditionRankTwoLink_of_inertia}.
\item\label{pf:modularity:E5} \emph{Descent.} $\sfname{HypothesesH'}$ implies $\sfname{DownCsd}$
(\lref{UpperTheory.V.Z.Win.Dn.downCsdK_of_printsH}, through the solvable descent statement
\lref{UpperTheory.V.Af.solvable_descent_of_printed_hida}; cf.\ \cite{ArthurClozel1989}), the descent of automorphy
of rank two from abelian CM fields to $\Q$
(\lref{UpperTheory.V.Af.rankTwoAutomorphyDescends_of_printsH}), and the existence of newforms of weight $2e+1$ for
automorphic, ordinary two-dimensional representations of $G_\Q$ with Hodge--Tate weights $(e,-e)$
(\cref{pf:modularity:autnew}).
\end{enumerate}

\begin{proposition}\label{pf:modularity:linked}
$\sfname{HypothesesH'}$, $\sfname{Seed2}$ and $\sfname{U3Dense}$ imply $\sfname{LinkedEnd}$ and $\sfname{LinkedUnl}$.
\formallabel{UpperTheory.V.Z.Sym.Top.linkedEnd_alt}
\formallabel{UpperTheory.V.Z.J.linkedUnlE_holds_alt}
\end{proposition}

\begin{proof}
Both statements are instances of two general theorems on linked chains of the formal development:
$\sfname{LinkedEnd}$ is \lref{UpperTheory.Gal.projectedFormsSystemPart.linked_endomorphism_dense} and
$\sfname{LinkedUnl}$ is \cref{pf:modularity:linkedunl}. Their common hypotheses are:
global Artin reciprocity and local class field theory at all places (consequences of \ref{pf:vwin-he:H1},
\cref{pf:modularity:loc}(e)); the modularity statement
$\sfname{ModU}$ (\cref{pf:modularity:modU}) and the existence statements $\sfname{LevEx}$, $\sfname{LamEx}$,
$\sfname{ProModularRingExists}$ of \cref{pf:modularity:adjin}; the seeds $\sfname{U3Dense}$ and $\sfname{Seed2}$; and the
twist, cyclotomic and restriction properties of algebraic characters, consequences of $\sfname{Hypotheses'}$
(\lref{UpperTheory.V.Sd.twist_of_printed_results}, \cref{pf:modularity:algcyc},
\lref{UpperTheory.V.Sd.alg_restriction_of_printed_results}).
The second theorem has three further hypotheses, all consequences of $\sfname{HypothesesH'}$: modularity in rank three
for universal ordinary determinant rings over CM fields not containing $\mu_\ell$, $\sfname{ModN}$
(\cref{pf:modularity:modN}, which uses the multiplicity statement $\sfname{Mult}$ of \cref{pf:modularity:mult} and the
local-component statement $\sfname{CompOf}$, \lref{UpperTheory.V.Z.Mult.compOf_free}); the tangent-space statement
$\sfname{TanA}$ (\cref{pf:modularity:tanA}, from the global compatibility statement and the local theorems,
\cref{pf:modularity:loc}); and the vanishing of ordinary Selmer groups at the automorphic, locally generic points of the
universal ordinary determinant rings over CM fields, $\sfname{DenseThresholdMainCase}$ (\cref{pf:modularity:core}, from
$\sfname{HypothesesH'}$ and the cohomological consequence $\sfname{CohP}$ of \ref{pf:vwin-he:H1}).
\end{proof}

\begin{definition}[Local conditions]\label{pf:modularity:loccond}
Let $F$ be a number field and $r:G_F\to\GL_n(\Qlbar)$.
\begin{enumerate}[label=(\alph*)]
\item $\sfname{IneAb}(r)$: for every finite place $v\nmid\ell$ and every valuation ring $A$ over $v$ there are continuous
characters $\kappa_1,\dots,\kappa_n$ of $D_A$ such that the characteristic polynomial of $r(g)$ is
$\prod_i(X-\kappa_i(g))$ for $g\in I_A$.
\item $\sfname{Flagged}(r)$: for every place $v\mid\ell$ and every $A$ over $v$, $r|_{D_A}$ stabilises a full flag on
whose graded pieces $D_A$ acts by characters $\kappa_1,\dots,\kappa_n$.
\item $\sfname{Keeps}(r,L)$, for a finite extension $L/F$: for every place $w\mid\ell$ of $L$, if $r$ satisfies the
condition (i) of \cref{pf:modularity:aut}(b) at the place below $w$, then $r|_{G_L}$ satisfies it at $w$.
\end{enumerate}
$\sfname{FlagLaws}$ is the statement that for every number field $F$ and every $r:G_F\to\GL_n(\Qlbar)$: (1) if $r'\cong r$,
then $\sfname{IneAb}$, $\sfname{Flagged}$, unlinkedness and potential unlinkedness at each place above $\ell$
(\cref{pf:modularity:aut}(b)) pass from $r$ to $r'$; (2) for every character $\mu$ of $G_F$, $\sfname{IneAb}$ passes from
$r$ to $r\otimes\mu$ if $\mu$ is continuous, $\sfname{Flagged}$ passes from $r$ to $r\otimes\mu$, and $r$ is unlinked,
resp.\ potentially unlinked, at a place if and only if $r\otimes\mu$ is; (3) a representation of level $M>0$ and type $m$
satisfies $\sfname{IneAb}$ and $\sfname{Flagged}$; (4) for a finite extension $L/F$, $\sfname{IneAb}$ and $\sfname{Flagged}$
pass from $r$ to $r|_{G_L}$, and for every place $w\mid\ell$ of $L$ with restriction $v$ to $F$: if $r$ is potentially
unlinked at $v$ then $r|_{G_L}$ is potentially unlinked at $w$, and if $r$ satisfies $\sfname{Flagged}$, then potential
unlinkedness of $r|_{G_L}$ at $w$ implies that of $r$ at $v$, and unlinkedness of $r|_{G_L}$ at $w$ implies that of $r$ at $v$.
\formallabel{UpperTheory.V.Z.Win.FlagLaws}
\formal{UpperTheory.V.Z.Win.IneAb}
\formal{UpperTheory.V.Z.Win.Flagged}
\formal{UpperTheory.V.Z.Win.Keeps}
\end{definition}

$\sfname{FlagLaws}$ is a theorem of the formal development (\lref{UpperTheory.V.Z.Win.flagLaws}).

\begin{lemma}[Finiteness from the frame]\label{pf:modularity:finendpt}
Let $k$ be a Galois link from $(\eta,\Psi,c,g_1,B,M)$ to a frame $\Phi$ over $L$ satisfying $\sfname{Ok}(C_0)$, with
$\ell\mid M$, $c^2=1$, and the residue field of $\OO$ of characteristic $\ell$. If $\sfname{LinkedEnd}$(i) holds for
$\Phi$ (the characteristic polynomials of all rank-three lifts of the residual representation of $\Phi$ to complete
discrete valuation rings of characteristic $\ell$ are Eisenstein), then $\sfname{FinEnd}(g_1,B)$ holds.
\formallabel{UpperTheory.V.Pr.finEnd_of_finPt}
\end{lemma}

\begin{proof}
Take $U=G_L$, of index $[L:\Q]$ in $G_\Q$. Let $A$ be a complete discrete valuation ring over $\OO$ with finite residue
field in which $\mathfrak m_\OO$ maps to $0$, so that $A$ is an algebra over $k_\OO$, and let $r$ be an ordinary point
over $A$. By \lref{UpperTheory.V.Pr.exists_nrep3}, there is a lift $\rho$ of rank three of the residual
representation of $\Phi$, with inertial characters the powers $\bar\Psi^{\,1-i}$ ($i=0,1,2$), such that
$r(s)^\flat=\rho(p(s))$ for $s\in G_L$. By hypothesis the characteristic polynomial of $\rho(p(s))$ is
$\prod_i(X-\bar\Psi_\Phi(p(s))^{1-i})=\prod_i(X-\Psi(s)^{1-i})$, the second equality because $k$ is compatible with
$\bar\Psi$. This is the characteristic polynomial of $r_0^{\eta,\Psi}(s)^\flat=\operatorname{diag}(\Psi(s),1,\Psi(s)^{-1})$.
\end{proof}

\begin{lemma}[Automorphy over $K$ by descent]\label{pf:modularity:autend3}
Assume $\sfname{DownCsd}$. Let $k$ be a Galois link from $(\eta,\Psi,c,g_1,B,M)$ to a frame $\Phi$ over $L$ satisfying
$\sfname{Ok}(C_0)$, with $\ell\mid M$ and $c$ a complex conjugation; let $\iota$ be given, and let $K$ be a CM field such
that $G_K=\ker\eta$ for every embedding of $K$ in $\Omega$. Assume that for every discrete valuation ring
$A\subset\Qlbar$ finite over $\OO$ and every ordinary point $r$ over $A$ whose matrix part is irreducible on all open
subgroups of $\ker\eta$, the representation $r^\flat|_{G_K}$ satisfies $\sfname{IneAb}$, $\sfname{Flagged}$ and
$\sfname{Keeps}(\cdot,L)$; and that every irreducible point $r$ of rank $3$ of $k$ is such that
$r\otimes\modularitycyc^{-1}$ is conjugate self-dual automorphic via $\iota$. Then $\sfname{AutEnd}_3(g_1,B,\iota)$ holds for
$K$.
\formallabel{UpperTheory.V.Z.Win.Dn.autEnd3_alt}
\end{lemma}

\begin{proof}
Since $G_L\subset\ker\eta=G_K$, the field $K$ is contained in $L$, and $L/K$ is abelian because $L/\Q$ is; in particular
$\Gal(L/K)$ is solvable. We have $c^2=1$, and $\eta(c)=-1$ because $p(c)=c_\Phi$ and $\delta(c_\Phi)=-1$ by
$\sfname{Ok}(C_0)$.

Let $A$ and $r$ be as in $\sfname{AutEnd}_3$, $\rho=r^\flat|_{G_K}$ and $\rho_1=\rho\otimes\modularitycyc^{-1}$. By
\lref{UpperTheory.V.Pr.exists_nrep3}, $r$ gives a point of rank three of $k$ whose restriction to $G_L$ is
$\rho|_{G_L}$, with inertial characters $\Psi^{1-i}$. Since $G_L$ is an open subgroup of $\ker\eta$, $\rho|_{G_L}$ is
irreducible, so by hypothesis $\rho|_{G_L}\otimes\modularitycyc^{-1}=\rho_1|_{G_L}$ is conjugate self-dual automorphic
via $\iota$. We check the hypotheses of $\sfname{DownCsd}$ for $\rho_1$ and $L/K$:
\begin{itemize}
\item $\rho_1$ is continuous: $r$ is continuous for the $\mathfrak m_A$-adic topology and the entries lie in the integral
elements of $\Qlbar$;
\item $\rho_1$ is polarised with respect to $c$ and $\modularitycyc^{1-3}$: since $r(c)^\flat=1$ and $\epsilon(r(c))=-1$,
the relation $r(s')r(c)=r(c)r(s)$ for $s'=csc^{-1}$ gives, by the definition of the involution $\Theta$,
${}^t\rho(s)J\rho(s')=J$ for the antidiagonal matrix $J$ defining $\Theta$
(\lref{orthogonal_pairing_from_intertwiner}); twisting by $\modularitycyc^{-1}$, and using
$\modularitycyc(s')=\modularitycyc(s)$, gives ${}^t\rho_1(s)J\rho_1(s')=\modularitycyc(s)^{-2}J$;
\item $\rho_1|_{G_L}$ is irreducible, being a twist of $\rho|_{G_L}$;
\item $\rho$ satisfies $\sfname{IneAb}$, $\sfname{Flagged}$ and $\sfname{Keeps}(\cdot,L)$ by hypothesis, hence so does
the twist $\rho_1$ by $\sfname{FlagLaws}$.
\end{itemize}
By $\sfname{DownCsd}$, $\rho_1$ is conjugate self-dual automorphic via $\iota$.
\end{proof}

\begin{corollary}\label{pf:modularity:linkedH}
$\sfname{HypothesesH'}$ implies $\sfname{LinkedUnl}$.
\formallabel{UpperTheory.V.Z.J.linkedUnlE_alt}
\end{corollary}

\begin{proof}
By \ref{pf:modularity:E1}, $\sfname{HypothesesH'}$ implies $\sfname{Seed2}$ and $\sfname{U3Dense}$; apply
\cref{pf:modularity:linked}.
\end{proof}

The statement $\sfname{LinkedUnl}$ concerns the points of a frame. Other sections use it in the following form, which
concerns the ordinary points of the unitary deformation problems themselves. For a number field $K$ and $M\ge1$, a
character $\chi$ of $G_K$ \emph{kills $\mu_M$} if it is trivial on the elements of $G_K$ fixing all $M$-th roots of unity.

\begin{definition}[Levels at $\ell$ and the statement $(\mathrm{Unl})$]\label{pf:modularity:unldef}
Let $K$ be a number field, $\rho:G_K\to\GL_n(\Qlbar)$, $M\ge1$ and $m\in\Z^n$.
\begin{enumerate}[label=(\alph*)]
\item $\rho$ has \emph{level $M$ and type $m$ at $\ell$} if for every place $w\mid\ell$ of $K$ there are a valuation ring $A$
over $w$, characters $\chi_1,\dots,\chi_n$ of $G_K$ killing $\mu_M$, and a full flag of $\rho|_{D_A}$ with graded
characters $\kappa_1,\dots,\kappa_n$ such that $\kappa_i=\modularitycyc^{m_i}\chi_i$ on $I_A$, and such that either
$\kappa_i\ne\kappa_j\modularitycyc$ for all $i,j$, or $\rho$ is potentially crystalline at $A$.
\item $\rho$ is \emph{unlinked of level $M$ and type $m$} if for every place $w\mid\ell$ of $K$ there are $A$, $\chi_i$
and a flag as in (a) with $\kappa_i\ne\kappa_j\modularitycyc$ for all $i,j$.
\item $(\mathrm{Unl})$ is the statement: for all $\Omega$, $\ell$, $\OO$ with residue field of characteristic $\ell$,
$N\ge1$, $e$, $\psi$, $\eta:G_\Q\to\{\pm1\}$, $c,g_1\in G_\Q$, a valuation ring $B$ and $M\ge1$ such that the link
hypotheses hold for $(N,e,\eta,c,B,M)$, $N$ is even, $\Psi(c)=1$ for $\Psi=\Psi_{e,\psi}$, and a frame exists (there are
a CM field $L$, $g_2\in G_\Q$, a frame $\Phi$ satisfying $\sfname{Ok}(C_0)$ and $\sfname{Big}(C_0)$ for some $C_0$, a
Galois link from $(\eta,\Psi,c,g_2,B,M)$ to $\Phi$ over $L$, and the seed clause for $L$): for every discrete valuation
ring $A\subset\Qlbar$ finite over $\OO$, every ordinary point $r$ over $A$ for $(\eta,\Psi,c,g_1,B)$ whose matrix part
is irreducible on every open subgroup of $\ker\eta$, every imaginary quadratic field $K$ with $G_K=\ker\eta$, and every
$N'\ge1$: if $\rho=r^\flat|_{G_K}$ has level $N'$ and type $(e,0,-e)$ at $\ell$, then $\rho$ is unlinked of level $N'$
and type $(e,0,-e)$.
\end{enumerate}
\formal{UpperTheory.V.Af.LevelAt}
\formal{UpperTheory.V.Th.Unl}
\formal{UpperTheory.V.Pr.HasFrame}
\formallabel{UpperTheory.V.Z.J.UnlEndLogNorm}
\end{definition}

\begin{lemma}\label{pf:modularity:flagmatch}
Let $k$ be a field, $H$ a set, $T:H\to M_3(k)$ and $\psi_0,\psi_1,\psi_2,D_0,D_1,D_2:H\to k$. Assume that there is
$Q\in\GL_3(k)$ such that $QT(h)Q^{-1}$ is upper triangular with diagonal $(D_0(h),D_1(h),D_2(h))$ for all $h\in H$, and
that $(T(h_0)-\psi_0(h_0))(T(h_1)-\psi_1(h_1))(T(h_2)-\psi_2(h_2))=0$ for all $h_0,h_1,h_2\in H$. Then for every $b$
there is $a$ with $D_b=\psi_a$.
\formallabel{UpperTheory.V.Z.J.UnitaryRankTwoFlags.flag_match}
\end{lemma}

\begin{proof}
Otherwise there is $b$ such that for each $a$ there is $h_a\in H$ with $D_b(h_a)\ne\psi_a(h_a)$. The matrices
$U_a=Q(T(h_a)-\psi_a(h_a))Q^{-1}=QT(h_a)Q^{-1}-\psi_a(h_a)$ are upper triangular with entry $D_b(h_a)-\psi_a(h_a)\ne0$ at
$(b,b)$, and $U_0U_1U_2=Q\,(T(h_0)-\psi_0(h_0))(T(h_1)-\psi_1(h_1))(T(h_2)-\psi_2(h_2))\,Q^{-1}=0$. But the entry $(b,b)$
of a product of upper triangular matrices is the product of their entries $(b,b)$, which is non-zero; a contradiction.
\end{proof}

\begin{lemma}\label{pf:modularity:move}
Let $K$ be an imaginary quadratic field, $R:G_\Q\to\modularityG(\Qlbar)$ a homomorphism whose sign is trivial on $G_K$,
and $\rho=R^\flat|_{G_K}$. Let $u\in G_\Q$, $P$ a valuation ring of $\Omega$ and $P'=u^{-1}(P)$. Suppose that $\rho$ is
potentially crystalline at $P$ and that $\rho|_{D_P}$ (decomposition group in $G_K$) stabilises a full flag with graded
characters $\kappa_0,\kappa_1,\kappa_2$ such that $\kappa_i=\kappa_j\modularitycyc$ for some $i,j$. Then $\rho$ is
potentially crystalline at $P'$, and $\rho|_{D_{P'}}$ stabilises a full flag with graded characters
$\kappa'_0,\kappa'_1,\kappa'_2$ such that $\kappa'_{i'}=\kappa'_{j'}\modularitycyc$ for some $i',j'$.
\formallabel{UpperTheory.V.Z.J.UnitaryRankTwoFlags.move_all}
\end{lemma}

\begin{proof}
As $G_K$ is normal in $G_\Q$, $c(s)=usu^{-1}$ defines an automorphism of $G_K$; it maps $D_{P'}$ into $D_P$, and
$\modularitycyc(c(s))=\modularitycyc(s)$. By \lref{UpperTheory.potentiallyCrystalline_comp_conjugation} (potential
crystallinity at $P$ is transported by conjugation by $u$ to potential crystallinity at $P'$), $\rho\circ c$ is
potentially crystalline at $P'$. The characters $\kappa_a\circ c$ of $D_{P'}$ satisfy
$\kappa_i\circ c=(\kappa_j\circ c)\modularitycyc$. Write $R(u)=(x,\epsilon)$. If $\epsilon=1$, then
$\rho(c(s))=x\rho(s)x^{-1}$; so $\rho\cong\rho\circ c$ is potentially crystalline at $P'$, and the given flag transported by
$x^{-1}$ is a full flag of $\rho|_{D_{P'}}$ with graded characters $\kappa'_a=\kappa_a\circ c$; take $(i',j')=(i,j)$. If
$\epsilon=-1$, then $\rho(c(s))=x\Theta(\rho(s))x^{-1}$; so $\rho$ is isomorphic (by the matrix $\Theta(x)^{-1}J$) to the
contragredient $s\mapsto{}^t\rho(c(s))^{-1}$ of $\rho\circ c$, which is potentially crystalline at $P'$
(\lref{UpperTheory.potentiallyCrystalline_contragredient}), hence so is $\rho$; and the given flag yields a full flag
of $\rho|_{D_{P'}}$ with graded characters $\kappa'_a=(\kappa_{2-a}\circ c)^{-1}$, for which
$\kappa'_{2-j}=\kappa'_{2-i}\modularitycyc$.
\end{proof}

\begin{theorem}\label{pf:modularity:unlend}
$\sfname{LinkedUnl}$ implies $(\mathrm{Unl})$. In particular, by \cref{pf:modularity:linkedH}, $\sfname{HypothesesH'}$
implies $(\mathrm{Unl})$.
\formallabel{UpperTheory.V.Z.J.unlEnd3E}
\end{theorem}

\begin{proof}
Let the data of $(\mathrm{Unl})$ be given, and fix $L$, $g_2$, $\Phi$, $C_0$ and a Galois link $k$ from
$(\eta,\Psi,c,g_2,B,M)$ to $\Phi$ over $L$ as in the existence of a frame; write $p:G_{L^+}\to\Gamma_\Phi$ for its
surjection, $B_v$ ($v\in V_\Phi$) for its valuation rings over $\ell$ and $j_v:G_v\to G_{L^+}$ for its local embeddings.

\emph{Step 1: renormalisation.} We have $\ell\mid M$ (as $\ell N\mid M$), $c^2=1$, $g_2$ is marked for $(\eta,\Psi,c)$
(\lref{UpperTheory.V.Pr.marks_of_link}), and $2\in\OO^\times$ ($\ell$ is odd by $\sfname{Ok}(C_0)$). By
\lref{UpperTheory.V.Th.End3.ordPt_renorm} (an ordinary point becomes, after conjugation by a matrix congruent to $1$,
an ordinary point for any other marked element) there is $h\in\GL_3(A)$, $h\equiv1$, such that $r'=hrh^{-1}$ is an
ordinary point for $(\eta,\Psi,c,g_2,B)$. Let $h'$ be the image of $h$ in $\GL_3(\Qlbar)$ and $\rho'=h'\rho h'^{-1}$.
Since the sign of $r$ is trivial on $\ker\eta=G_K$, $\rho'=r'^\flat|_{G_K}$; and $r'^\flat$ is irreducible on every open
subgroup of $\ker\eta$, as $r^\flat$ is. Moreover $G_L\subset\ker\eta=G_K$ (\lref{UpperTheory.V.Th.Db.galL_le}).

\emph{Step 2: the point of the frame.} By \lref{UpperTheory.V.Pr.exists_nrep3}, there is a lift $\nu$ of rank three of
the residual representation of $\Phi$ over $A$ with $\nu(p(s))=r'(s)^\flat$ for $s\in G_L$ and with local characters
$\nu_{v,i}:G_v\to A^\times$ satisfying $\nu_{v,i}(g)=\Psi(j_v(g))^{1-i}$ for $g\in I_v$. Being a lift of rank three, $\nu$
satisfies the ordinary condition
$\prod_{a=0}^2\bigl(\nu(p(j_v(h_a)))-\nu_{v,a}(h_a)\bigr)=0$ for all $v$ and $h_0,h_1,h_2\in G_v$. The representation
$\rho'|_{G_L}=\nu\circ p|_{G_L}$ is irreducible on every open subgroup $U$ of $G_L$, because $U$ is open in $\ker\eta$
(\lref{UpperTheory.V.Th.Db.irr_of_matIrr}).

\emph{Step 3.} Let $w\mid\ell$ be a place of $K$, and let $A_w=P$, the characters $\chi_i$ and a flag of $\rho|_{D_P}$
with graded characters $\kappa_i$ be given by the hypothesis that $\rho$ has level $N'$ and type $(e,0,-e)$. We show that
the same data witness that $\rho$ is unlinked at $w$; only $\kappa_i\ne\kappa_j\modularitycyc$ needs proof, and by the
hypothesis we may assume that $\rho$ is potentially crystalline at $P$. Suppose $\kappa_i=\kappa_j\modularitycyc$. Since
$P$ lies over $\ell$, the link gives $v\in V_\Phi$ and $u\in G_{L^+}$ with $P=u(B_v)$. Transporting by $h'$, $\rho'$ is
potentially crystalline at $P$ and $\rho'|_{D_P}$ has a flag with the same graded characters. By
\cref{pf:modularity:move}, applied to $r'$ (composed with $A\to\Qlbar$), $\rho'$ is potentially crystalline at $B_v$, and
$\rho'|_{D_{B_v}}$ has a full flag with graded characters $\kappa'_a$ satisfying $\kappa'_{i'}=\kappa'_{j'}\modularitycyc$
for some $i',j'$.

\emph{Step 4: comparison with the frame.} As $B_v$ lies over $\ell\mid M$, the link gives
$D_{B_v}\cap G_{L^+}\subset\ker\eta$. Hence $j_v(g)\in D_{B_v}\cap G_L$ for $g\in G_v$; in particular $j_v(g)$ lies in the
decomposition group of $B_v$ in $G_K$, and $\rho'(j_v(g))=\nu(p(j_v(g)))$ by Step 2. Apply \cref{pf:modularity:flagmatch}
(over $\Qlbar$) to $H=G_v$, $T(g)=\nu(p(j_v(g)))$, $\psi_a=\nu_{v,a}$, and $D_b(g)=\kappa'_b(j_v(g))$, the diagonal of
$\rho'(j_v(g))$ in a basis adapted to the flag: there are $a,b$ with $\kappa'_{i'}\circ j_v=\nu_{v,a}$ and
$\kappa'_{j'}\circ j_v=\nu_{v,b}$. Hence $\nu_{v,a}(g)=\nu_{v,b}(g)\modularitycyc(j_v(g))$ for all $g\in G_v$.

\emph{Step 5: contradiction.} $\nu$ is a point of rank three of $k$ whose representation $\rho'|_{G_L}$ is irreducible on
every open subgroup (Step 2) and potentially crystalline at $B_v$, since $\rho'$ is and potential crystallinity is
preserved by restriction to $G_L$ (\lref{UpperTheory.V.Br.potRes}). As $N$ is even, $\sfname{LinkedUnl}$ applies to
$k$, $\nu$ and $v$ and gives $g\in G_v$ with $\nu_{v,a}(g)\ne\nu_{v,b}(g)\modularitycyc(j_v(g))$, contradicting Step 4.
\end{proof}

\begin{proposition}\label{pf:modularity:gdend}
The statements \ref{pf:modularity:E3}--\ref{pf:modularity:E5}, $\sfname{LinkedEnd}$ and $\sfname{HypothesesH'}$ imply
$(\mathrm{U}_{\mathrm{ord}})$.
\formallabel{UpperTheory.V.Z.Win.Dn.gdEnd_of_frame_linkedEnd_downCsdK_descent}
\end{proposition}

\begin{proof}
Let $D$ be a unitary datum for $(K,\ell,\OO,e,\psi)$, $\iota$, and $B$ over $\ell$. The residue field of $\OO$ has
characteristic $\ell$. By \ref{pf:modularity:E3} there are $L$, $g_1$, $\Phi$, $C_0$ and a Galois link $k$ from
$(\eta,\Psi,c,g_1,B,M)$, $M=\ell N|d_K|$, to $\Phi$ over $L$. Since $c$ is a complex conjugation, $c^2=1$, and
$\Psi(c)=\modularitycyc(c)^e\psi(-1)=(-1)^e(-1)^e=1$ by the parity condition of $\sfname{GD}$. The character $\eta$ is
ramified only at primes dividing $d_K$, since $\ker\eta=G_K$; so the link hypotheses hold for $(N,e,\eta,c,B,M)$.
$\sfname{LinkedEnd}$ applied to $k$ gives (i) finiteness, (ii) automorphy of the irreducible points of rank $3$ of $k$
twisted by $\modularitycyc^{-1}$, and (iii) an algebraic $\kappa$ with automorphy of the irreducible points of rank $2$
twisted by $\kappa$. Now:
\begin{itemize}
\item $g_1$ is marked (\lref{UpperTheory.V.Pr.marks_of_link}), and (i) gives $\sfname{FinEnd}(g_1,B)$ by
\cref{pf:modularity:finendpt}.
\item As $K$ is quadratic, all embeddings of $K$ in $\Omega$ have the same image, so $G_K=\ker\eta$ for each of them.
Hence $\sfname{AutEnd}_3(g_1,B,\iota)$ follows from \cref{pf:modularity:autend3}, whose hypotheses are (ii), the local
conditions of \ref{pf:modularity:E4} and $\sfname{DownCsd}$.
\item $\sfname{AutEnd}_2(g_1,B,\iota)$ follows from (iii), the two-dimensional local conditions of
\ref{pf:modularity:E4}, the statements of \ref{pf:modularity:E5}, and the properties of algebraic characters
(\lref{UpperTheory.V.Z.Win.Dn.autEnd2_alt}; the passage from the formulation there to that of
\cref{pf:modularity:Uord} is \lref{UpperTheory.V.Z.Win.Dn.end2'}).
\end{itemize}
This gives $(\mathrm{U}_{\mathrm{ord}})$ for $D$.
\end{proof}

\begin{proposition}\label{pf:modularity:hostend}
\ref{pf:modularity:E3}, \ref{pf:modularity:E4}, $\sfname{LinkedEnd}$ and $\sfname{DownCsd}$ imply $\sfname{HostEnd}$.
\formallabel{UpperTheory.V.Z.Win.Dn.hostEnd_of_loc3_frame_linkedEnd_downCsdK}
\end{proposition}

\begin{proof}
The argument is that of the first two items of the proof of \cref{pf:modularity:gdend}, with the frame of
\ref{pf:modularity:E3} for host data and $M=\ell N$: the frame and link give $g_1$; $c^2=1$; $\sfname{LinkedEnd}$ gives
finiteness and automorphy of the rank-three points over $L$; $g_1$ is marked; finiteness gives $\sfname{FinEnd}$; and
descent from $L$ to $\Q(i)$ gives $\sfname{AutEnd}_3$.
\end{proof}

\begin{theorem}\label{pf:modularity:symcube}
$\sfname{HostEnd}$ implies $(\mathrm{U}_i)$.
\formallabel{UpperTheory.V.Pr.symmetric_cube_theorem_of_endomorphisms_alt}
\end{theorem}

\begin{proof}
Let $K=\Q(i)$, the data $\Omega,\ell,\iota,\OO,N,e,\psi,o_1$, $S$, the valuation rings $A_v$ and a host frame $F$ be
given. The frame comes with a complex conjugation $c$ mapping to its conjugation element, and the valuation ring
$A_{v_\ell}$ lies over $\ell$. By $\sfname{HostEnd}$ (for $c$ and $B=A_{v_\ell}$) there is $g_1$, marked for
$(\eta_4,\Psi,c)$, with $\sfname{FinEnd}(g_1,B)$ (with a subgroup $U$ of finite index) and $\sfname{AutEnd}_3(g_1,B,\iota)$.
The multiplier and character of $F$ are $\eta_4$ and $\Psi$ composed with $G_\Q\to G_{\Q,S}$, so the image of $g_1$ in
$G_{\Q,S}$ is marked for $F$. Let $u:G_{\Q,S}\to\modularityG(R)$ be a universal framed deformation of $F$ at $g_1$.

\emph{Finiteness.} By \lref{UpperTheory.V.UnitaryThreeInertia.universal_diagonal_of_triangular_points} (applied
with \lref{UpperTheory.V.UnitaryThreeInertia.hostFrame_ok} and \lref{UpperTheory.V.UnitaryThreeInertia.frame_class_one}),
the part of weight $0$ of $R$ is a finite $\OO$-module provided every point of $u$ with values in a complete discrete
valuation ring of characteristic $\ell$ is upper or lower triangular on $\ker\eta_4$. Such a point, composed with
$G_\Q\to G_{\Q,S}$, is an ordinary point (\lref{UpperTheory.V.Pr.ordPt_of_isPt}); by $\sfname{FinEnd}(g_1,B)$ its
characteristic polynomials on $U\cap\ker\eta_4$ are Eisenstein, and then it is triangular (\lref{tri_of_char}).

\emph{Automorphy.} Let $\xi$ be a $\Qlbar$-point of the subalgebra generated by the homogeneous parts, with
$r=\xi_\ast u$ of type I. By the finiteness just proved, $r$ is a characteristic-zero point
(\lref{UpperTheory.V.UnitaryThreeInertia.isCharZeroPt_of_isPtHom}): there are a discrete valuation ring $A'$ finite
over $\OO$, a unit $t$ and a point $r'$ over $A'$ with $r'=\mathrm{tor}(t)\,r\,\mathrm{tor}(t)^{-1}$; $r'$ is an ordinary
point (\lref{UpperTheory.V.Pr.ordPt_of_isPt}); and $r^\flat$ is irreducible on every open subgroup of
$\ker\eta_4$ (\lref{UpperTheory.V.UnitaryThreeInertia.matIrr_open}), hence so is $r'^\flat$. Since
$\ker\eta_4=G_{\Q(i)}$ (\lref{UpperTheory.V.UnitaryThreeInertia.gaussian_field_range_eq_ker_quartic_char}),
$\sfname{AutEnd}_3(g_1,B,\iota)$ shows that $r'^\flat|_{G_K}\otimes\modularitycyc^{-1}$ is conjugate self-dual automorphic,
and $r^\flat|_{G_K}\otimes\modularitycyc^{-1}$ is isomorphic to it via $\mathrm{tor}(t)$; we conclude by
\cref{pf:modularity:iso}.
\end{proof}

\begin{theorem}\label{pf:modularity:endsof}
$\sfname{HypothesesH'}$, $\sfname{LinkedEnd}$ and $\sfname{LinkedUnl}$ imply $(\mathrm{U}_{\mathrm{ord}})$ and $(\mathrm{U}_i)$.
\formallabel{UpperTheory.V.Z.Win.Dn.ends_of_alt}
\end{theorem}

\begin{proof}
By \ref{pf:modularity:E3}--\ref{pf:modularity:E5} (where \ref{pf:modularity:E4} uses $\sfname{LinkedUnl}$),
\cref{pf:modularity:gdend} gives $(\mathrm{U}_{\mathrm{ord}})$, and \cref{pf:modularity:hostend} and
\cref{pf:modularity:symcube} give $(\mathrm{U}_i)$.
\end{proof}

\begin{corollary}\label{pf:modularity:endsu}
$\sfname{HypothesesH'}$ and $\sfname{LinkedUnl}$ imply $(\mathrm{U}_{\mathrm{ord}})$ and $(\mathrm{U}_i)$.
\formallabel{UpperTheory.V.Z.Win.Dn.ends_u_alt}
\end{corollary}

\begin{proof}
By \ref{pf:modularity:E1} and the first part of \ref{pf:modularity:E2}, $\sfname{HypothesesH'}$ implies
$\sfname{LinkedEnd}$; apply \cref{pf:modularity:endsof}.
\end{proof}

\begin{theorem}[Step \ref{pf:vwin-he:b} of \cref{sec:pf:vwin-he:proof}]\label{pf:modularity:ends0}
$\sfname{HypothesesH'}$ implies $(\mathrm{U}_{\mathrm{ord}})$ and $(\mathrm{U}_i)$.
\formallabel{UpperTheory.V.Z.Win.Dn.ends0_alt}
\end{theorem}

\begin{proof}
By \ref{pf:modularity:E1}, $\sfname{HypothesesH'}$ implies $\sfname{Seed2}$ and $\sfname{U3Dense}$; by the second part of
\ref{pf:modularity:E2} it then implies $\sfname{LinkedUnl}$; and \cref{pf:modularity:endsu} gives the conclusion.
\end{proof}

\begin{remark}\label{pf:modularity:rem-ends}
(1) The statement $\sfname{U3Dense}$, the density of automorphic points for unitary groups in three variables, is the
root of a different part of the proof (\lref{UpperTheory.V.Z.Sym.W.unitary_three_automorphy_dense_of_prints}); it
is used here only through \ref{pf:modularity:E2}.
(2) $\sfname{LinkedUnl}$, in the form of \cref{pf:modularity:linkedH} and \cref{pf:modularity:unlend}, is also used in the
proofs of the counts for cases I and V. The results entering its proof are proved in \cref{sec:pf:modularity:modU}, as
assemblies of long results of the formal development which we state there.
\end{remark}

\subsection{Complements on automorphy}\label{sec:pf:modularity:compl}

This subsection proves auxiliary statements on automorphy which are used in this and other sections: algebraicity of
cyclotomic and finite order characters, potential unlinkedness of ordinary representations, the behaviour of the window
condition under isomorphisms and twists, and the existence of newforms attached to ordinary automorphic representations
of $G_\Q$ of dimension two.

\begin{definition}[Algebraic characters]\label{pf:modularity:algchar}
Let $F$ be a number field and $\iota:\Qlbar\cong\C$. A \emph{global Artin map} for $F$ is a continuous open surjection
$\mathrm{art}:\A_F^\times\to G_F^{\mathrm{ab}}$ whose kernel is the closure of the product of $F^\times$ and the identity
component of $(F\otimes\R)^\times$, and which maps the local units, resp.\ uniformisers, resp.\ local multiplicative groups
at each finite place onto the images of the inertia groups, resp.\ (inverse) Frobenius elements, resp.\ (topologically)
the decomposition groups. A Hecke character $\chi$ of $F$ is \emph{algebraic of type} $a=(a_\tau)_{\tau:F\to\C}$ if
$\chi_\infty(x)=\prod_\tau\tau(x)^{a_\tau}$ for every totally positive $x\in F^\times$. A character
$\mu:G_F\to\Qlbar^\times$ is \emph{algebraic} (via $\iota$) if it is continuous and there are a global Artin map
$\mathrm{art}$, a Hecke character $\chi$ and $a$ such that $\chi$ is algebraic of type $a$ and $\mu$ is the
\emph{$\ell$-adic character of $(\chi,a)$}: for $x\in\A_F^\times$ and $\sigma\in G_F$ with $\mathrm{art}(x)$ equal to the
image of $\sigma$, $\mu(\sigma)$ is obtained from $\iota^{-1}(\chi(x)x_\infty^{-a})$ by multiplication by $x_\ell^{a}$,
where $x_\infty^{a}=\prod_\tau\tau(x_\infty)^{a_\tau}$ and $x_\ell^a$ is its $\ell$-adic analogue (through
$\iota^{-1}\circ\tau$). $\sfname{AlgCyc}$ is the statement: for every number field $F$, $\ell$ and $\iota$, (a)
$\modularitycyc^j\chi$ is algebraic for every $j\in\Z$ and every character $\chi:G_F\to\Qlbar^\times$ with open kernel;
(b) products of algebraic characters are algebraic; (c) inverses of algebraic characters are algebraic.
\formal{GlobalTheory.HeckeGal.IsGlobalArtin}
\formal{GlobalTheory.HeckeGal.IsAlgPos}
\formal{GlobalTheory.HeckeGal.IsLAdicOf}
\formal{UpperTheory.V.Sd.IsAlgChar}
\formal{UpperTheory.V.Af.IsAlgChar}
\formal{UpperTheory.V.Sd.AlgCyc}
\end{definition}

\begin{proposition}\label{pf:modularity:algcyc}
$\sfname{Hypotheses'}$ implies $\sfname{AlgCyc}$.
\formallabel{UpperTheory.V.Sd.alg_cyclotomic_of_printed_results}
\end{proposition}

\begin{proof}
The global Artin reciprocity statement contained in \ref{pf:vwin-he:H1}(b) provides a global Artin map for every number
field. The defining identities of algebraicity and of the $\ell$-adic character are multiplicative in $(\chi,a,\mu)$.

(c) If $\mu$ is algebraic with data $(\mathrm{art},\chi,a)$, then $\mu^{-1}$ is continuous, $\chi^{-1}$ is algebraic of
type $-a$, and $\mu^{-1}$ is the $\ell$-adic character of $(\chi^{-1},-a)$.

(b) If $\mu,\mu'$ are algebraic with data $(\mathrm{art},\chi,a)$ and $(\mathrm{art}',\chi',a')$, then
$\mathrm{art}=\mathrm{art}'$ by the uniqueness of global Artin maps (\lref{GlobalTheory.ArtinUniq.globalArtin_unique}),
$\chi\chi'$ is algebraic of type $a+a'$, and $\mu\mu'$ is the $\ell$-adic character of $(\chi\chi',a+a')$.

(a) By (b) it suffices to show that $\modularitycyc^j$ and $\chi$ are algebraic. For $\modularitycyc^j$ this is
\lref{UpperTheory.exists_algebraic_heckeChar_cyclotomic} (given global Artin reciprocity, every integral power of the
cyclotomic character is the $\ell$-adic character of an algebraic Hecke character). As $\ker\chi$ is open,
$d=[G_F:\ker\chi]$ is finite and $\chi^d=1$; by class field theory \cite{Neukirch1999}
(\lref{GlobalTheory.HeckeGal.exists_unique_hchar_of_finite_order}) there is a Hecke character $\chi_0$ of order
dividing $d$ such that $\chi$ is the $\ell$-adic character of $(\chi_0,0)$, and a Hecke character of finite order is
algebraic of type $0$ (\lref{GlobalTheory.HeckeGal.isAlgPos_zero_of_orderDvd}).
\end{proof}

\begin{lemma}\label{pf:modularity:potflag}
Let $F$ be a number field, $r:G_F\to\GL_n(\Qlbar)$, $m\in\Z^n$ with $m_i\ne m_j+1$ for all $i,j$, and $v\mid\ell$ a place of
$F$. Assume that for every valuation ring $A$ over $v$ there are a full flag of $r|_{D_A}$, with graded characters
$\kappa_1,\dots,\kappa_n$, and $a\ge1$ such that $\kappa_i(s)^a=\modularitycyc(s)^{am_i}$ for all $i$ and all $s\in I_A$.
Then $r$ is potentially unlinked at $v$.
\formallabel{UpperTheory.V.Z.ImaginaryQuadraticEtaCharacter.potOfFlag}
\formal{UpperTheory.V.Z.ImaginaryQuadraticEtaCharacter.PotOfFlag}
\end{lemma}

\begin{proof}
Let $A$ be over $v$, and $\kappa$, $a$ as in the hypothesis. Suppose that $\kappa_i^b=(\kappa_j\modularitycyc)^b$ for some
$i,j$ and $b\ge1$. As $A$ lies over $\ell$, the image of $I_A$ under $\modularitycyc$ is an open subgroup of $\Zl^\times$,
so there is $s\in I_A$ such that $\modularitycyc(s)$ has infinite order. Raising
$\kappa_i(s)^b=(\kappa_j(s)\modularitycyc(s))^b$ to the power $a$ gives
$\modularitycyc(s)^{abm_i}=\modularitycyc(s)^{ab(m_j+1)}$, hence $ab(m_i-m_j-1)=0$ and $m_i=m_j+1$, a contradiction.
\end{proof}

\begin{lemma}\label{pf:modularity:potord}
Let $F$ be a number field, $r:G_F\to\GL_n(\Qlbar)$, $m\in\Z^n$ with $m_i\ne m_j+1$ for all $i,j$, $\psi_1,\dots,\psi_n$
characters of $G_F$, and $\kappa$ a character of $G_F$. Assume that for every valuation ring $A$ of $\Omega$ over $\ell$,
$r|_{D_A}$ stabilises a full flag whose graded characters are equal to $\modularitycyc^{m_i}$ on $I_A$ and congruent to
$\psi_i$. Then $r\otimes\kappa$ is potentially unlinked at every place $v\mid\ell$ of $F$.
\formallabel{UpperTheory.V.Z.Win.Birth.potOrd}
\formal{UpperTheory.V.Z.Win.Birth.PotOrd}
\formal{UpperTheory.V.Sd.OrdAt}
\end{lemma}

\begin{proof}
For every $A$ over $v$, the flag given by the hypothesis satisfies the assumption of \cref{pf:modularity:potflag} with
$a=1$; hence $r$ is potentially unlinked at $v$. By $\sfname{FlagLaws}$ (\cref{pf:modularity:loccond}), which is a theorem
(\lref{UpperTheory.V.Z.Win.flagLaws}), potential unlinkedness at $v$ is invariant under twisting, so $r\otimes\kappa$ is
potentially unlinked at $v$. (The congruences with the $\psi_i$ are not used.)
\end{proof}

\begin{definition}[Twisting statements]\label{pf:modularity:wintw}
\begin{enumerate}[label=(\alph*)]
\item $\sfname{WinTw}$ is the statement: for every totally complex number field $K$ and every Hecke character $\chi$ of $K$
there is $a\in\R$ such that for all $n\ge1$, all cuspidal automorphic representations $\pi$ of $\GL_n(\A_K)$, all $t\in\R$
and all finite places $v$: if $\pi$ satisfies the window condition with centre $t$ at $v$, then $\pi\otimes\chi$ satisfies
it with centre $t+a$.
\item $\sfname{AIsoTw}$ is the statement: $\sfname{Hypotheses'}$ implies that for every CM field $F$, $\iota$, $n$ and
$r:G_F\to\GL_n(\Qlbar)$ satisfying $\sfname{WitA}(\iota)$, every $r'\cong r$ satisfies $\sfname{WitA}(\iota)$, and
$r\otimes\mu$ satisfies $\sfname{WitA}(\iota)$ for every algebraic character $\mu$ of $G_F$.
\end{enumerate}
\formallabel{UpperTheory.V.Z.ImaginaryQuadraticEtaCharacter.WinTw}
\formallabel{UpperTheory.V.Z.Win.AIsoTwAlt}
\end{definition}

\begin{proposition}[Twisting property of the window condition]\label{pf:modularity:aisotw}
$\sfname{FlagLaws}$ and $\sfname{WinTw}$ imply $\sfname{AIsoTw}$.
\formallabel{UpperTheory.V.Z.Win.aIsoTwOf_alt}
\formal{UpperTheory.V.Z.Win.AIsoTwOfAlt}
\end{proposition}

\begin{proof}
Let $F$, $\iota$, $n$ and $r$ satisfying $\sfname{WitA}(\iota)$ be given.

\emph{Isomorphic representations.} Let $r'\cong r$. If $r'$ is unlinked at $v\mid\ell$, then so is $r$ ($\sfname{FlagLaws}$
(1) applied to $r'$), hence $r$ is potentially unlinked at $v$ by condition (i) for $r$, and so is $r'$ ($\sfname{FlagLaws}$
(1) applied to $r$). For condition (ii), let $\pi$ be regular algebraic cuspidal and attached to $r'$. Then $\pi$ is
attached to $r$, since attachment only involves semisimplicity and characteristic polynomials. Condition (ii) for $r$ gives
$t$ such that $\pi$ satisfies the window condition with centre $t$ at the places $v\nmid\ell$ and at the places where $r$ is
potentially unlinked, and the latter include the places where $r'$ is potentially unlinked ($\sfname{FlagLaws}$ (1)).

\emph{Twists.} Let $\mu$ be algebraic, with data $(\mathrm{art},\chi,a)$. If $n=0$ then $r\otimes\mu=r$. Let $n\ge1$. By
$\sfname{FlagLaws}$ (2), for every $v\mid\ell$, $r$ is unlinked at $v$ if and only if $r\otimes\mu$ is, and similarly for
potential unlinkedness; so condition (i) for $r\otimes\mu$ follows from condition (i) for $r$. For condition (ii), let
$\pi'$ be regular algebraic cuspidal and attached to $r\otimes\mu$. Since $\mu^{-1}$ is the $\ell$-adic character of
$(\chi^{-1},-a)$, the twist $\pi'\otimes\chi^{-1}$ is attached to $(r\otimes\mu)\otimes\mu^{-1}=r$
(\lref{UpperTheory.V.Af.Cq.attTwist}: if $\rho$ is attached to $\pi'$ and $\mu'$ is the $\ell$-adic character of a
Hecke character $\chi'$, then $\rho\otimes\mu'$ is attached to $\pi'\otimes\chi'$), and $\pi'\otimes\chi^{-1}$ is regular
algebraic because $\chi^{-1}$ is algebraic (\lref{UpperTheory.V.Z.Win.regAlg_tw}: twisting by an algebraic Hecke
character preserves regular algebraicity). Condition (ii) for $r$ gives $t$ such that $\pi'\otimes\chi^{-1}$ satisfies the
window condition with centre $t$ at the places $v\nmid\ell$ and at the places where $r$, equivalently $r\otimes\mu$, is
potentially unlinked. A CM field is totally complex, so $\sfname{WinTw}$ applied to $\chi$ gives $b$ such that
$\pi'=(\pi'\otimes\chi^{-1})\otimes\chi$ satisfies the window condition with centre $t+b$ at the same places.
\end{proof}

\begin{remark}\label{pf:modularity:rem-aisotw}
In the formal statement $\sfname{AIsoTw}$ the hypothesis $\sfname{Hypotheses'}$ is not used: the proof derives the
conclusion from $\sfname{FlagLaws}$ and $\sfname{WinTw}$ alone.
\end{remark}

Finally we record the passage from automorphic representations of $\GL_2(\A_\Q)$ to newforms used in
\ref{pf:modularity:E5}.

\begin{definition}\label{pf:modularity:autnewdef}
Let $\iota:\Qlbar\cong\C$. A continuous $W:G_\Q\to\GL_2(\Qlbar)$ is \emph{weakly unlinked} (via $\iota$) if $\det W$ has
cyclotomic weight $0$ and there are an abelian CM field $L$ such that $W|_{G_L}$ is irreducible, a character
$\eta:G_\Q\to\{\pm1\}$ trivial on $G_L$ with $\eta(c)=-1$ for some complex conjugation $c$, and an algebraic character
$\kappa$ of $G_L$ such that $W|_{G_L}\otimes\kappa$ satisfies $\sfname{WitA}(\iota)$. $\sfname{AutNewform}$ is the statement:
for all $\Omega$, $\ell$, $\iota$ and $e\ge1$, every continuous irreducible $W:G_\Q\to\GL_2(\Qlbar)$ which is attached via
$\iota$ to a regular algebraic cuspidal automorphic representation of $\GL_2(\A_\Q)$ and is weakly unlinked, every
valuation ring $B$ over $\ell$ and every open subgroup $H\subset G_\Q$ such that $W|_{H\cap I_B}$ stabilises a full flag
with graded characters $\modularitycyc^{e}$, $\modularitycyc^{-e}$: there is a newform $f$ of weight $2e+1$ attached via
$\iota$ to $W\otimes\modularitycyc^{-e}$.
\formallabel{UpperTheory.V.Z.Win.RR.AutNewformA}
\formallabel{UpperTheory.V.Z.Win.RR.WitUp}
\formal{UpperTheory.V.Af.IsAutR}
\end{definition}

\begin{theorem}\label{pf:modularity:autnew}
$\sfname{HypothesesH'}$ implies $\sfname{AutNewform}$.
\formallabel{UpperTheory.V.Br.autNewformA_of_printsH}
\end{theorem}

\begin{proof}
Consider the following statements, for $k\ge3$ and cuspidal automorphic representations $\pi$ of $\GL_2(\A_\Q)$.
\begin{itemize}
\item $\sfname{HermEx}$: for every totally real field $F$ and $a\ge1$, every cuspidal automorphic representation of
$\GL_a(\A_F)$ annihilated by the central element of the Lie algebra used in \cite{LeanSource2026} carries an invariant
Hermitian form.
\item $\sfname{RepEig}$: if $\pi$ has infinitesimal character $\{\frac{k-1}2,-\frac{k-1}2\}$, there are $N\ge1$, a
Dirichlet character $\chi$ modulo $N$ and an eigenform $g$ of weight $k$, level $N$ and character $\chi$ such that for almost
all primes $p$, $\pi_p$ has Satake parameters $\{\alpha,\beta\}$ with $a_p(g)=p^{(k-1)/2}(\alpha+\beta)$ and
$\alpha\beta=\chi(p)$.
\item $\sfname{RepNew}$: if $\pi$ has infinitesimal character $\{\frac12+j,\frac32-k+j\}$ for some $j\in\Z$, there is a
newform $f$ of weight $k$ such that for almost all primes $p$ the Hecke polynomial of $f$ at $p$ is
$\prod_b(X-p^jb)$, the product over the renormalised Satake parameters $b$ of $\pi_p$ as in \cref{pf:modularity:aut}(a).
\end{itemize}
$\sfname{HermEx}$ holds (\lref{UpperTheory.V.Br.hermEx}); $\sfname{HypothesesH'}$ and $\sfname{HermEx}$ imply
$\sfname{RepEig}$ (\lref{UpperTheory.V.Br.repEig_of_printsH_of_hermEx}); $\sfname{HypothesesH'}$ and $\sfname{RepEig}$
imply $\sfname{RepNew}$ (\lref{UpperTheory.V.Br.repNewAE_of_printsH_of_repEig}); and $\sfname{HypothesesH'}$ and
$\sfname{RepNew}$ imply $\sfname{AutNewform}$ (\lref{UpperTheory.V.Z.Win.RR.autNewformA_of_printsH_of_repNewAE}).
\end{proof}

\subsection{Modularity of universal ordinary determinant rings}\label{sec:pf:modularity:modU}

This subsection proves the remaining hypotheses of \cref{pf:modularity:linked} and \cref{pf:modularity:adjin}: the
modularity statements $\sfname{ModU}$ and $\sfname{ModN}$ for universal ordinary determinant rings, the multiplicity
statement $\sfname{Mult}$, the vanishing of ordinary adjoint Selmer groups at automorphic points
($\sfname{DenseThresholdMainCase}$), the tangent-space statement $\sfname{TanA}$, and the general theorem behind
$\sfname{LinkedUnl}$. The proofs assemble long results of the formal development, which we state.

\emph{Setting.} We use the notation of \cref{sec:pf:vwin:patch}: an odd prime $\ell$, a CM field $L$ with maximal totally
real subfield $L^+$, a Hida setting with set of places $S$, a level with $S'=S$ (so that $\Gamma=G_{L^+,S}$ and
$G_0=\ker\eta$), a coefficient ring $\OO$, $n\ge1$, a residual datum $r$ of dimension $n$, and the universal ordinary
determinant ring $B$ of $r$, with universal local characters $\psi^{\mathrm u}_{v,i}$ at the places $v\mid\ell$; for the
auxiliary places $u$ of the level, with inertia subgroups $I_u\subset G_0$, let $I'_u=I_u\cap\bigcap_i\ker\bar\chi_i$, where
$\bar\chi_i$ are the residual characters of $r$. The
coefficient ring is \emph{large} if it contains the images of the integers of $L$ under all embeddings $L\to\Qlbar$. A
\emph{point} of $B$ is an $\OO$-algebra map $x:B\to\Qlbar$; it has a determinant $D_x$ of $G_0$ and local characters
$\psi_{x,v,i}=x\circ\psi^{\mathrm u}_{v,i}$, and a representation $\rho$ of $G_L$ \emph{realises} $x$ if $D_x$, pulled back
to $G_L$, is the determinant of $\rho$. A \emph{weight} $h$ assigns integers $h_{v,i,\tau}$ to the places $v\mid\ell$,
the indices $i$ and the embeddings $\tau$ of $L_v$ in $\Qlbar$; it is \emph{increasing} if $i\mapsto h_{v,i,\tau}$ is
strictly increasing for all $v,\tau$. A point $x$ is \emph{arithmetic of weight $h$} if $h$ is increasing and each
$\psi_{x,v,i}$ is locally algebraic of weight $h_{v,i}$ (on inertia it agrees, up to a character with open kernel, with
$\sigma\mapsto\prod_\tau\tau(\mathrm{a}(\sigma))^{-h_{v,i,\tau}}$ for every local Artin map $\mathrm a$); it is \emph{good
of weight $h$} if moreover these characters with open kernel can be chosen pairwise distinct for each $v$; and it is
\emph{irreducible} if it is realised by a continuous irreducible representation. A \emph{weight algebra} $\Lambda$ for $r$
is a complete local $\OO$-algebra with characters of the inertia groups $I_v$ ($v\mid\ell$) lifting those of $r$, universal
for such characters; it maps to $B$. For an ideal $J\subset B$ we write $V(J)$ for the set of points $x$ with $x(J)=0$.

\begin{definition}[Modular ideals]\label{pf:modularity:moddef}
In this setting let $\iota:\Qlbar\cong\C$, $\Lambda$ a weight algebra and $J\subset B$ an ideal. $B$ is \emph{modular for
$(\iota,J,\Lambda)$} if:
\begin{enumerate}[label=(Mod\arabic*)]
\item (ring properties) $B/J$ is finite over $\Lambda$ and $\Lambda$-torsion free; its determinant is polarised, ordinary
at the places $v\mid\ell$ with respect to the characters $\psi^{\mathrm u}_{v,i}$ modulo $J$, trivial on the subgroups
$I'_u$, given on the inertia groups $I_u$ by the Teichm\"uller types of $r$, and, if $r$ satisfies the residual condition
$\sfname{V0ok}$, trivial on the inertia group $I_0$ at the distinguished place $v_0$;
\item (density) for every increasing weight $h$, every element of $B$ vanishing at all good points of weight $h$ in
$V(J)$ lies in $J$;
\item (point theorem) if $2\le n\le4$ and $x$ is an arithmetic point realised by a continuous $\rho$ which is conjugate
self-dual automorphic via $\iota$, trivial on the inertia group at $v_0$, and such that for every $v\mid\ell$,
$\rho|_{D_v}$ has a full flag with graded characters $\psi_{x,v,i}$ satisfying
$\psi_{x,v,i}\ne\psi_{x,v,j}\modularitycyc$ for all $i,j$, then $x\in V(J)$;
\item (decomposition) every arithmetic $x\in V(J)$ is realised by a continuous semisimple $\rho$ whose characteristic
polynomials are those of $\bigoplus_i\bigoplus_{0\le j<m_i}\rho_i\otimes\modularitycyc^{-j}$ for some continuous $\rho_i$
of dimensions $d_i\ge1$ and multiplicities $m_i\ge1$; and if all $m_i=1$, then at almost all places $w$ of $L$ the
Frobenius eigenvalues $z,z'$ of $\rho$ satisfy $|\iota(z)|<q_w|\iota(z')|$;
\item (automorphy away from $v_0$) for every arithmetic $x\in V(J)$, every irreducible $\rho$ realising $x$ is attached via
$\iota$ to a regular algebraic conjugate self-dual cuspidal representation, and, if continuous, satisfies the variant of
the window condition away from $v_0$;
\item (twists) if $x\in V(J)$ is good of weight $h$, $\beta$ is a character of $G_0$ of $\ell$-power order with open kernel,
$\beta(cgc^{-1})=\beta(g)^{-1}$, trivial on the groups $I_u$ and on $I_0$, and $x'$ is a point with
$D_{x'}=D_x\otimes\beta$ and $\psi_{x',v,j}=\psi_{x,v,j}\beta$, then $x'\in V(J)$;
\item (patching conclusions) the conclusions (a) and (b) of \cref{pf:vwin:pfromaut} hold for $J$.
\end{enumerate}
$\sfname{ModU}$ is the statement: for every Hida setting over a CM field $L$ not containing a primitive $\ell$-th root of
unity, with level $S'=S$, every large $\OO$, $n\ge1$, residual datum $r$, universal ring $B$, weight algebra $\Lambda$ and
$\iota$, there is an ideal $J$ such that $B$ is modular for $(\iota,J,\Lambda)$. $\sfname{HidaU}$ is the same statement with
(Mod7) replaced by: $J$ is automorphic at all Taylor--Wiles levels (\cref{pf:vwin:pfromaut}).
\formallabel{UpperTheory.V.W.ModUAlt}
\formallabel{UpperTheory.V.W.HidaUAlt}
\formallabel{UpperTheory.V.W.UniversalOrdinaryDeterminantRing.Pt.Arith}
\formallabel{UpperTheory.V.W.UniversalOrdinaryDeterminantRing.Pt.Good}
\formal{UpperTheory.V.W.UniversalOrdinaryDeterminantRing.ModularAlt}
\formal{UpperTheory.V.W.UniversalOrdinaryDeterminantRing.HidaPropsAlt}
\formal{UpperTheory.V.W.LamObj}
\end{definition}

\begin{lemma}\label{pf:modularity:modUof}
Assume the conclusion of the patching theorem \cref{pf:vwin:pfromaut} for all data, $\sfname{CohP}$, local class field
theory at all places, and $\sfname{HidaU}$. Then $\sfname{ModU}$ holds. In particular \ref{pf:vwin-he:H1}, the
patching theorem and $\sfname{HidaU}$ imply $\sfname{ModU}$.
\formallabel{UpperTheory.V.W.modU_of_printsH}
\formallabel{UpperTheory.V.W.modularity_unitary_of_automorphy_hida_unitary}
\end{lemma}

\begin{proof}
Given the data of $\sfname{ModU}$, $\sfname{HidaU}$ gives $J$ satisfying (Mod1)--(Mod6) and automorphic at all Taylor--Wiles
levels. The hypotheses of the patching theorem hold: its Galois cohomology statements are part of $\sfname{CohP}$; the
class field theory condition on the local degrees follows from local class field theory; and $\varepsilon\bmod\ell$ is
non-trivial on $G_0$ because $L$ contains no primitive $\ell$-th root of unity (\lref{UpperTheory.V.W.muL}). Hence
(Mod7) holds. For the second statement, $\sfname{CohP}$ and local class field theory are consequences of \ref{pf:vwin-he:H1}
(\cref{pf:modularity:loc}(c),(e)).
\end{proof}

\begin{proposition}\label{pf:modularity:hidaU}
$\sfname{HypothesesH'}$ implies $\sfname{HidaU}$.
\formallabel{UpperTheory.V.Hi.hida_unitary_of_printed_hida}
\end{proposition}

\begin{proof}
Let the data of $\sfname{HidaU}$ be given. The ideal $J$ is a Hecke ideal for ordinary automorphic forms on a definite
unitary group; all ingredients are results of the formal development proved from $\sfname{HypothesesH'}$.
\begin{enumerate}[label=(\arabic*)]
\item There are \emph{unitary Hecke data} $R$ for the setting, $\iota$ and $\Lambda$, with the empty set of Taylor--Wiles
primes (\lref{UpperTheory.V.Hi.nonempty_unitaryHeckeData}): an algebraic weight type $\tau_0$ with values in some
$\GL_m(\OO)$, a weight character $\theta$ with values in $\Lambda^\times$, and eigenvalue data for the Hecke operators,
subject to type, weight and residual compatibilities with $r$; they come with a level $K$.
\item There is an ideal $J\subset B$ which is the \emph{Hecke ideal} of $R$ at level $K$
(\lref{UpperTheory.V.Hi.exists_hecke_ideal}): either $J=B$, or $J$ is the kernel of a surjection from $B$ onto the
Hecke algebra of $R$ at level $K$, compatible with the universal determinant and with $\Lambda$.
\item For this $J$, (Mod1) holds by \lref{UpperTheory.V.Hi.ring_props_of_is_j_unitary_hecke_data}, (Mod2) by
\lref{UpperTheory.V.Hi.dense_of_is_j_unitary_hecke_data}, (Mod3) by
\lref{UpperTheory.V.Hi.hecke_ideal_point_properties}, (Mod4) and (Mod5) by
\lref{UpperTheory.Gal.projectedFormsSystemPart.decomposition_and_automorphicAway_of_isJ}, and (Mod6) by
\lref{UpperTheory.V.Z.Desc.W.twistOk_of_isJ}.
\item $J$ satisfies the level statement of the patching argument
(\lref{UpperTheory.V.Hi.level_of_is_j_unitary_hecke_data}), which, together with the local theorems (a consequence of
\ref{pf:vwin-he:H1}), implies that $J$ is automorphic at all Taylor--Wiles levels
(\lref{UpperTheory.V.Pa.has_automorphy_of_statement}).\qedhere
\end{enumerate}
\end{proof}

\begin{theorem}\label{pf:modularity:modU}
$\sfname{HypothesesH'}$ implies $\sfname{ModU}$.
\formallabel{UpperTheory.V.HV.Core.modU_of_printsH}
\end{theorem}

\begin{proof}
Apply \cref{pf:modularity:modUof} with the component \ref{pf:vwin-he:H1} of $\sfname{HypothesesH'}$, the patching theorem
\cref{pf:vwin:pfromaut}, and $\sfname{HidaU}$, which holds by \cref{pf:modularity:hidaU}.
\end{proof}

\begin{definition}[Multiplicities]\label{pf:modularity:modNdef}
In the setting above let $\Lambda$ be a weight algebra and $J\subset B$ an ideal. \emph{Hecke data} for $(J,\Lambda)$ is
a $B$-module $M$ on which $\Lambda$ acts through $\Lambda\to B$, free and finite over $\Lambda$, such that $J$ is radical,
a prime of $B$ contains the annihilator of $M$ if and only if it contains $J$, and $B/J$ is $\Lambda$-torsion free. For a
$\Lambda$-linear idempotent $e$ of $M$ commuting with $B$ and a point $x$, let $\mu_e(x)$ be the dimension over $\Qlbar$ of
the intersection of the image of $e$ with the generalised eigenspaces of all $b\in B$, for the eigenvalues $x(b)$, in
$M\otimes_{\Lambda,x}\Qlbar$. $\sfname{Num}(B,\Lambda,J)$ is the statement: there are Hecke data $M$ for $(J,\Lambda)$ and
an idempotent $e$ as above such that $1\le\mu_e(x)\le\mu_e(y)$ for all irreducible points $x,y\in V(J)$ with $x$
arithmetic of some weight $h$ and $y$ good of weight $h$. $\sfname{ModN}$ is the statement: for every Hida setting over a
CM field $L$ not containing a primitive $\ell$-th root of unity, with level $S'=S$, every large $\OO$, every residual datum
$r$ of dimension $3$ satisfying $\sfname{V0ok}$, every $B$, $\Lambda$ and $\iota$, there is $J$ such that $B$ is modular
for $(\iota,J,\Lambda)$ and $\sfname{Num}(B,\Lambda,J)$ holds. $\sfname{NumJ}$ is the statement: $\sfname{HypothesesH'}$,
$\sfname{Mult}$ and $\sfname{CompOf}$ (\cref{pf:modularity:multdef}) imply that $\sfname{Num}(B,\Lambda,J)$ holds for every
Hida setting as above, $1\le n\le4$, unitary Hecke data $R$ (for any set of Taylor--Wiles primes), level $K$ of $R$, and
ideal $J$ which is the Hecke ideal of $R$ at level $K$ and is radical, provided $r$ satisfies $\sfname{V0ok}$ and $\Lambda$
is a domain.
\formallabel{UpperTheory.V.J.Num}
\formallabel{UpperTheory.V.J.ModNAlt}
\formallabel{UpperTheory.V.Z.J.D.NumJAlt}
\formal{UpperTheory.V.J.Dat}
\formal{UpperTheory.V.Z.Desc.V0.muE}
\end{definition}

\begin{theorem}\label{pf:modularity:modNnum}
$\sfname{NumJ}$, $\sfname{HypothesesH'}$, $\sfname{Mult}$ and $\sfname{CompOf}$ imply $\sfname{ModN}$.
\formallabel{UpperTheory.V.Z.J.D.modularity_n_of_num_j_printed_mult_comp}
\end{theorem}

\begin{proof}
Let the data of $\sfname{ModN}$ be given ($n=3$). As in the proof of \cref{pf:modularity:hidaU}, there are unitary Hecke
data $R$ with no Taylor--Wiles primes, a level $K$ and the Hecke ideal $J$ of $R$ at level $K$, and $J$ satisfies
(Mod1)--(Mod6) and is automorphic at all Taylor--Wiles levels. As in the proof of \cref{pf:modularity:modUof}, the hypotheses of
the patching theorem hold ($\sfname{CohP}$ and local class field theory follow from \ref{pf:vwin-he:H1}, and
$\varepsilon\bmod\ell$ is non-trivial on $G_0$ since $L$ contains no primitive $\ell$-th root of unity), so (Mod7) holds and
$B$ is modular for $(\iota,J,\Lambda)$. For $\sfname{Num}$: $\Lambda$ is a domain
(\lref{UpperTheory.V.Hi.isDomain_and_denseWt_of_isTheta}, from \ref{pf:vwin-he:H1} and the properties of the weight
character of $R$); and $J$ is radical: if $b^k\in J$, then $x(b)=0$ for every $x\in V(J)$, so $b\in J$ by (Mod2) applied to the
increasing weight $h_{v,i,\tau}=i$. Hence $\sfname{NumJ}$ applies (with $n=3\le4$) and gives $\sfname{Num}(B,\Lambda,J)$.
\end{proof}

\begin{theorem}\label{pf:modularity:modN}
$\sfname{HypothesesH'}$, $\sfname{Mult}$ and $\sfname{CompOf}$ imply $\sfname{ModN}$. Since $\sfname{Mult}$
(\cref{pf:modularity:mult}) and $\sfname{CompOf}$ (\lref{UpperTheory.V.Z.Mult.compOf_free}) hold,
$\sfname{HypothesesH'}$ implies $\sfname{ModN}$.
\formallabel{UpperTheory.V.Z.J.D.modularity_n_of_printed_mult_comp}
\end{theorem}

\begin{proof}
By \cref{pf:modularity:modNnum} it suffices that $\sfname{NumJ}$ holds; this is \lref{UpperTheory.V.Z.J.D.numJ'_alt}.
In its formal proof, $M$ is the Hecke module of $R$ at level $K$ (which is Hecke data for $(J,\Lambda)$), $e$ is an idempotent
commuting with $B$ constructed from data provided by $\sfname{Mult}$ and $\sfname{CompOf}$, and it is shown that
$\mu_e(x)=\mu_e(y)\ge1$, by computing both multiplicities in terms of these data.
\end{proof}

\begin{definition}[Multiplicity statements]\label{pf:modularity:multdef}
Let $E$ be a CM field such that $E/E^+$ is unramified at all finite places, and $1\le n\le4$. We use the definite unitary
group in $n$ variables attached to $E/E^+$ of \cite{LeanSource2026}: for a dominant weight $\lambda$, a set $S$ of places of
$E$ containing exactly one place above each place of $E^+$ split in $E$ (a \emph{section}), and a level $K'=(K'_w)_w$ (open
compact subgroups, equal to the standard maximal compact subgroup at the non-split places) we have the finite-dimensional
space $\mathcal A(S,\lambda,K')$ of automorphic forms of weight $\lambda$ and level $K'$. An automorphic representation
$\sigma$ of the unitary group is \emph{attached to} a cuspidal $\pi$ of $\GL_n(\A_E)$ if at almost all split places $w$, the
local component of $\sigma$ at $w$ (a representation of $\GL_n(E_w)$) and $\pi_w$ have the same Satake parameters.
\begin{enumerate}[label=(\alph*)]
\item $\sfname{Mult}$: for all such $E$, $n$ and $\lambda$ there is $a\ge1$ such that for every regular algebraic conjugate
self-dual cuspidal $\pi$ of $\GL_n(\A_E)$ with infinitesimal character given by $\lambda$ and with Satake parameters at
all non-split places, every section $S$, every level $K'$ such that the local components of $\pi$ at the split places of $S$
have non-zero $K'_w$-fixed vectors, and every decomposition of $\mathcal A(S,\lambda,K')$ as the direct sum of the
$K'$-invariants of finitely many automorphic representations $\sigma_i$ of weight $\lambda$, exactly $a$ of the $\sigma_i$
are attached to $\pi$.
\item $\sfname{CompOf}$: for such $E$, $n$ and $\pi$, every automorphic representation $\sigma$ of the unitary group which is
unramified at the non-split places and attached to $\pi$ has, at every split place $w$, local component $\pi_w$.
\end{enumerate}
\formallabel{UpperTheory.V.Z.Mult.Mult}
\formallabel{UpperTheory.V.Z.Mult.CompOf}
\end{definition}

\begin{theorem}\label{pf:modularity:mult}
$\sfname{Mult}$ holds.
\formallabel{UpperTheory.V.Z.Mult.mult_free}
\end{theorem}

\begin{proof}
The proof compares traces of Hecke operators on the spaces $\mathcal A(S,\lambda,K')$. For a finite set $Z$ of split places
in $S$, Hecke functions $f_w$ at $w\in Z$, a further finite set $T$ of split places disjoint from $Z$, and auxiliary data
$e_1,e_2,k$ at $T$, the formal development defines a test function and the trace of the corresponding Hecke operator on
$\mathcal A(S,\lambda,K')$; $\mathrm{val}_T$ denotes the explicit function of Satake parameters at $T$ entering it. The
proof is the combination of three statements.
\begin{enumerate}[label=(\roman*)]
\item ($\sfname{CountAt}$, by \lref{UpperTheory.V.Z.Mult.rCountAt} from $\sfname{CompOf}$,
\lref{UpperTheory.V.Z.Mult.evalD}, \lref{UpperTheory.V.Z.Mult.cpUniq} and
\lref{UpperTheory.V.Z.Tf.GlobalTraceComparison.Asm.places}.) For $\pi$ as in $\sfname{Mult}$, a section $S$, a finite
set $Z$ of split places, Satake parameters $B_w$ of $\pi$ off $Z$, local components $\tau_w$ of $\pi$ at $w\in Z$, a level $K'$
equal to the maximal one off $Z$, and a decomposition $(\sigma_i)$, there are finitely many families $C_j$ of Satake
parameters, each differing from $(B_w)$ at some split place off $Z$, such that for Hecke functions $f_w$ ($w\in Z$) and
auxiliary data with $\mathrm{val}_T(B)=1$ and $\mathrm{val}_T(C_j)=0$, the trace of the $k$-th test function is
$\dim W_\lambda\cdot\#\{i:\sigma_i\text{ attached to }\pi\}\cdot\prod_{w\in Z}\tr\tau_w(f_w)$ for $k\ge1$, where
$W_\lambda$ is the representation of weight $\lambda$. Here \lref{UpperTheory.V.Z.Mult.evalD} is the spectral expansion
of the trace as $\dim W_\lambda\sum_i\prod_{w\in Z}\tr t_{i,w}(f_w)\,\mathrm{val}_T(B_i)^k$ in terms of local components
$t_{i,w}$ and Satake parameters $B_i$ of the $\sigma_i$, and \lref{UpperTheory.V.Z.Mult.cpUniq} is the uniqueness of
local components at split places.
\item ($\sfname{Ident}$, by \lref{UpperTheory.V.Z.Mult.rIdent}.) There is $a>0$, depending only on $(E,n,\lambda)$,
such that $\#\{i:\sigma_i\text{ attached to }\pi\}=a$ for all $\pi$ as in $\sfname{Mult}$ and all data as in (i) together
with an algebraic Hecke character unramified off $Z$, Hecke functions $f_w$ with $\tr\tau_w(f_w)\ne0$ ($w\in Z$), one of them
supported on elements with separable characteristic polynomial, and a level adapted to the $f_w$. It is deduced from (i)
and the following results: the geometric expansion of the trace for such test functions as $\dim W_\lambda$ times an
explicit geometric expression, with $\dim W_\lambda\ne0$ (\lref{UpperTheory.V.Z.Mult.geomTr}); a stabilisation
statement providing $c_0>0$, depending only on $(E,n,\lambda)$, which controls the geometric expression in terms of $\pi$
and finitely many other automorphic data of $\GL_n$ (\lref{UpperTheory.V.Z.Mult.stableBlockSA}); the statements of
the trace formula comparison used as the hypotheses $\sfname{G2s}$, $\sfname{MatchUnit}$, $\sfname{MatchSplit}$ and
$\sfname{EndoBlockS}$ of \lref{UpperTheory.V.Z.Mult.rIdent} (\lref{UpperTheory.ConvolutionTheory.G2.g2s},
\lref{UpperTheory.ConvolutionTheory.G2.matchUnit} and \lref{UpperTheory.ConvolutionTheory.G2.matchSplit},
\lref{UpperTheory.V.Z.Tf.GaloisFromTraceFormula.endoS}); and the fact that a regular
automorphic datum of $\GL_n$ not containing $\pi$, with Satake parameters stable under inversion at almost all non-split
places, cannot have the Satake parameters of $\pi$ at almost all split places
(\lref{UpperTheory.V.Z.Mult.noData_free}).
\item ($\sfname{Mult}$, by \lref{UpperTheory.V.Z.Mult.rMult}.) $\sfname{Mult}$ follows from (i), (ii) and the following
auxiliary statements: given $\pi$, $S$ and $K'$, there are a finite set $Z\subset S$ of split places off which $K'$ is
maximal, Satake parameters, local components, and a Hecke function with non-zero trace supported on elements with separable
characteristic polynomial (\lref{UpperTheory.V.Z.Mult.choiceM}); an algebraic Hecke character unramified off a given
finite set exists (\lref{UpperTheory.V.Z.Tf.GlobalTraceComparison.muEx}); a smaller level at a split place carries a
Hecke function with non-zero trace on every representation having fixed vectors, and every locally constant function is
bi-invariant under some level (\lref{UpperTheory.V.Z.Mult.levInd} and \lref{UpperTheory.V.Z.Mult.lcLev}); the trace does not
depend on the level (\lref{UpperTheory.V.Z.Mult.levTr}); the geometric expansion of (ii); and the global level only
depends on the components in $S$ (\lref{UpperTheory.V.Z.Mult.offSec}). With these, the auxiliary data of (ii) can be
chosen for every $\pi$, $S$, $K'$ and $(\sigma_i)$ as in $\sfname{Mult}$, which gives the claim.\qedhere
\end{enumerate}
\end{proof}

\begin{definition}[Automorphic points and ordinary cocycles]\label{pf:modularity:coredef}
In the setting above let $x$ be a point of $B$ and $\rho:G_0\to\GL_n(\Qlbar)$; let $G'\subset G_0$ be the subgroup of the
patching theorem (\cref{pf:vwin:pfromaut}, (P2)), on which the cyclotomic character is trivial.
\begin{enumerate}[label=(\alph*)]
\item $(x,\rho)$ is an \emph{automorphic point} via $\iota$ if $\rho$ is continuous and realises $x$ (through $G_L\to
G_0$), $x$ is arithmetic of some weight, $\rho|_{G_L}$ is conjugate self-dual automorphic via $\iota$, $\rho$ is trivial on
the inertia group $I_0$, and for every $v\mid\ell$, $\rho|_{G_v}$ has a full flag with graded characters $\psi_{x,v,i}$. It
is \emph{locally generic} if moreover: for every $v\mid\ell$ there is $g\in G_v$ mapping into $G'$ such that
$\psi_{x,v,i}(g)^m\ne\psi_{x,v,j}(g)^m$ for all $i\ne j$ and $m\ge1$; $\rho|_{G'}$ is irreducible; and every open subset
of $G_0$ meeting $G'$ contains an element of $G'$ at which $\rho$ has separable characteristic polynomial.
\item Let $P$ be a polarisation of $\rho$ and $R$ an extension of the adjoint representation of $\rho$ to $\Gamma$
compatible with $P$ (\lref{UpperTheory.V.Ws.Eb.IsAdExt}). A $1$-cocycle $z$ of $\Gamma$ with values in $R$ is
\emph{ordinary} if it vanishes on the subgroups $I'_u$ and is ordinary at each $v\mid\ell$ for the
characters $\psi_{x,v,i}$, in the sense of \cite{LeanSource2026} (in a basis adapted to the flag, it is upper triangular
on $G_v$ and has zero diagonal on $I_v$, up to a coboundary).
\item $\sfname{DenseThresholdMainCase}$ is the statement: for every Hida setting over a CM field $L$ not containing a
primitive $\ell$-th root of unity, with $S'=S$, every large $\OO$, $2\le n\le4$, residual datum $r$ satisfying
$\sfname{V0ok}$, $B$, $\iota$, and every automorphic, locally generic point $(x,\rho)$, every ordinary $1$-cocycle (for every
$P$ and $R$ as in (b)) is a coboundary.
\item $\sfname{OrdVan2}$ is the statement: for every Hida setting, $\OO$, $2\le n\le4$ and $r$ as in (c), every arithmetic
point $x$ and continuous $\rho:G_0\to\GL_n(\Qlbar)$ with $D_x=\det\rho$, such that $\rho|_{G_L}$ is conjugate self-dual
automorphic via $\iota$, the image of $\rho$ spans $M_n(\Qlbar)$, $\rho$ is polarised by $P$, $\rho$ is trivial on $I_0$,
$\psi_{x,v,i}\ne\psi_{x,v,j}\modularitycyc$ for all $v\mid\ell$ and all $i,j$, and $\rho$ satisfies the regularity condition
$\sfname{RegDet}$ (there are finitely many elements of $G'$ at which $\rho$ has separable characteristic polynomial and
which satisfy a determinant condition), every continuous $1$-cocycle $z$ of $G_0$ with values in $M_n(\Qlbar)$ with the
adjoint action, which is polarised ($Pz(cgc^{-1})=-{}^tz(g)P$), vanishes on the subgroups $I'_u$ and is
ordinary at each $v\mid\ell$, has the form $z(g)=\rho(g)Y\rho(g)^{-1}-Y$.
\end{enumerate}
\formallabel{UpperTheory.V.Th.AutomorphicPointGalois}
\formallabel{UpperTheory.V.Th.AutomorphicPointLocallyGeneric}
\formallabel{UpperTheory.V.Th.DenseThresholdMainCase}
\formallabel{UpperTheory.V.Th.OrdVan2Alt}
\formal{UpperTheory.V.Th.OrdZero}
\formal{UpperTheory.V.Th.OrdCoc}
\formal{UpperTheory.V.Th.IsOrdCoc}
\end{definition}

\begin{lemma}\label{pf:modularity:coreof}
$\sfname{OrdVan2}$ and $\sfname{PhiAll}$ (for every Hida setting and level, every open subgroup of $\Gamma$, of $G_0$ and of
each $G_v$ has only finitely many open homomorphisms to $\Z/\ell$) imply $\sfname{DenseThresholdMainCase}$.
\formallabel{UpperTheory.V.Th.core_of_ord_vanishing_two_phi_all}
\end{lemma}

\begin{proof}
Let $(x,\rho)$ be automorphic and locally generic, $P$, $R$ as in \cref{pf:modularity:coredef}(b) and $z$ an ordinary
cocycle; let $h$ be a weight of $x$. By \lref{UpperTheory.V.Th.Det.redB}, $D_x=\det\rho$. As $\rho|_{G'}$ is
irreducible, so is $\rho$, and by Burnside's theorem the image of $\rho$ spans $M_n(\Qlbar)$
(\lref{KStar.Schema.span_eq_top_of_irreducible}); therefore the subgroups $I'_u$ act trivially through $\rho$
(\lref{UpperTheory.V.Ws.Eb.relKill}). By \lref{UpperTheory.V.Th.Det.polNorm} there is $X_0$ such that
$z'=z-dX_0$, where $dX(g)=\rho(g)X\rho(g)^{-1}-X$, is polarised on $G_0$, and such that $z$ is a coboundary on $\Gamma$ as
soon as $z'|_{G_0}=dY$ for some $Y$. The condition $\sfname{V0ok}$ gives $g\in G_0$ with $\varepsilon(g)\ne1$, so
\lref{UpperTheory.V.Th.regDet} (from $\sfname{PhiAll}$, continuity, irreducibility on $G'$ and the separability
condition) gives $\sfname{RegDet}$. For $v\mid\ell$ and indices $i,j$ we have $\psi_{x,v,i}\ne\psi_{x,v,j}\modularitycyc$: if $i=j$, because
$\modularitycyc$ is non-trivial on the decomposition group at $v$ (\lref{UpperTheory.V.Th.Det.epsDec}); if $i\ne j$,
local genericity gives $g\in G_v$ mapping into $G'$, where $\modularitycyc$ is trivial, with
$\psi_{x,v,i}(g)\ne\psi_{x,v,j}(g)$. Finally $z'|_{G_0}$ is continuous, a cocycle for the adjoint action, polarised, vanishes
on the subgroups $I'_u$ (as $z$ does and $\rho$ is trivial there), and is ordinary at each $v\mid\ell$
(ordinarity is insensitive to coboundaries). By $\sfname{OrdVan2}$, $z'|_{G_0}=dY$ for some $Y$; hence $z$ is a coboundary.
\end{proof}

\begin{theorem}\label{pf:modularity:core}
$\sfname{HypothesesH'}$ and $\sfname{CohP}$ imply $\sfname{DenseThresholdMainCase}$.
\formallabel{UpperTheory.V.Th.core_alt_bis}
\end{theorem}

\begin{proof}
$\sfname{CohP}$ implies $\sfname{PhiAll}$ (\lref{UpperTheory.V.W.phiAll}), and $\sfname{HypothesesH'}$ implies
$\sfname{OrdVan2}$ (\lref{UpperTheory.V.Th.ordinary_vanishing_two_of_printed_hida}); apply \cref{pf:modularity:coreof}.
The second implication carries the weight of the argument. In its formal proof, after an extension of the coefficient ring
and a change of residual datum preserving the local characters, one considers the Hecke ideal $J$ given by the Hida family
statement proved from $\sfname{HypothesesH'}$ (\lref{UpperTheory.V.Th.hida_x_of_printed_hida}); the point $x$ lies on
$V(J)$ by the point theorem for $J$, which applies because $\rho$ is conjugate self-dual automorphic, unramified at $I_0$,
and has unlinked flags at the places above $\ell$; and the vanishing then follows from the ordinary main case
(\lref{UpperTheory.V.Th.core_o_of_literal_automorphy_n}, proved from \ref{pf:vwin-he:H1}) applied to the
corresponding point of $V(J)$.
\end{proof}

We now turn to the frames and links of \cref{pf:modularity:link}. A \emph{realisation} of a Galois link $k$ (from
$(\eta,\Psi,c,g_1,B,M)$ to $\Phi$ over $L$) with coefficients $\OO_1$, in dimension $3$ and with twist
$\modularitycyc^{-1}$, is a structure defined in \cite{LeanSource2026} which contains in particular: a Hida setting over $L$
and the universal ordinary determinant ring $B_0$ of a three-dimensional residual datum for it; a universal deformation ring
$S$ of the residual representation of rank three of the frame $\Phi$ (extended to the residue field of $\OO_1$), with
universal point $\rho^{\mathrm u}$ and universal local characters $\psi^{\mathrm u}_{v,i}$ ($v\in V_\Phi$); and the subring
$S_B\subset S$ generated by the coefficients of the determinant.

\begin{definition}\label{pf:modularity:tanAdef}
$\sfname{TanA}$ is the statement: for every Galois link $k$ from $(\eta,\Psi_{e,\psi},c,g_1,B,M)$ to a frame $\Phi$ over $L$
for which the link hypotheses hold, every coefficient ring $\OO_1$ with a map $k_\OO\to k_{\OO_1}$, every realisation of $k$
with coefficients $\OO_1$ in dimension $3$ and twist $\modularitycyc^{-1}$, every $\OO_1$-point $\xi$ of $S$ such that the
representation $r:G_L\to\GL_3(\Qlbar)$, $r(s)=\xi(\rho^{\mathrm u}(p(s)))$, is irreducible and
$\xi(\psi^{\mathrm u}_{v,i}(g))=\Psi(j_v(g))^{1-i}$ for $g\in I_v$, every $v\in V_\Phi$ at which $r$ is potentially
crystalline, and all $i,j$ with $\xi(\psi^{\mathrm u}_{v,i}(g))=\xi(\psi^{\mathrm u}_{v,j}(g))\modularitycyc(j_v(g))$ for
$g\in G_v$: there are two $\OO_1$-linearly independent point derivations $\delta,\delta':S\to\OO_1$ at $\xi$ vanishing on all
$\psi^{\mathrm u}_{v',a}(g)$ with $g\in I_{v'}$.
\formallabel{UpperTheory.V.J.TanA}
\formal{UpperTheory.V.J.CReal}
\end{definition}

\begin{theorem}\label{pf:modularity:tanA}
The global compatibility statement of \cref{pf:modularity:loc}(d) and the local theorems imply $\sfname{TanA}$.
\formallabel{UpperTheory.V.J.tanA}
\formallabel{UpperTheory.V.J.tanA_of}
\end{theorem}

\begin{proof}
Let the data of $\sfname{TanA}$ be given.
\begin{enumerate}[label=(\arabic*)]
\item (\lref{UpperTheory.V.Z.J.unitOf}, from the local theorems.) There is $i'\in\{0,1\}$ such that the specialised
point $\xi\circ\rho^{\mathrm u}$ has a \emph{unit at $(v,i')$}: after conjugation by some $g\in\GL_3(\OO_1)$ its restriction
to $G_v$ is upper triangular with diagonal $(\xi\circ\psi^{\mathrm u}_{v,a})_a$, and its entries at the step $(i',i'+1)$
satisfy a unit condition on $I_v$ relative to the cyclotomic character \cite{LeanSource2026}.
\item (\lref{UpperTheory.V.Pr.degEq}.) The degree of $\Phi$ equals $[L^+:\Q]$, and $L$ contains no primitive $\ell$-th
root of unity.
\item (\lref{UpperTheory.V.Z.J.linkOf}, using (2).) The extended frame is realised by a Galois-theoretic frame over
$L^+$: there are elements $c_w\in\Gamma_\Phi$, for the infinite places $w$ of $L^+$, with $c_w^2=1$ and $\delta(c_w)=-1$,
and an identification of the extended frame with one attached to the Galois group of the maximal extension of $L^+$
unramified outside a finite set, compatible with the local data.
\item (\lref{UpperTheory.V.Z.J.gwlOf} and \lref{UpperTheory.V.Z.J.leOf2}, from the global compatibility statement.) The
frame of (3) satisfies the Greenberg--Wiles inequality for finite modules killed by a power of $\ell$ (an upper bound for
$\#N\prod_v\#Z_v$ in terms of the corresponding Selmer set, the $\ell$-power roots of unity, the invariants of the $c_w$ and
$\#N^{|V_\Phi|}$) and the local Euler characteristic formula at the places $v\in V_\Phi$.
\item (\lref{UpperTheory.V.Z.J.tan}.) For a frame satisfying $\sfname{Ok}(C_0)$, with elements $c_w$ as in (3), the
Greenberg--Wiles inequality and the local Euler characteristic formula, every $\OO_1$-point of the universal deformation ring
of rank three having a unit at some $(v,i')$ admits two linearly independent point derivations vanishing on the universal
inertial characters.
\end{enumerate}
Applying (5) to the frame of (3), with (1) and (4), gives $\delta$ and $\delta'$.
\end{proof}

\begin{theorem}\label{pf:modularity:linkedunl}
Global Artin reciprocity, local class field theory at all places, $\sfname{ModU}$, $\sfname{ModN}$, $\sfname{LevEx}$,
$\sfname{LamEx}$, $\sfname{ProModularRingExists}$, $\sfname{U3Dense}$, $\sfname{Seed2}$, the twist, cyclotomic and
restriction properties of algebraic characters, $\sfname{TanA}$ and $\sfname{DenseThresholdMainCase}$ imply
$\sfname{LinkedUnl}$.
\formallabel{UpperTheory.Gal.projectedFormsSystemPart.linkedUnlE_alt}
\formallabel{UpperTheory.Gal.projectedFormsSystemPart.UnitaryRankTwoFlagsAlt}
\end{theorem}

\begin{proof}
Let the data of $\sfname{LinkedUnl}$ be given: a Galois link $k$ from $(\eta,\Psi,c,g_1,B,M)$ to $\Phi$ over $L$ with
$\sfname{Ok}(C_0)$, $\sfname{Big}(C_0)$ and the seed clause, $N$ even, a discrete valuation ring $A$ finite over $\OO$, a
point $\nu$ of rank three over $A$ with inertial characters $\Psi^{1-i}\circ j_v$, the representation $r:G_L\to\GL_3(\Qlbar)$
it defines, irreducible on every open subgroup, a place $v\in V_\Phi$ at which $r$ is potentially crystalline, and indices
$i,j$. Suppose, for a contradiction, that $\nu_{v,i}(g)=\nu_{v,j}(g)\modularitycyc(j_v(g))$ for all $g\in G_v$. Fix an
isomorphism $\iota:\Qlbar\cong\C$.
\begin{enumerate}[label=(\arabic*)]
\item (\lref{UpperTheory.Gal.projectedFormsSystemPart.openEx2E_alt}, from all hypotheses except $\sfname{TanA}$ and
$\sfname{DenseThresholdMainCase}$.) There are a large coefficient ring $\OO_1\supset A$, a realisation of $k$ with
coefficients $\OO_1$ in dimension $3$ and twist $\modularitycyc^{-1}$, a weight algebra $\Lambda\cong\OO_1[[T_1,\dots,T_m]]$
for $B_0$, an ideal $J\subset B_0$ such that $B_0$ is modular for $(\iota,J,\Lambda)$ and $\sfname{Num}(B_0,\Lambda,J)$
holds, and a surjection $\varphi:B_0\to S_B$ compatible with the universal determinants and local characters, such that all
points of $S$ on the ideal generated by $\varphi(J)$ are good, and a dimension statement for $J$ and an endpoint statement
hold \cite{LeanSource2026}.
\item (\lref{UpperTheory.V.J.xiEx}.) There is an $\OO_1$-point $\xi$ of $S$ specialising the universal point and
characters to $\nu$.
\item (\lref{UpperTheory.V.J.xiTI}.) As $r$ is irreducible, no column of $\xi\circ\rho^{\mathrm u}$ vanishes off the
diagonal, and the induced point of a subring $S_A\subset S$ into which $S_B$ maps \cite{LeanSource2026} is integral and
has the same property.
\item (\lref{UpperTheory.V.J.xiArith}, from global Artin reciprocity.) The point $x=\xi\circ\varphi$ of $B_0$ is
arithmetic of some weight and irreducible.
\item ($\sfname{TanA}$.) Since $r$ is potentially crystalline at $B_v$ and $\nu_{v,i}=\nu_{v,j}\cdot(\modularitycyc\circ
j_v)$, there are two linearly independent point derivations at $\xi$ vanishing on the universal inertial characters.
\item (\lref{UpperTheory.V.J.unramified_algebraic_weights}.) Since $\Lambda$ is a power series ring over $\OO_1$ and
(Mod2), (Mod4) and $\sfname{Num}$ hold for $J$: at every arithmetic irreducible point of $V(J)$, every point derivation which
vanishes on $J$ and on the universal inertial characters is zero.
\item (\lref{UpperTheory.Gal.projectedFormsSystemPart.red2_alt}, from global Artin reciprocity and
$\sfname{DenseThresholdMainCase}$.) At the integral point induced by $\xi$ in (3), the ideal generated by $\varphi(J)$
vanishes locally: every $a\in J$ is annihilated by some element not vanishing at that point.
\item (\lref{UpperTheory.V.J.endJ}.) Under (3), (4), (6) and (7), any two point derivations at $\xi$ vanishing on the
universal inertial characters are linearly dependent.
\end{enumerate}
(5) and (8) contradict each other.
\end{proof}

\subsection{The Eisenstein congruence number for \texorpdfstring{$\Uni(3)$}{U(3)}}\label{sec:pf:modularity:numU3}

The counts for cases I and V (\lref{UpperTheory.V.Z.J.count_i_of_printed_bound_k_kato_automorphy},
\lref{UpperTheory.V.Z.J.count_v_of_printed_bound_k_reduced_unitary_lift}) use the analogue of
$\sfname{NumberSp4}$ for three-dimensional representations of $G_K$, $K$ imaginary quadratic. For a prime power $q$,
$e\in\N$ and $z\in\C^\times$ put $\mathrm{eis}_3(q,e;z)=q^{-e}z+1+q^ez^{-1}$.

\begin{definition}\label{pf:modularity:numberU3}
Let $K$ be an imaginary quadratic field, $\iota:\Qlbar\cong\C$, $e\ge1$, $N\ge1$ and $\rho:G_K\to\GL_3(\Qlbar)$.
\begin{enumerate}[label=(\alph*)]
\item $\rho$ is \emph{admissible of unitary type} (level $N$) if it is irreducible, $\rho\otimes\modularitycyc^{-1}$ is
conjugate self-dual automorphic via $\iota$, $\rho$ is of level $N$ and type $(N,(e,0,-e))$, and it satisfies the
condition $\sfname{Unl}$ of \cite{LeanSource2026} at $\ell$ for the weights $(e,0,-e)$.
\item $\rho$ is \emph{endoscopic} (level $M$) if there is an irreducible $W:G_\Q\to\GL_2(\Qlbar)$ such that
$W\otimes\modularitycyc^{-e}$ is attached via $\iota$ to a newform of weight $2e+1$, $W$ satisfies
$\sfname{WitAQ}(\iota)$, $W$ is of level $M$ and type $(M,(e,-e))$, $\det W$ is of finite order and factors through a
cyclotomic field, and $\rho$ is conjugate to the representation
$s\mapsto\left(\begin{smallmatrix}W_{11}&0&W_{12}\\0&1&0\\W_{21}&0&W_{22}\end{smallmatrix}\right)(s)$ of $G_K$.
\item $\sfname{NumberU3}$ is the statement: \emph{for every imaginary quadratic field $K$ there are $c,A\in\N$ such that
for all $e\ge1$, $N\ge1$, primes $p\ge7$ with $p\nmid Nd_K$, places $w\mid p$ of $K$ with norm $q_w$, and finite families
$(z_o)_{o\in\mathcal O}$ of roots of unity stable under $\Aut(\C)$, there is an integer $\mathcal N\ne0$ with
$\log|\mathcal N|\le c\,\#\mathcal O\,(N(e+1))^A\log p$ such that for every prime $\ell\ne2,p$, every $\iota$, every
Frobenius element $\sigma\in G_K$ at a valuation ring over $w$, and every family $(X_o)$ of finite sets of pairwise
non-isomorphic representations which are admissible of unitary type of level $N$ or endoscopic of level $N|d_K|$:
$\tr\rho(\sigma^{-1})\ne\iota^{-1}\mathrm{eis}_3(q_w,e;z_o)$ for $\rho\in X_o$, and
\[
 \|\mathcal N\|\le\prod_{o}\prod_{\rho\in X_o}\bigl\|\tr\rho(\sigma^{-1})-\iota^{-1}\mathrm{eis}_3(q_w,e;z_o)\bigr\|.
\]}
\end{enumerate}
A \emph{trace datum} for $(q,e)$ is a finite set $\mathcal S$ with complex numbers $t_s$ ($s\in\mathcal S$), stable
under $\Aut(\C)$, with $|t_s|\le\sqrt q+2$ and $q^et_s$ an algebraic integer; it \emph{classifies} a class of
representations at $g$ (via $\iota$) if there is a map $\mathrm{cl}$ from the class to $\mathcal S$, with
$\mathrm{cl}(\rho)=\mathrm{cl}(\rho')$ only if $\rho\cong\rho'$, such that $\tr\rho(g)=\iota^{-1}t_{\mathrm{cl}(\rho)}$.
\formallabel{UpperTheory.V.Z.J.NumberU3U_s}
\formallabel{UpperTheory.V.Z.J.Adm3U_s}
\formallabel{UpperTheory.V.Af.Adm3Alt}
\formallabel{UpperTheory.V.Z.Win.Dn.Adm2A}
\formallabel{UpperTheory.V.Z.Win.Dn.Adm2B}
\formallabel{UpperTheory.GalN.V.Win.TrData}
\formal{UpperTheory.V.Af.eis3}
\end{definition}

\begin{proposition}\label{pf:modularity:numU3}
$\sfname{Hypotheses'}$ implies $\sfname{NumberU3}$.
\formallabel{UpperTheory.GalN.V.Win.N.numberU3U_of_prints}
\formallabel{UpperTheory.GalN.V.Win.N.numberU3U_of_adm2A_imp_adm2B}
\end{proposition}

\begin{proof}
We use the following consequences of $\sfname{Hypotheses'}$, proved in the formal development.
\begin{enumerate}[label=(\roman*)]
\item (\lref{UpperTheory.GalN.V.Win.N.frobenius_trace_count_bound_of_cm_eigensystem_count_alt}, from the counting
statement for conjugate self-dual eigensystems over CM fields \lref{UpperTheory.GalN.V.Win.N.countK_alt} and the
level statement \lref{UpperTheory.V.Keep.levelKA_alt}) For $K$ there are $c_1,A_1$ such that for $e,N,p,w$ as in
$\sfname{NumberU3}$ there is a trace datum $D_1$ for $(q_w,e)$ with $\#\mathcal S_1\le c_1(N(e+1))^{A_1}$ which, for every
$\ell\ne2,p$, $\iota$ and Frobenius element $\sigma$ at $w$, classifies at $\sigma^{-1}$ the representations admissible
in the sense $\sfname{Adm3N}$ of \cite{LeanSource2026} (level $N$, weights $(e,0,-e)$).
\item (\lref{UpperTheory.GalN.V.Win.N.frobenius_trace_count_bound_of_rat_eigensystem_count_alt}, from the counting
statement over $\Q$ of the proof of \cref{pf:modularity:numberSp4} and the statement
\lref{UpperTheory.V.Z.Win.A4.newAdmS_alt}) There are $c_2,A_2$ such that for $e$, a level $M$, $p\nmid M$, a place
$w$ of $K$ above $p$ of norm $p^h$, there is a trace datum $D_2$ for $(p^h,e)$ with $\#\mathcal S_2\le c_2(M(e+1))^{A_2}$
classifying at $\sigma^{-1}$ the endoscopic representations of level $M$.
\item (\lref{UpperTheory.GalN.V.Win.frobenius_trace_norm_nonzero_log_bound}) For $q\ge7$, $e\ge1$, a trace datum $D$
for $(q,e)$ and an $\Aut(\C)$-stable family $(z_o)$ of roots of unity, there is an integer $\mathcal N\ne0$ with
$\log|\mathcal N|\le\#\mathcal O\cdot\#\mathcal S\cdot(2e\log q+1)$ such that for every $\ell\nmid q$, $\iota$, group
element $g$ and class of representations classified by $D$ at $g$, the non-vanishing and divisibility assertions of
$\sfname{NumberU3}$ hold at $g$ for that class. (Here $\mathcal N$ is the product over $(o,s)$ of the algebraic integers
$q^e(t_s-\mathrm{eis}_3(q,e;z_o))$, each of which is non-zero with logarithm at most $2e\log q+1$ by the bound
$|t_s|\le\sqrt q+2$.)
\end{enumerate}
Let $d=|d_K|$, $c=5(c_1+c_2d^{A_2})$ and $A=A_1+A_2+1$. Let $e,N,p,w,(z_o)$ be given. The norm of $w$ is $q_w=p^h$ with
$h\le[K:\Q]=2$ (\lref{UpperTheory.exists_normV_eq_prime_pow_le_finrank}); in particular $q_w\ge7$. Let $D_1$ be as in
(i) and $D_2$ as in (ii) for the level $M=Nd$ (note $p\nmid Nd$), and let $D=D_1\sqcup D_2$; it classifies at $\sigma^{-1}$
both classes of representations (each restricted to its own part). Let $\mathcal N$ be given by (iii) for $D$. With
$Y=N(e+1)$:
\[
 \#\mathcal S\le c_1Y^{A_1}+c_2(dY)^{A_2}\le(c_1+c_2d^{A_2})Y^{A_1+A_2},\qquad 2e\log q_w+1\le4e\log p+1\le5Y\log p,
\]
using $\log p\ge1$ (as $p\ge7$) and $e+1\le Y$; hence $\log|\mathcal N|\le c\,\#\mathcal O\,Y^{A}\log p$. For the
divisibility, let $\ell\ne2,p$ (so $\ell\nmid q_w$), $\iota$, $\sigma$ and $(X_o)$ be given; every representation which is
admissible of unitary type is admissible in the sense $\sfname{Adm3N}$ (\lref{UpperTheory.GalN.V.csdN'_alt},
\lref{UpperTheory.V.Keep.odd3'}), and endoscopic representations are classified by $D_2$; so (iii) applied at
$g=\sigma^{-1}$ gives the assertions of $\sfname{NumberU3}$. Finally, $\sfname{NumberU3}$ follows from
$\sfname{Hypotheses'}$ by this argument with the level statements
\lref{UpperTheory.V.Keep.levelKA_alt}, \lref{UpperTheory.V.Keep.levelPA_alt}, and \lref{UpperTheory.V.Z.Win.A4.admissible_restriction_level_n_of_printed_results}
and \lref{UpperTheory.V.Z.Win.A4.newAdmS_alt}.
\end{proof}

\begin{remark}\label{pf:modularity:rem-numU3}
The formal statement \cref{pf:modularity:numU3} has, besides $\sfname{Hypotheses'}$ and three level statements, the
hypothesis that every endoscopic representation is endoscopic in the sense $\sfname{Adm2B}$; since $\sfname{Adm2B}$ is
defined to be $\sfname{Adm2A}$, this hypothesis is a tautology (it is supplied by
\lref{UpperTheory.GalN.V.adm2B_of_adm2B}).
\end{remark}

\subsection{Standard \texorpdfstring{$L$}{L}-functions of cuspidal representations of \texorpdfstring{$\GL_n$}{GL(n)}}\label{sec:pf:modularity:GJ}

The Eisenstein congruence arguments for $\Uni(3)$ use the analytic theory of the standard $L$-function of cuspidal
representations of $\GL_n$ over a number field, in the form of Godement and Jacquet \cite{GodementJacquet1972}. The
formal development proves the results below; we state them precisely, and since they are classical we refer to
\cite{GodementJacquet1972} for the proofs of the two main inputs, indicating how the formal proof assembles them.

Let $K$ be a number field, $n\ge1$, and for a finite place $v$ of $K$ let $K_v$, $\OO_v$ and $q_v$ be the completion,
its ring of integers and the order of the residue field. For an irreducible admissible representation $\pi_v$ of
$\GL_n(K_v)$, $L(s,\pi_v)=P_{\pi_v}(q_v^{-s})^{-1}$ is its Godement--Jacquet $L$-factor \cite{GodementJacquet1972},
$P_{\pi_v}$ being a polynomial with constant term $1$. If $\pi_v$ has non-zero $\GL_n(\OO_v)$-fixed vectors, a
\emph{Satake parameter} of $\pi_v$ is a tuple $z\in(\C^\times)^n$ such that $\pi_v$ is a constituent of the unramified
principal series with these parameters, in the normalisation of \cite{LeanSource2026}. For a Schwartz--Bruhat function
$\Phi$ on $M_n(K_v)$ and a matrix coefficient $f$ of $\pi_v$,
$Z(\Phi,f,s)=\int_{\GL_n(K_v)}\Phi(g)f(g)|\det g|^{s+(n-1)/2}\,dg$ is the local zeta integral.

\begin{theorem}[Unramified $L$-factors]\label{pf:modularity:lfunprod}
Let $F$ be a non-archimedean local field whose residue field has $q$ elements, $\OO_F$ its ring of integers, $\pi$ an
irreducible admissible representation of $\GL_n(F)$ ($n\ge1$) with non-zero $\GL_n(\OO_F)$-fixed vectors, $\mu$ a Haar
measure on $\GL_n(F)$ and $z$ a Satake parameter of $\pi$ (with respect to $\mu$). Then $P_\pi(X)=\prod_{i=1}^n(1-z_iX)$;
hence $L(s,\pi)=\prod_{i=1}^n(1-z_iq^{-s})^{-1}$ for all $s\in\C$.
\formallabel{GroundTheory.Automorphic.GodementJacquet.LFunction_eq_prod_of_isSatakeParameter}
\formallabel{GroundTheory.Automorphic.GodementJacquet.LPolynomial_eq_prod_of_isSatakeParameter}
\formal{GroundTheory.Automorphic.GodementJacquet.LPolynomial}
\formal{GroundTheory.Automorphic.GodementJacquet.LFunction}
\end{theorem}

\begin{proof}
Recall that $P_\pi$ is defined from the local zeta integrals $Z(\Phi,f,s)$, for Schwartz--Bruhat functions $\Phi$ and
smooth matrix coefficients $f$ of $\pi$, viewed as rational functions of $X=q^{-s}$: $P_\pi$ is the polynomial with
constant term $1$ such that these zeta integrals span $P_\pi^{-1}\C[X,X^{-1}]$ (and $P_\pi=1$ if there is no such
polynomial), and $L(s,\pi)=P_\pi(q^{-s})^{-1}$. The identity $P_\pi=\prod_i(1-z_iX)$ is the unramified computation of
Godement and Jacquet \cite{GodementJacquet1972}: the zeta integral of the characteristic function of $M_n(\OO_F)$ against
the spherical matrix coefficient equal to $1$ at the identity is $\mu(\GL_n(\OO_F))\prod_i(1-z_iq^{-s})^{-1}$, and the zeta
integrals of all pairs $(\Phi,f)$ span $\prod_i(1-z_iX)^{-1}\C[X,X^{-1}]$. Evaluating at $X=q^{-s}$ and inverting the finite
product gives the formula for $L(s,\pi)$.
\end{proof}

\begin{theorem}[Unramified computation]\label{pf:modularity:locUnr}
Let $v$ be a finite place of $K$, $\pi_v$ an irreducible admissible representation of $\GL_n(K_v)$ with non-zero
$\GL_n(\OO_v)$-fixed vectors, and $z$ a Satake parameter of $\pi_v$. Then for all $s\in\C$
\[
 L(s,\pi_v)=\prod_{i=1}^n\bigl(1-z_iq_v^{-s}\bigr)^{-1},\qquad L(s,\pi_v^\vee)=\prod_{i=1}^n\bigl(1-z_i^{-1}q_v^{-s}\bigr)^{-1}.
\]
\formallabel{UpperTheory.V.Z.Inf.locUnr}
\formallabel{UpperTheory.V.Z.Inf.LocUnr}
\end{theorem}

\begin{proof}
The first identity is \cref{pf:modularity:lfunprod}; the second follows from it, since the contragredient of $\pi_v$ is
unramified with Satake parameter $z^{-1}$ (\lref{UpperTheory.V.Z.Inf.LocalUnramifiedTorus.lfun_contra}).
\end{proof}

\begin{littheorem}[Global functional equation of zeta integrals {\cite{GodementJacquet1972}}]\label{pf:modularity:GJglobal}
Let $n\ge2$, $\varphi_1,\varphi_2$ cusp forms on $\GL_n(\A_K)$, with $\varphi_1$ invariant under the central subgroup
$\R_{>0}$ embedded diagonally at the archimedean places, and let
$f(g)=\int\varphi_1(xg)\overline{\varphi_2(x)}\,dx$ be their matrix coefficient (integral over a fundamental domain of
$\GL_n(K)$ in the norm-one subgroup). For every Schwartz--Bruhat function $\Phi$ on $M_n(\A_K)$ there are $\sigma_0\in\R$
and entire functions $\Lambda,\Lambda'$ such that for $\Re s>\sigma_0$ the global zeta integrals
$Z(\Phi,f,s)=\int_{\GL_n(\A_K)}\Phi(g)f(g)|\det g|^{s+(n-1)/2}dg$ and $Z(\widehat\Phi,f^\vee,s)$, where
$f^\vee(g)=f(g^{-1})$ and $\widehat\Phi$ is the adelic Fourier transform, converge absolutely and equal $\Lambda(s)$ and
$\Lambda'(s)$, and $\Lambda(s)=\Lambda'(1-s)$ for all $s$.
\formallabel{GroundTheory.Automorphic.GJGlobal.globalZeta_entire_functionalEquation}
\end{littheorem}

\begin{theorem}[Unramified places]\label{pf:modularity:exunr}
Let $n\ge2$ and $\varphi_1,\varphi_2$, $f$ be as in \cref{pf:modularity:GJglobal}. Let $S$ be a finite set of finite places
of $K$, and for every finite place $v$ let $\pi_v$ be a representation of $\GL_n(K_v)$, $\xi_v$ a smooth vector and
$\lambda_v$ a smooth vector of the contragredient, such that for $v\notin S$: $\pi_v$ is irreducible and admissible, $\xi_v$
and $\lambda_v$ are $\GL_n(\OO_v)$-fixed with $\lambda_v(\xi_v)=1$, and $z_v$ is a Satake parameter of $\pi_v$ (for the
Haar measure giving $\GL_n(\OO_v)$ volume $1$). Let $c_v(g)=\lambda_v(\pi_v(g)\xi_v)$ and $c_v^\vee(g)=c_v(g^{-1})$, and
assume that $f$ is not identically zero and $f(g)=f_\infty(g_\infty)\prod_vc_v(g_v)$ for a function $f_\infty$. Then there
is $\sigma\in\R$ such that for $\Re s>\sigma$:
\begin{enumerate}[label=(\alph*)]
\item for $v\notin S$, $Z(\mathbf 1_{M_n(\OO_v)},c_v,s)=L(s,\pi_v)$ and $Z(\mathbf 1_{M_n(\OO_v)},c_v^\vee,s)=L(s,\pi_v^\vee)$;
\item the products $\prod_{v\notin S}L(s,\pi_v)$ and $\prod_{v\notin S}L(s,\pi_v^\vee)$ converge to $L^S(s,\pi)$ and
$L^S(s,\pi^\vee)$;
\item for $v\notin S$, $|z_{v,i}q_v^{-s}|<1$ for all $i$ and $L(t,\pi_v)=\prod_i(1-z_{v,i}q_v^{-t})^{-1}$ for all $t\in\C$;
and there is $z'_v\in(\C^\times)^n$ with $|z'_{v,i}q_v^{-s}|<1$ for all $i$ and
$L(t,\pi_v^\vee)=\prod_i(1-z'_{v,i}q_v^{-t})^{-1}$ for all $t\in\C$.
\end{enumerate}
\formallabel{GroundTheory.Automorphic.GJStandardL.exists_unramified}
\end{theorem}

\begin{proof}
The matrix coefficient $f$ is continuous
(\lref{GroundTheory.Automorphic.GJStandardL.continuous_globalMatrixCoeff}), and the global zeta integrals of $f$ and of
$f^\vee$ against Schwartz--Bruhat functions converge absolutely for $\Re s$ large, as in \cite{GodementJacquet1972}
(\lref{GroundTheory.Automorphic.GJStandardL.integrable_zetaIntegrand_coeff} and \lref{GroundTheory.Automorphic.GJStandardL.integrable_zetaIntegrand_coeff_inv}).
For a factorisable non-zero continuous $f$ whose global zeta integrals converge absolutely for $\Re s$ large, there is
$\sigma_p$ such that for $\Re s>\sigma_p$ the local zeta integrals $Z(\mathbf 1_{M_n(\OO_v)},c_v,s)$, $v\notin S$, converge
absolutely and their product over $v\notin S$ converges
(\lref{GroundTheory.Automorphic.GJStandardL.exists_sphericalIntegrable}). This applies to $f$, and to $f^\vee$, which
factors as $f^\vee_\infty\prod_vc_v^\vee$ with $c_v^\vee$ the matrix coefficient of the contragredient for the vectors
$\lambda_v$ and $\xi_v$ (\lref{GroundTheory.Automorphic.GJStandardL.inv_factorization}); this gives $\sigma_p$ and
$\sigma_d$, and we put $\sigma=\max(\sigma_p,\sigma_d)$. Let $\Re s>\sigma$ and $v\notin S$. Since $\lambda_v(\xi_v)=1$,
$\xi_v$ is a non-zero spherical vector. By the unramified computation \cite{GodementJacquet1972}, an absolutely convergent
$Z(\mathbf 1_{M_n(\OO_v)},c_v,s)$ equals $\vol(\GL_n(\OO_v))\prod_i(1-z_{v,i}q_v^{-s})^{-1}$ with $|z_{v,i}q_v^{-s}|<1$; as
the volume is $1$ (\lref{GroundTheory.Automorphic.GJStandardL.localHaar_integralGL}), \cref{pf:modularity:lfunprod}
gives the first parts of (a) and (c). Similarly, the absolutely convergent integral $Z(\mathbf 1_{M_n(\OO_v)},c_v^\vee,s)$
equals $L(s,\pi_v^\vee)$, and there is $z'_v$ with $|z'_{v,i}q_v^{-s}|<1$ for all $i$ and
$L(t,\pi_v^\vee)=\prod_i(1-z'_{v,i}q_v^{-t})^{-1}$ for all $t$
(\lref{GroundTheory.Automorphic.GJStandardL.zetaIntegral_inv_sphericalCoeff_of_integrable} and \lref{GroundTheory.Automorphic.GJStandardL.zetaIntegral_inv_sphericalCoeff_of_integrable'}).
Then (b) follows from the convergence of the products of the local zeta integrals.
\end{proof}

\begin{theorem}[Functional equation of the standard $L$-function]\label{pf:modularity:genFE}
Let $n\ge2$, $\varphi_1,\varphi_2$, $f$ be as in \cref{pf:modularity:GJglobal}, and let $S$ be a finite set of finite
places of $K$ containing the places dividing the different of $K$. For every finite place $v$ let $\pi_v$ be an
irreducible admissible representation of $\GL_n(K_v)$, $\xi_v$ a smooth vector and $\lambda_v$ a smooth vector of the
contragredient, such that for $v\notin S$, $\xi_v$ and $\lambda_v$ are $\GL_n(\OO_v)$-fixed with $\lambda_v(\xi_v)=1$
and $z_v$ is a Satake parameter of $\pi_v$. Let $c_v(g)=\lambda_v(\pi_v(g)\xi_v)$ and assume that $f$ is not identically
zero and factors as $f(g)=f_\infty(g_\infty)\prod_vc_v(g_v)$ for a function $f_\infty$ on $\GL_n(K\otimes\R)$. Let
$\Phi_v$ be Schwartz--Bruhat functions on $M_n(K_v)$ for $v\in S$, and $\Phi_v=\mathbf 1_{M_n(\OO_v)}$ for $v\notin S$.
Put $L^S(s,\pi)=\prod_{v\notin S}L(s,\pi_v)$ and $L^S(s,\pi^\vee)=\prod_{v\notin S}L(s,\pi_v^\vee)$.

Then there are constants $c,c'\ne0$ such that the local Fourier transforms $\widehat\Phi_v$ ($v\in S$) are
Schwartz--Bruhat, and for every Schwartz function $\phi$ on $M_n(K\otimes\R)$ there are $\sigma\in\R$ and an entire
function $\Lambda$ such that for $\Re s>\sigma$: $L^S(s,\pi)\ne0$,
\[
 \Lambda(s)=c\,Z_\infty(\phi,f_\infty,s)\prod_{v\in S}Z(\Phi_v,c_v,s)\,L^S(s,\pi),
\]
\[
 \Lambda(1-s)=c'\,Z_\infty(\widehat\phi,f_\infty^\vee,s)\prod_{v\in S}Z(\widehat\Phi_v,c_v^\vee,s)\,L^S(s,\pi^\vee),
\]
where $Z_\infty$ is the archimedean zeta integral, $\widehat\phi$ the archimedean Fourier transform, and
$h^\vee(g)=h(g^{-1})$.
\formallabel{UpperTheory.V.Z.Inf.genFE}
\end{theorem}

\begin{proof}
We use four results of the formal development, besides \cref{pf:modularity:GJglobal}.
\begin{enumerate}[label=(\roman*)]
\item (\cref{pf:modularity:exunr}; unramified theory) There is $\sigma_u$ such
that for $\Re s>\sigma_u$ and $v\notin S$ the local zeta integrals $Z(\mathbf 1_{M_n(\OO_v)},c_v,s)$ and
$Z(\mathbf 1_{M_n(\OO_v)},c_v^\vee,s)$ converge and equal $L(s,\pi_v)$ and $L(s,\pi_v^\vee)$, the products over
$v\notin S$ converge, and the Satake parameters satisfy a bound giving $|z_{v,i}q_v^{-s}|\le q_v^{-2}$ whenever
$\Re s\ge s_0+2$ for any $s_0>\sigma_u$ (\lref{GroundTheory.Automorphic.GJStandardL.norm_mul_qs_le}); since
$\sum_vq_v^{-2}<\infty$, the Euler product of $L^S(s,\pi)$ then converges to a non-zero limit.
\item (\lref{GroundTheory.Automorphic.GJStandardL.matrixTensor_mem}, \lref{GroundTheory.Automorphic.GJStandardL.exists_adelicMatrixFourier_matrixTensor}, and \lref{GroundTheory.Automorphic.GJStandardL.tensor_hyps_fourier})
The pure tensor $\Phi=\phi\otimes\bigotimes_v\Phi_v$ is a Schwartz--Bruhat function on $M_n(\A_K)$; there is $C\ne0$,
independent of $\phi$ and the $\Phi_v$, such that $\widehat\Phi=C\,\widehat\phi\otimes\bigotimes_v\widehat\Phi_v$; and
$\widehat\Phi_v$ is Schwartz--Bruhat for $v\in S$ and equal to $\mathbf 1_{M_n(\OO_v)}$ for $v\notin S$ (because the local
component of the additive character is unramified outside the places dividing the different).
\item (\lref{GroundTheory.Automorphic.GJStandardL.globalZeta_eq_euler}) For a pure tensor and a factorisable
matrix coefficient, whenever the global zeta integral converges absolutely,
$Z(\Phi,f,s)=c_E\,Z_\infty(\phi,f_\infty,s)\prod_vZ(\Phi_v,c_v,s)$, with $c_E>0$ a constant comparing Haar measures.
\item (\lref{GroundTheory.Automorphic.GJStandardL.inv_factorization}) $f^\vee$ factors as
$f_\infty^\vee\prod_vc_v^\vee$, where $c_v^\vee(g)=\lambda_v(\pi_v(g^{-1})\xi_v)$ is a matrix coefficient of
$\pi_v^\vee$, with the roles of $\xi_v$ and $\lambda_v$ exchanged.
\end{enumerate}
Put $c=c_E$ and $c'=Cc_E$. Given $\phi$, let $\sigma_0,\Lambda,\Lambda'$ be given by \cref{pf:modularity:GJglobal} for
$\Phi$, and $\sigma=\max(\sigma_0,\sigma_u)+3$. For $\Re s>\sigma$, $L^S(s,\pi)$ is a convergent product of non-zero
factors with summable deviations from $1$, hence non-zero
(\lref{GroundTheory.Automorphic.GJStandardL.ne_zero_of_hasProd_inv_prod}). By (iii) and (i),
$\Lambda(s)=Z(\Phi,f,s)=c_EZ_\infty(\phi,f_\infty,s)\prod_{v\in S}Z(\Phi_v,c_v,s)\prod_{v\notin S}L(s,\pi_v)$. By the
functional equation, (ii), (iv), (iii) applied to $\widehat\Phi$ and $f^\vee$, and (i) for the contragredients,
$\Lambda(1-s)=\Lambda'(s)=Z(\widehat\Phi,f^\vee,s)=Cc_E\,Z_\infty(\widehat\phi,f_\infty^\vee,s)\prod_{v\in S}
Z(\widehat\Phi_v,c_v^\vee,s)\prod_{v\notin S}L(s,\pi_v^\vee)$.
\end{proof}

\begin{remark}\label{pf:modularity:rem-GJ}
In \cref{pf:modularity:genFE} the cuspidal representation itself does not appear: the hypotheses are formulated
through a factorisation of the matrix coefficient of two cusp forms, with local data $(\pi_v,\xi_v,\lambda_v)$ for which
only irreducibility, admissibility and the spherical normalisation outside $S$ are assumed. The constants $c,c'$ do not
depend on $\phi$, while $\sigma$ and $\Lambda$ do.
\end{remark}

